% Ejercicios y soluciones copiados literalmente de los capítulos.

\chapter{Distancias, normas y topología métrica}

\begin{ejercicio}[Lectura 1, ejercicio 1]\label{ej:lectura-1-1}
Sea $d:X\times X\to\R$ tal que $d(x,y)=0$ si y solo si $x=y$ y $d(x,y)\le d(x,z)+d(y,z)$ para todo $x,y,z$. Demostrar que $d$ es una métrica.
\end{ejercicio}
\begin{solucion}
Tomando $z=x$ se obtiene $d(x,y)\le d(y,x)$, pues $d(x,x)=0$. Intercambiando $x,y$ se obtiene la desigualdad opuesta, y por tanto la simetría. Tomando $y=x$ en la desigualdad dada resulta $0\le2d(x,z)$: así se deduce la no negatividad, que no figuraba entre las hipótesis. La condición de separación ya está supuesta. Finalmente, por simetría la desigualdad dada es $d(x,y)\le d(x,z)+d(z,y)$, la desigualdad triangular usual.
\end{solucion}

\begin{ejercicio}[Lectura 1, ejercicio 2]\label{ej:lectura-1-2}
Sean $d_1,\ldots,d_n$ métricas en $X$. Demostrar que (i) su suma y (ii) su máximo son métricas. (iii) Dar un ejemplo en el que $\min(d_1,d_2)$ no lo sea.
\end{ejercicio}
\begin{solucion}
\begin{enumerate}[label=(\roman*)]
\item La suma es no negativa y simétrica. Si se anula, cada sumando no negativo se anula y $x=y$; si $x=y$, todos son cero. Al sumar las $n$ desigualdades triangulares se obtiene la de la suma.
\item El máximo es no negativo, simétrico y se anula exactamente en la diagonal. Para cada $i$,
\[
d_i(x,y)\le d_i(x,z)+d_i(z,y)\le \max_j d_j(x,z)+\max_jd_j(z,y).
\]
El lado derecho no depende de $i$, por lo que también acota el máximo del lado izquierdo.
\item En $X=\{a,b,c\}$, definamos las distancias entre puntos distintos por
\[
\begin{array}{c|ccc}
&ab&bc&ac\\\hline
d_1&1&3&3\\
d_2&3&1&3
\end{array}
\]
y extendamos por simetría, con diagonal cero. Cada terna de longitudes es $(1,3,3)$; su longitud mayor es menor o igual que la suma de las otras dos, así que ambas funciones son métricas. Su mínimo da distancias $1,1,3$, y falla $d(a,c)\le d(a,b)+d(b,c)$.
\end{enumerate}
\end{solucion}

\begin{ejercicio}[Lectura 1, ejercicio 3]\label{ej:lectura-1-3}
Si $d$ es una métrica, demostrar que $d'(x,y)=d(x,y)/(1+d(x,y))$ también lo es.
\end{ejercicio}
\begin{solucion}
La función $h(t)=t/(1+t)$ es creciente en $[0,\infty)$, es no negativa y solo se anula en $0$. Además, para $s,t\ge0$,
\[
\frac{s+t}{1+s+t}=\frac{s}{1+s+t}+\frac{t}{1+s+t}
\le\frac{s}{1+s}+\frac{t}{1+t}.
\]
Por tanto $h(d(x,y))\le h(d(x,z)+d(z,y))\le h(d(x,z))+h(d(z,y))$. La simetría y la separación se heredan de $d$. Nótese que la nueva métrica está acotada por $1$ aunque la original no lo esté.
\end{solucion}

\begin{ejercicio}[Lectura 1, ejercicio 4]\label{ej:lectura-1-4}
Sean $(d_n)$ métricas en $X$ y $a_n>0$ con $\sum_na_n<\infty$. Demostrar que
\[
D(x,y)=\sum_{n=1}^{\infty}a_n\frac{d_n(x,y)}{1+d_n(x,y)}
\]
es una métrica.
\end{ejercicio}
\begin{solucion}
Cada término está entre $0$ y $a_n$, así que la serie converge y define un número real finito. La simetría se verifica término a término. Si $x=y$, todos los términos son cero. Si $x\ne y$, el primer término es estrictamente positivo, por lo que $D(x,y)>0$. Cada cociente es una métrica por el ejercicio~\ref{ej:lectura-1-3}; al multiplicar su desigualdad triangular por $a_n$ y sumar de $1$ a $N$, obtenemos la desigualdad para sumas parciales. El paso al límite conserva la desigualdad porque las tres series convergen. No se está intercambiando una suma con un ínfimo o un supremo: solo se pasa al límite en una desigualdad entre números reales.
\end{solucion}

\begin{ejercicio}[Lectura 1, ejercicio 5]\label{ej:lectura-1-5}
Para métricas $d_1,\ldots,d_n$ y $p>1$, demostrar que
\[
d_p(x,y)=\left(\sum_{i=1}^n d_i(x,y)^p\right)^{1/p}
\]
es una métrica.
\end{ejercicio}
\begin{solucion}
La no negatividad, simetría y separación siguen de las de cada $d_i$. Poniendo $a_i=d_i(x,z)$ y $b_i=d_i(z,y)$, la desigualdad triangular de cada métrica y la monotonía de $t\mapsto t^p$ dan
\[
d_p(x,y)\le\left(\sum_i(a_i+b_i)^p\right)^{1/p}
\le\left(\sum_i a_i^p\right)^{1/p}+\left(\sum_i b_i^p\right)^{1/p}.
\]
La última desigualdad es Minkowski, demostrada en el teorema~\ref{ch01:minkowski}.
\end{solucion}

\begin{ejercicio}[Lectura 1, ejercicio 6]\label{ej:lectura-1-6}
Demostrar que $|d(x,z)-d(z,y)|\le d(x,y)$.
\end{ejercicio}
\begin{solucion}
De $d(x,z)\le d(x,y)+d(y,z)$ se deduce $d(x,z)-d(y,z)\le d(x,y)$. Intercambiando $x,y$ resulta $d(y,z)-d(x,z)\le d(x,y)$. Estas dos desigualdades equivalen a la del valor absoluto. Esta \emph{desigualdad triangular inversa} será la base de la continuidad de las funciones distancia.\index{desigualdad triangular!inversa}
\end{solucion}

\begin{ejercicio}[Lectura 1, ejercicio 7]\label{ej:lectura-1-7}\index{ultrametrica@ultramétrica}
Una métrica es ultramétrica si $d(x,y)\le\max\{d(x,z),d(y,z)\}$. Demostrar que, cuando $d(x,z)\ne d(y,z)$, se tiene igualdad.
\end{ejercicio}
\begin{solucion}
Supongamos $a=d(x,z)>b=d(y,z)$. La desigualdad ultramétrica da $d(x,y)\le a$. Aplicándola ahora al par $x,z$ con punto intermedio $y$, obtenemos $a\le\max\{d(x,y),b\}$. Como $b<a$, el máximo solo puede ser al menos $a$ si $d(x,y)\ge a$. Por tanto $d(x,y)=a$. Si $b>a$, se intercambian $x$ y $y$ en este mismo argumento. Geométricamente, en un triángulo ultramétrico los dos lados más largos tienen igual longitud.
\end{solucion}

\begin{ejercicio}[Lectura 2, ejercicio 1]\label{ej:lectura-2-1}
En $\R^{\N}$, sea $A=\{a:|a_n|<1/n\text{ para todo }n\}$. Decidir si $A$ es vecindad de $0$ para (i) la métrica discreta; (ii) la métrica $d(a,b)=1/\min\{n:a_n\ne b_n\}$ cuando $a\ne b$, y $d(a,a)=0$; (iii) la métrica
\[
d'(a,b)=\sum_{n=1}^{\infty}\frac{|a_n-b_n|}{2^n(1+|a_n-b_n|)}.
\]
\end{ejercicio}
\begin{solucion}
\begin{enumerate}[label=(\roman*)]
\item Sí: $B(0,1/2)=\{0\}\subset A$.
\item No. Para cada $r>0$, elijamos $N$ con $1/N<r$ y pongamos $a_N=1/N$ y $a_n=0$ si $n\ne N$. La primera coordenada diferente de cero es la $N$, así que $d(0,a)=1/N<r$, pero $a\notin A$. Por tanto ninguna bola de radio positivo está contenida en $A$.

La función propuesta sí es una métrica: si $a$ y $c$ coinciden hasta el índice $m-1$, y $c$ y $b$ también, entonces $a$ y $b$ coinciden hasta ese índice. En consecuencia, el primer desacuerdo entre $a,b$ no puede preceder al menor de los primeros desacuerdos de los otros dos pares. Al invertir los índices se obtiene la desigualdad ultramétrica. La separación y la simetría están en la definición.
\item Tampoco. Con la misma sucesión concentrada en $N$,
\[
d'(0,a)=\frac{1}{2^N(N+1)}.
\]
Este número tiende a cero, de modo que para cualquier $r>0$ podemos elegir $N$ con $d'(0,a)<r$, manteniendo $a\notin A$.

Para verificar que $d'$ es una métrica, la serie converge por comparación con $\sum2^{-n}$. En cada coordenada, $|a_n-b_n|/(1+|a_n-b_n|)$ satisface la desigualdad triangular por el argumento del ejercicio~\ref{ej:lectura-1-3}. Es una pseudométrica en todo $\R^{\N}$, pues puede anularse aunque otras coordenadas difieran. La suma sí separa puntos: si $a\ne b$, alguna coordenada aporta un término positivo. La simetría y la triangular se conservan al sumar y pasar al límite.
\end{enumerate}
Este ejercicio muestra que una condición impuesta simultáneamente a infinitas coordenadas puede impedir la existencia de una sola tolerancia métrica válida para todas ellas.
\end{solucion}

\begin{ejercicio}[Lectura 2, ejercicio 2]\label{ej:lectura-2-2}
Demostrar que una intersección finita de vecindades de un mismo punto es una vecindad de él.
\end{ejercicio}
\begin{solucion}
Para $V_1,\ldots,V_m$ elegimos $r_j>0$ con $B(x,r_j)\subset V_j$. Al ser una familia finita, $r=\min_jr_j>0$, y $B(x,r)\subset\bigcap_jV_j$. La intersección de una familia vacía se interpreta como $X$, que también es vecindad de $x$. La condición de finitud importa: en $\R$, $\bigcap_{n\ge1}(-1/n,1/n)=\{0\}$ no es vecindad de $0$.
\end{solucion}

\begin{ejercicio}[Lectura 2, ejercicio 3]\label{ej:lectura-2-3}
Probar $\cl{A\cup B}=\cl A\cup\cl B$ y $\cl{A\cap B}\subset\cl A\cap\cl B$.
\end{ejercicio}
\begin{solucion}
La monotonía de la clausura da $\cl A\cup\cl B\subset\cl{A\cup B}$. Si $x\notin\cl A\cup\cl B$, existen radios $r,s>0$ cuyas bolas centradas en $x$ no intersectan a $A$ y a $B$, respectivamente. La bola de radio $\min\{r,s\}$ no intersecta a la unión; por tanto $x\notin\cl{A\cup B}$. Esto prueba la otra inclusión por contraposición. Para la intersección, $A\cap B\subset A$ y $A\cap B\subset B$ dan, por monotonía, ambas inclusiones necesarias.
\end{solucion}

\begin{ejercicio}[Lectura 2, ejercicio 4]\label{ej:lectura-2-4}
Dar ejemplos en los que (i) $\interior{(A\cup B)}\not\subset\interior A\cup\interior B$ y (ii) $\cl A\cap\cl B\not\subset\cl{A\cap B}$.
\end{ejercicio}
\begin{solucion}
\begin{enumerate}[label=(\roman*)]
\item Tomemos $A=\Q$ y $B=\R\setminus\Q$. Ambos tienen interior vacío porque todo intervalo contiene racionales e irracionales. Su unión es $\R$, cuyo interior es $\R$.
\item Con los mismos conjuntos, ambas clausuras son $\R$, pero $A\cap B=\varnothing$ tiene clausura vacía. Como alternativa finita, $A=(0,1)$ y $B=(1,2)$ tienen clausuras que se intersectan en $1$, mientras su intersección es vacía.
\end{enumerate}
\end{solucion}

\begin{ejercicio}[Lectura 2, ejercicio 5]\label{ej:lectura-2-5}
En un espacio ultramétrico demostrar: (i) si $y\in B(x,r)$, entonces $B(y,r)=B(x,r)$; (ii) las bolas abiertas y las bolas cerradas de radio positivo son conjuntos abiertos y cerrados; (iii) dos bolas que se intersectan están ordenadas por inclusión.
\end{ejercicio}
\begin{solucion}
\begin{enumerate}[label=(\roman*)]
\item Si $z\in B(y,r)$, entonces $d(x,z)\le\max\{d(x,y),d(y,z)\}<r$. Esto da una inclusión. Como $x\in B(y,r)$, intercambiar los centros da la otra. Para bolas cerradas, el mismo argumento con $\le r$ muestra que cualquier punto de una bola cerrada es también centro de ella con el mismo radio.
\item Las bolas abiertas son abiertas por el ejercicio~\ref{ej:lectura-2-6}. Para ver que $B(x,r)$ es cerrada, sea $z\notin B(x,r)$. Si $w\in B(z,r)$, entonces $d(z,w)<r\le d(z,x)$; el ejercicio~\ref{ej:lectura-1-7} da $d(x,w)=d(x,z)\ge r$. Por tanto $B(z,r)$ no intersecta a $B(x,r)$, y $z$ no pertenece a su clausura.

La bola cerrada es cerrada por el ejercicio~\ref{ej:lectura-2-6}. Si $y\in\overline B(x,r)$, todo $z\in B(y,r)$ satisface $d(x,z)\le\max\{d(x,y),d(y,z)\}\le r$. Así $B(y,r)\subset\overline B(x,r)$ y esta bola es abierta. La formulación «una bola es abierta si y solo si es cerrada» se entiende aquí en sentido topológico; no afirma que una bola definida con $<$ sea idéntica a la definida con $\le$ para el mismo radio.
\item Si dos bolas contienen a $z$, por (i) y su versión cerrada podemos tomar $z$ como centro de ambas. Con centro común, la bola de radio menor está contenida en la de radio mayor. Si los radios coinciden, las bolas son iguales cuando son del mismo tipo; si una es abierta y otra cerrada, la abierta está contenida en la cerrada. Esto cubre también el caso de tipos distintos.
\end{enumerate}
\end{solucion}

\begin{ejercicio}[Lectura 2, ejercicio 6]\label{ej:lectura-2-6}
En cualquier espacio métrico, demostrar que (i) las bolas abiertas son abiertas, (ii) las bolas cerradas son cerradas y (iii) las esferas son cerradas.
\end{ejercicio}
\begin{solucion}
\begin{enumerate}[label=(\roman*)]
\item Si $y\in B(x,r)$, tomamos $s=r-d(x,y)>0$. Para $z\in B(y,s)$, la triangular da $d(x,z)<d(x,y)+s=r$. Así la bola original es vecindad de cada punto suyo y coincide con su interior.
\item Si $y\notin\overline B(x,r)$, sea $s=d(x,y)-r>0$. Para $z\in B(y,s)$, la triangular inversa da $d(x,z)\ge d(x,y)-d(y,z)>r$. Esta bola no intersecta a $\overline B(x,r)$, así que ningún punto exterior pertenece a su clausura.
\item Si $y\notin S(x,r)$, sea $s=|d(x,y)-r|/2>0$. Para $z\in B(y,s)$, la triangular inversa da $|d(x,z)-d(x,y)|<s$. Si $d(x,z)$ fuera $r$, esto contradiría $|r-d(x,y)|=2s$. Por tanto la bola no intersecta a la esfera y esta es cerrada.
\end{enumerate}
\end{solucion}

\begin{ejercicio}[Lectura 2, ejercicio 7]\label{ej:lectura-2-7}
Demostrar que (i) $A$ es abierto si y solo si es vecindad de todos sus puntos; (ii) $A$ es cerrado si y solo si todo punto cuyas vecindades intersectan a $A$ pertenece a $A$.
\end{ejercicio}
\begin{solucion}
\begin{enumerate}[label=(\roman*)]
\item Ser vecindad de cada punto de $A$ significa $A\subset\interior A$. Como siempre $\interior A\subset A$, equivale a la igualdad que define apertura.
\item Los puntos cuyas vecindades intersectan a $A$ son precisamente los de $\cl A$. La condición es $\cl A\subset A$, que, junto con $A\subset\cl A$, equivale a la igualdad que define cerradura.
\end{enumerate}
\end{solucion}

\begin{ejercicio}[Lectura 3, ejercicio 1]\label{ej:lectura-3-1}
Demostrar que (i) $\interior A$ es la unión de todos los abiertos contenidos en $A$ y (ii) $\cl A$ es la intersección de todos los cerrados que contienen a $A$.
\end{ejercicio}
\begin{solucion}
\begin{enumerate}[label=(\roman*)]
\item Si $G$ es abierto y $G\subset A$, la monotonía del interior da $G=\interior G\subset\interior A$. Por otra parte, $\interior A$ es abierto y está contenido en $A$, de modo que él mismo figura en la familia de la unión. Esto prueba las dos inclusiones.
\item Si $F$ es cerrado y $A\subset F$, entonces $\cl A\subset\cl F=F$. Además, $\cl A$ es un cerrado que contiene a $A$, por lo que figura entre los conjuntos que se intersectan. Así es exactamente esa intersección.
\end{enumerate}
\end{solucion}

\begin{ejercicio}[Lectura 3, ejercicio 2]\label{ej:lectura-3-2}
Demostrar que (i) $c_0$, el conjunto de sucesiones que convergen a cero, es cerrado en $\ell^\infty$; (ii) fijado $t\in K$, el conjunto $\{f\in C(K):f(t)=0\}$ es cerrado en la norma uniforme.
\end{ejercicio}
\begin{solucion}
\begin{enumerate}[label=(\roman*)]
\item Sean $a^{(m)}\in c_0$ y $\norm{a^{(m)}-a}_\infty\to0$. Debemos mostrar que $a_n\to0$ cuando cambia el índice de coordenada $n$, no el índice $m$ de aproximación. Dado $\eps>0$, elegimos un solo $m$ con $\norm{a^{(m)}-a}_\infty<\eps/2$. Como $a^{(m)}\in c_0$, existe $N$ con $|a^{(m)}_n|<\eps/2$ para $n\ge N$. Entonces
\[
|a_n|\le |a_n-a^{(m)}_n|+|a^{(m)}_n|<\eps.
\]
Así $a\in c_0$, y el corolario~\ref{ch03:cerradosecuencial} prueba cerradura.
\item Si $f_m(t)=0$ y $\norm{f_m-f}_\infty\to0$, entonces $|f(t)|=|f(t)-f_m(t)|\le\norm{f-f_m}_\infty$. Al pasar al límite obtenemos $|f(t)|=0$. Todo límite en $C(K)$ permanece en el conjunto, por lo que este es cerrado.
\end{enumerate}
\end{solucion}

\begin{ejercicio}[Lectura 3, ejercicio 3]\label{ej:lectura-3-3}\index{Hausdorff, propiedad de}
Para $x\ne y$, encontrar abiertos disjuntos $U,V$ con $x\in U$ y $y\in V$. Deducir la unicidad del límite de una sucesión.
\end{ejercicio}
\begin{solucion}
Sea $r=d(x,y)/3>0$. Las bolas $U=B(x,r)$ y $V=B(y,r)$ son abiertas. Si $z$ perteneciera a ambas, $d(x,y)\le d(x,z)+d(z,y)<2r<d(x,y)$, contradicción. Si una sucesión convergiera a $x$ y a $y$, a partir de un índice estaría en $U$ y a partir de otro estaría en $V$. Desde el mayor de ambos índices estaría en la intersección vacía. Por tanto dos límites deben coincidir. En el enunciado original se sustituye la expresión accidental $y\in X$ por la condición necesaria $y\in V$.
\end{solucion}

\begin{ejercicio}[Lectura 3, ejercicio 4]\label{ej:lectura-3-4}
Demostrar que $\{x\}$ es la intersección de todas las vecindades cerradas de $x$, y deducir que todo conjunto finito es cerrado.
\end{ejercicio}
\begin{solucion}
Todas esas vecindades contienen a $x$. Si $y\ne x$, la bola cerrada $\overline B(x,d(x,y)/2)$ es vecindad de $x$ y no contiene a $y$. Así ningún otro punto pertenece a la intersección. La intersección de cerrados es cerrada, por lo que cada singleton $\{x\}$ es cerrado. Una unión finita de estos singletons es cerrada; el conjunto finito vacío también es cerrado.
\end{solucion}

\begin{ejercicio}[Lectura 3, ejercicio 5]\label{ej:lectura-3-5}
En $\ell^1$, sea $e_k=(\delta_{n,k})_{n\ge1}$ y $A=\{e_k:k\in\N\}$. Demostrar que $A$ no tiene puntos de acumulación.
\end{ejercicio}
\begin{solucion}
Para $j\ne k$, $\norm{e_j-e_k}_1=2$, pues hay exactamente dos coordenadas de valor absoluto $1$. Una bola de radio $1/2$ contiene a lo sumo un elemento de $A$: si contuviera dos, su distancia sería menor que $1$ por la desigualdad triangular. Pero una bola centrada en un punto de acumulación tendría que contener infinitos elementos de $A$, por el lema~\ref{ch03:infinitospuntos}. No existe tal punto. Este es un conjunto infinito sin puntos de acumulación, incluso aunque todos sus elementos tienen norma $1$.
\end{solucion}

\begin{ejercicio}[Lectura 3, ejercicio 6]\label{ej:lectura-3-6}
Demostrar que $A$ es cerrado si y solo si $A'\subset A$.
\end{ejercicio}
\begin{solucion}
Por \eqref{ch03:clderivada}, $\cl A=A\cup A'$. Si $A$ es cerrado, $A'\subset\cl A=A$. Si $A'\subset A$, entonces $\cl A=A\cup A'=A$, de modo que $A$ es cerrado. El enunciado no dice que todo punto de un cerrado sea de acumulación: puede tener puntos aislados.
\end{solucion}

\begin{ejercicio}[Lectura 3, ejercicio 7]\label{ej:lectura-3-7}\index{distancia!a un conjunto}
Para $A\ne\varnothing$, sea $d(x,A)=\inf\{d(x,a):a\in A\}$. Probar $x\in\cl A$ si y solo si $d(x,A)=0$.
\end{ejercicio}
\begin{solucion}
El ínfimo existe y es finito: las distancias son no negativas y alguna de ellas es finita porque $A$ no es vacío. Si $x\in\cl A$, para todo $\eps>0$ hay $a\in A$ con $d(x,a)<\eps$, así que $0\le d(x,A)<\eps$ y el ínfimo es cero. Si el ínfimo es cero, para cada $r>0$ debe existir $a\in A$ con $d(x,a)<r$; de otro modo $r$ sería una cota inferior positiva. Por tanto toda bola intersecta a $A$.

Además, esta distancia es $1$-Lipschitz: para $\eps>0$ elijamos $a\in A$ con $d(y,a)<d(y,A)+\eps$. Entonces $d(x,A)\le d(x,a)\le d(x,y)+d(y,A)+\eps$. Haciendo $\eps\downarrow0$ e intercambiando $x,y$, obtenemos $|d(x,A)-d(y,A)|\le d(x,y)$. No se supuso que el ínfimo se alcanzara.
\end{solucion}

\begin{ejercicio}[Lectura 3, ejercicio 8]\label{ej:lectura-3-8}
Para una familia no vacía $\mathcal S$ de subconjuntos de $X$, demostrar
\[
\left(\bigcap_{A\in\mathcal S}A\right)'\subset\bigcap_{A\in\mathcal S}A',\qquad
\bigcup_{A\in\mathcal S}A'\subset\left(\bigcup_{A\in\mathcal S}A\right)',
\]
y dar ejemplos de inclusiones estrictas.
\end{ejercicio}
\begin{solucion}
Si $E\subset F$, cualquier bola que contiene un punto de $E$ distinto de su centro contiene también un punto de $F$ distinto del centro; por tanto $E'\subset F'$. Aplicamos esta monotonía a $\bigcap\mathcal S\subset A$ para cada $A$ y a $A\subset\bigcup\mathcal S$ para cada $A$.

Para la primera inclusión, en $\R$ tomemos $A=(-1,0)$ y $B=(0,1)$. La intersección es vacía y tiene derivado vacío, mientras que $0\in A'\cap B'$. Para la segunda, sea $\mathcal S=\{\{1/n\}:n\ge1\}$. Cada singleton tiene derivado vacío, pero $0$ pertenece al derivado de la unión, pues $1/n\to0$ y los puntos son distintos.
\end{solucion}

\begin{ejercicio}[Lectura 3, ejercicio 9]\label{ej:lectura-3-9}
En un espacio ultramétrico demostrar: (i) las relaciones $d(x,y)<r$ y $d(x,y)\le r$, con $r>0$, son de equivalencia; (ii) las bolas abiertas de radio $r$ particionan $X$, y lo mismo las cerradas; (iii) si $0<s<r$, una bola $B(x,r)$ es particionada por las bolas abiertas de radio $s$ que la intersectan.
\end{ejercicio}
\begin{solucion}
\begin{enumerate}[label=(\roman*)]
\item La reflexividad usa $d(x,x)=0<r$; la simetría es la de la métrica. Si $d(x,y)<r$ y $d(y,z)<r$, la desigualdad ultramétrica da $d(x,z)\le\max\{d(x,y),d(y,z)\}<r$. Sustituir ambas desigualdades estrictas por $\le r$ demuestra la transitividad de la segunda relación.
\item La clase de equivalencia de $x$ para la primera relación es exactamente $B(x,r)$; para la segunda es $\overline B(x,r)$. Las clases de una relación de equivalencia cubren el conjunto y dos clases son iguales o disjuntas, así que proporcionan las particiones afirmadas. La colección de bolas se entiende sin repetir conjuntos iguales.
\item Si $B(y,s)$ intersecta a $B(x,r)$, elijamos $z$ en la intersección. Entonces $d(x,y)\le\max\{d(x,z),d(z,y)\}<r$, y para $w\in B(y,s)$ se obtiene $d(x,w)\le\max\{d(x,y),d(y,w)\}<r$. Por tanto $B(y,s)\subset B(x,r)$. Todo punto $y$ de la bola grande pertenece a su bola $B(y,s)$, que la intersecta. Las bolas pequeñas son iguales o disjuntas por (ii), de modo que cubren la bola grande sin superposición entre elementos distintos de la colección.
\end{enumerate}
\end{solucion}

\begin{ejercicio}[Lectura 4, ejercicio 1]\label{ej:lectura-4-1}
Demostrar que los abiertos y cerrados relativos son intersecciones con abiertos y cerrados del espacio ambiente.
\end{ejercicio}
\begin{solucion}
La demostración completa es el teorema~\ref{ch04:relativa}. Su construcción explícita, útil al aplicar el resultado, consiste en elegir para cada $a$ del abierto relativo una bola $B_Y(a,r_a)$ contenida en él y tomar en $X$ la unión $G=\bigcup_aB_X(a,r_a)$. Así se obtiene el abierto ambiental requerido. Para un cerrado relativo $A$, se aplica esa construcción a $Y\setminus A$ y se toma el complemento ambiental del abierto obtenido. Ambas direcciones y el caso vacío quedaron comprobados allí.
\end{solucion}

\begin{ejercicio}[Lectura 4, ejercicio 2]\label{ej:lectura-4-2}
Sea $A\subset Y\subset X$. Probar: (i) si $A$ es abierto en $Y$ y $Y$ abierto en $X$, entonces $A$ es abierto en $X$; (ii) si $A$ es cerrado en $Y$ y $Y$ cerrado en $X$, entonces $A$ es cerrado en $X$.
\end{ejercicio}
\begin{solucion}
\begin{enumerate}[label=(\roman*)]
\item Por el teorema~\ref{ch04:relativa}, $A=G\cap Y$ para algún abierto $G$ de $X$. Ambos factores son abiertos en $X$ y una intersección finita de abiertos es abierta.
\item Ahora $A=F\cap Y$ para algún cerrado $F$ de $X$. La intersección de cerrados es cerrada en $X$.
\end{enumerate}
Las hipótesis sobre $Y$ no son superfluas: $[0,1]$ es abierto en sí mismo y no en $\R$; $(0,1)$ es cerrado en sí mismo y no en $\R$.
\end{solucion}

\begin{ejercicio}[Lectura 4, ejercicio 3]\label{ej:lectura-4-3}
Considerar $\Q\subset\R$ con la métrica usual. Para un intervalo abierto relativo no vacío y acotado $I=(a,b)\cap\Q$, con $a<b$, demostrar que $I$ es cerrado en $\Q$ si y solo si ambos extremos de su clausura en $\R$ son irracionales. Precisar el caso de intervalos no acotados.
\end{ejercicio}
\begin{solucion}
La densidad de $\Q$ en $\R$ implica $\cl I^{\,\R}=[a,b]$. En detalle, todo punto de $[a,b]$ puede aproximarse por racionales de $(a,b)$; un punto fuera de $[a,b]$ tiene una bola que no intersecta al intervalo. Por la proposición~\ref{ch04:clausura},
\[
\cl I^{\,\Q}=[a,b]\cap\Q.
\]
Este conjunto coincide con $(a,b)\cap\Q$ exactamente cuando ni $a$ ni $b$ son racionales. Si un extremo es racional, pertenece a la clausura relativa y falta en $I$, por lo que este no es cerrado.

Para $(-\infty,b)\cap\Q$, solo importa el extremo finito $b$: el conjunto es cerrado en $\Q$ exactamente cuando $b\notin\Q$. Para $(a,\infty)\cap\Q$ ocurre lo mismo con $a$. El intervalo $\Q$ entero es abierto y cerrado. Así, en la versión sin acotación debe decirse «todos los extremos finitos son irracionales»; $\pm\infty$ no son números irracionales. El intervalo vacío también es abierto y cerrado y se trata por separado.
\end{solucion}

\begin{ejercicio}[Lectura 4, ejercicio 4]\label{ej:lectura-4-4}
Para $A\subset Y\subset X$, demostrar que el derivado relativo $A'_Y$ está contenido en $A'_X$, y dar un ejemplo de inclusión estricta.
\end{ejercicio}
\begin{solucion}
En realidad se tiene la identidad $A'_Y=Y\cap A'_X$. Para $y\in Y$, la condición $B_Y(y,r)\cap(A\setminus\{y\})\ne\varnothing$ equivale a $B_X(y,r)\cap(A\setminus\{y\})\ne\varnothing$, porque $A\subset Y$. Esto prueba la identidad y la inclusión. Como ejemplo, tomemos $X=\R$ y $Y=A=\{1/n:n\ge1\}$. Cada punto $1/n$ es aislado en $Y$; por ello $A'_Y=\varnothing$. En $\R$, el punto $0$ es de acumulación, de modo que la inclusión es estricta.
\end{solucion}

\begin{ejercicio}[Lectura 4, ejercicio 5: corrección del enunciado]\label{ej:lectura-4-5}
La lectura propone que la imagen de una subsucesión convergente de una sucesión acotada es cerrada en la imagen de la sucesión original. Examinar la afirmación y establecer una versión válida.
\end{ejercicio}
\begin{solucion}
La afirmación escrita es falsa. En $\R$, definamos $x_1=0$ y $x_{n+1}=1/n$. La sucesión es acotada, su imagen es $Y=\{0\}\cup\{1/n:n\ge1\}$ y la subsucesión $(x_{n+1})$ converge a $0$. Su imagen $S=\{1/n:n\ge1\}$ no es cerrada en $Y$, pues $0\in\cl S^{\,Y}\setminus S$.

La versión correcta es la siguiente: si una subsucesión converge a $\ell\in X$ y $S$ es su imagen, entonces
\[
\cl S^{\,X}=S\cup\{\ell\},\qquad
\cl S^{\,Y}=Y\cap(S\cup\{\ell\}).
\]
Para demostrar la primera igualdad, $\ell\in\cl S^{\,X}$ por convergencia y $S\subset\cl S^{\,X}$. Si $z\notin S\cup\{\ell\}$, sea $r=d(z,\ell)/3>0$. Todos los términos suficientemente avanzados están en $B(\ell,r)$, que no intersecta a $B(z,r)$. Los términos anteriores forman un conjunto finito que no contiene a $z$. Podemos reducir $r$ a un radio positivo menor que las distancias de $z$ a esos términos. La bola reducida no intersecta a $S$; así $z\notin\cl S^{\,X}$. La segunda igualdad es la fórmula de clausura relativa.

Por tanto $S$ es cerrado en $Y$ exactamente cuando $\ell\notin Y\setminus S$. En particular, lo es si $\ell\in S$ o si $\ell\notin Y$; y $S\cup\{\ell\}$ siempre es cerrado en $X$. La acotación de la sucesión original no interviene en esta versión correcta.
\end{solucion}

\begin{ejercicio}[Lectura 4, ejercicio 6]\label{ej:lectura-4-6}
¿Es cierto que $A\subset Y$ es abierto en $Y$ si y solo si $A=Y\cap\operatorname{int}_X A$?
\end{ejercicio}
\begin{solucion}
No. Tomemos $X=\R$, $Y=[0,1]$ y $A=Y$. El conjunto $A$ es abierto en $Y$, pero $Y\cap\operatorname{int}_{\R}A=(0,1)\ne A$.

La igualdad propuesta sí es suficiente: si $A=Y\cap\operatorname{int}_X A$, es intersección de $Y$ con un abierto de $X$. De hecho, como $\operatorname{int}_X A\subset A\subset Y$, esa igualdad equivale a $A=\operatorname{int}_X A$, es decir, a apertura en el espacio ambiente, condición más fuerte que la apertura relativa.

La fórmula general correcta para el interior relativo es
\[
\operatorname{int}_Y A
=Y\cap\operatorname{int}_X\bigl(A\cup(X\setminus Y)\bigr).
\]
Para $y\in Y$, una bola relativa está contenida en $A$ si y solo si la bola ambiental correspondiente está contenida en $A\cup(X\setminus Y)$. La equivalencia se verifica punto a punto: los puntos de la bola que no están en $Y$ no imponen restricciones relativas; los que sí están deben pertenecer a $A$.
\end{solucion}

\chapter{Continuidad, densidad y uniformidad}

\begin{ejercicio}[Lectura 5, ejercicio 1]\label{ej:lectura-5-1}
Demostrar que cada una de las siguientes condiciones equivale a continuidad de $f:X\to Y$:
\begin{enumerate}[label=(\roman*)]
\item Para todo $A\subset Y$, $f^{-1}(\operatorname{int}A)\subset\operatorname{int}(f^{-1}(A))$.
\item Para todo $A\subset Y$, $\cl{f^{-1}(A)}\subset f^{-1}(\cl A)$.
\end{enumerate}
En (ii), se corrige la letra $B$, no definida en el original, por $A$.
\end{ejercicio}
\begin{solucion}
\begin{enumerate}[label=(\roman*)]
\item Si $f$ es continua, $f^{-1}(\operatorname{int}A)$ es abierto y está contenido en $f^{-1}(A)$. El interior de un conjunto contiene a todos sus subconjuntos abiertos; de ahí la inclusión. Si la condición vale y $U$ es abierto, entonces $\operatorname{int}U=U$ y $f^{-1}(U)\subset\operatorname{int}(f^{-1}(U))$. La inclusión inversa es automática, de modo que $f^{-1}(U)$ es abierto y $f$ continua.
\item Si $f$ es continua, $f^{-1}(\cl A)$ es cerrado y contiene a $f^{-1}(A)$; por ello contiene a su clausura. Recíprocamente, para un cerrado $F$ tenemos $\cl F=F$, y la condición da $\cl{f^{-1}(F)}\subset f^{-1}(F)$. Por tanto todas las preimágenes de cerrados son cerradas y $f$ es continua.
\end{enumerate}
\end{solucion}

\begin{ejercicio}[Lectura 5, ejercicio 2]\label{ej:lectura-5-2}
Demostrar que una función es continua en el espacio si y solo si es continua en cada punto.
\end{ejercicio}
\begin{solucion}
Las dos direcciones están desarrolladas en el teorema~\ref{ch05:criterios}. El mecanismo es local: la preimagen de un abierto $U$ es abierta exactamente cuando contiene una bola alrededor de cada $x$ con $f(x)\in U$. La continuidad en ese $x$ produce esa bola, y la continuidad global garantiza que la preimagen de una bola alrededor de $f(x)$ es un abierto que contiene a $x$. Así se recupera en cada dirección la definición requerida.
\end{solucion}

\begin{ejercicio}[Lectura 5, ejercicio 3]\label{ej:lectura-5-3}
Probar que $f$ es continua si y solo si las preimágenes de todas las bolas abiertas son abiertas. Dar una función no continua cuyas preimágenes de todas las bolas cerradas sean cerradas.
\end{ejercicio}
\begin{solucion}
La necesidad sigue de que las bolas abiertas son abiertas. Para la suficiencia, todo abierto $U$ del codominio es una unión $\bigcup_iB_i$ de bolas abiertas. Como las preimágenes preservan uniones,
\[
f^{-1}(U)=\bigcup_i f^{-1}(B_i),
\]
que es abierto por hipótesis.

Para la segunda parte, tomemos $X=\R$ con la métrica usual, $Y=\R$ con la discreta y $f(t)=t$. Una bola cerrada de radio $r<1$ en $Y$ es un singleton, y una de radio $r\ge1$ es todo $Y$. Sus preimágenes son singletons o $\R$, todos cerrados en $X$. Pero $\{0\}$ es abierto en $Y$ y su preimagen no es abierta en $X$; así $f$ no es continua. El argumento con bolas abiertas no se transfiere a bolas cerradas, porque una unión arbitraria de cerrados no tiene por qué ser cerrada.
\end{solucion}

\begin{ejercicio}[Lectura 5, ejercicio 4]\label{ej:lectura-5-4}
Sean $f,g:X\to\R$ continuas. Probar que (i) $\{x:f(x)<g(x)\}$ es abierto y (ii) $\{x:f(x)\le g(x)\}$ es cerrado. Se corrige el dominio de $g$ a $X$, necesario para que ambas expresiones estén definidas.
\end{ejercicio}
\begin{solucion}
\begin{enumerate}[label=(\roman*)]
\item Si $f(x)<g(x)$, sea $\eta=(g(x)-f(x))/3>0$. La continuidad de ambas funciones permite elegir $\delta>0$ tal que $d(x,y)<\delta$ implique $|f(y)-f(x)|<\eta$ y $|g(y)-g(x)|<\eta$. Entonces
\[
f(y)<f(x)+\eta<g(x)-\eta<g(y).
\]
La bola está contenida en el conjunto, que es abierto.
\item El complemento es $\{x:g(x)<f(x)\}$, abierto por (i) aplicado con los papeles intercambiados. Por tanto el conjunto original es cerrado.
\end{enumerate}
\end{solucion}

\begin{ejercicio}[Lectura 5, ejercicio 5]\label{ej:lectura-5-5}\index{extension@extensión!por cero}\index{frontera}
Sea $A\subset X$ cerrado y $f:A\to\R$ continua, con $f=0$ en la frontera $\partial A=\cl A\cap\cl{X\setminus A}$. Demostrar que la extensión $g:X\to\R$, igual a $f$ en $A$ y a cero fuera de $A$, es continua.
\end{ejercicio}
\begin{solucion}
Distinguimos tres posiciones del punto $x$.

Si $x\notin A$, como $A$ es cerrado existe una bola alrededor de $x$ contenida en $X\setminus A$. Allí $g$ es constante y es continua en $x$.

Si $x\in\operatorname{int}A$, existe $r>0$ con $B_X(x,r)\subset A$. Dado $\eps>0$, la continuidad relativa de $f$ da $\delta_0>0$ tal que $y\in A$ y $d(x,y)<\delta_0$ implican $|f(y)-f(x)|<\eps$. Tomando $\delta=\min\{r,\delta_0\}$ se obtiene la misma desigualdad para $g$ y todo $y\in X$ con $d(x,y)<\delta$.

Finalmente, si $x\in A\setminus\operatorname{int}A$, todo entorno de $x$ intersecta al complemento de $A$, así que $x\in\partial A$ y $g(x)=f(x)=0$. Por continuidad relativa existe $\delta>0$ con $|f(y)|<\eps$ para $y\in A\cap B(x,\delta)$. Para $y$ de esa bola fuera de $A$, $g(y)=0$; para los puntos dentro de $A$, $|g(y)|=|f(y)|<\eps$. La misma tolerancia sirve para ambos lados de la frontera. Por tanto $g$ es continua en todos los puntos.
\end{solucion}

\begin{ejercicio}[Lectura 5, ejercicio 6: corrección del recíproco]\label{ej:lectura-5-6}\index{soporte}
Sea $S=\supp(f)=\cl{\{x:f(x)\ne0\}}$. Examinar la afirmación «$f$ es continua si y solo si $f=0$ en $\partial S$» y establecer una versión válida.
\end{ejercicio}
\begin{solucion}
La implicación directa es verdadera. Si $x\in\partial S$ y $f(x)\ne0$, por continuidad hay una bola en la que $|f(y)-f(x)|<|f(x)|/2$, así que $|f(y)|>|f(x)|/2>0$. Esa bola está contenida en $S$, lo que convierte a $x$ en punto interior de $S$ y contradice $x\in\partial S$.

El recíproco es falso. En $\R$, tomemos $f(x)=1$ para $x<0$ y $f(x)=2$ para $x\ge0$. La función no es continua en $0$, pero nunca se anula, de modo que $S=\R$ y $\partial S=\varnothing$. La condición en la frontera se satisface de forma vacía y no controla el salto interior.

Una versión exacta es: $f$ es continua en $X$ si y solo si $f|_S$ es continua en el subespacio $S$ y $f=0$ en $\partial S$. La necesidad se sigue de la continuidad de las restricciones y de la prueba anterior. Para la suficiencia, $S$ es cerrado por definición, $f$ es cero en $X\setminus S$ y el ejercicio~\ref{ej:lectura-5-5} aplica a la función continua $f|_S$. Su extensión por cero es precisamente $f$.
\end{solucion}

\begin{ejercicio}[Lectura 5, ejercicio 7]\label{ej:lectura-5-7}
Probar: (i) $f$ es abierta si y solo si las imágenes de las bolas abiertas son abiertas; (ii) si $f$ es biyectiva, es abierta si y solo si es cerrada; (iii) construir una función suprayectiva abierta que no sea cerrada y una cerrada que no sea abierta.
\end{ejercicio}
\begin{solucion}
\begin{enumerate}[label=(\roman*)]
\item La necesidad es inmediata porque una bola abierta es un abierto. Para la suficiencia, escribimos un abierto $U$ como unión de bolas $B_i$; entonces $f(U)=\bigcup_if(B_i)$ es una unión de abiertos.
\item Para una biyección se cumple $f(X\setminus A)=Y\setminus f(A)$. Si $f$ es abierta y $F$ cerrado, el conjunto $Y\setminus f(F)=f(X\setminus F)$ es abierto, así que $f(F)$ es cerrado. Si $f$ es cerrada y $U$ abierto, el conjunto $Y\setminus f(U)=f(X\setminus U)$ es cerrado, de modo que $f(U)$ es abierto. La identidad de complementos utiliza tanto inyectividad como suprayectividad.
\item La proyección $\pi:\R^2\to\R$, $\pi(x,y)=x$, es suprayectiva y abierta. Si $a\in\pi(U)$ con $U$ abierto, existe $b$ con $(a,b)\in U$ y un disco de radio $r>0$ centrado en $(a,b)$ contenido en $U$. Entonces $(a-r,a+r)\subset\pi(U)$, pues $(t,b)$ pertenece a ese disco cuando $|t-a|<r$. Sin embargo, el conjunto $F=\{(x,y):xy=1\}$ es cerrado, como preimagen de $\{1\}$ por la función continua producto, y $\pi(F)=\R\setminus\{0\}$ no es cerrado.

Para una función incluso suprayectiva y cerrada que no sea abierta, tomemos $h:\R\to[0,\infty)$, $h(t)=\max\{t,0\}$. Si $F$ es cerrado, su imagen es $F\cap[0,\infty)$, con el punto $0$ añadido cuando $F$ intersecta $(-\infty,0)$. En ambos casos es un cerrado de $[0,\infty)$. Pero $h((-2,-1))=\{0\}$ no es abierto en $[0,\infty)$: toda bola relativa alrededor de $0$ contiene puntos positivos.
\end{enumerate}
\end{solucion}

\begin{ejercicio}[Lectura 5, ejercicio 8]\label{ej:lectura-5-8}\index{pegado!sobre abiertos}
Sea $(A_i)_{i\in I}$ una cubierta abierta de $X$ y sean $f_i:A_i\to Y$ continuas que coinciden en todas las intersecciones. Probar que la función obtenida pegándolas es continua.
\end{ejercicio}
\begin{solucion}
Para $x\in X$, elijamos $i$ con $x\in A_i$ y definamos $g(x)=f_i(x)$. Si también $x\in A_j$, la hipótesis da $f_i(x)=f_j(x)$; por tanto la definición no depende de la elección. Cualquier función con las restricciones prescritas tiene necesariamente este valor en cada $x$, de modo que es única.

Sea $U$ abierto en $Y$. Entonces
\[
g^{-1}(U)=\bigcup_{i\in I}f_i^{-1}(U).
\]
Cada preimagen es abierta en $A_i$ por continuidad. Como $A_i$ es abierto en $X$, el ejercicio~\ref{ej:lectura-4-2} garantiza que también es abierta en $X$. Su unión es abierta y el criterio global prueba continuidad de $g$. No se requiere que la cubierta sea finita, porque las uniones arbitrarias de abiertos preservan apertura.
\end{solucion}

\begin{ejercicio}[Lectura 5, ejercicio 9]\label{ej:lectura-5-9}
Demostrar que ser homeomorfos es una relación de equivalencia entre espacios métricos.
\end{ejercicio}
\begin{solucion}
La identidad de $X$ es una biyección continua cuya inversa es ella misma; esto prueba reflexividad. Si $h:X\to Y$ es un homeomorfismo, su inversa $h^{-1}:Y\to X$ es biyectiva y continua y tiene inversa continua $h$; esto prueba simetría. Si $h:X\to Y$ y $k:Y\to Z$ son homeomorfismos, $k\circ h$ es una biyección continua y su inversa es $h^{-1}\circ k^{-1}$, continua por composición. Así se prueba transitividad. Cuando se habla de la clase de todos los espacios, «relación de equivalencia» se entiende en ese sentido de clases; sobre cualquier colección que sea un conjunto se obtiene la relación de equivalencia usual.
\end{solucion}

\begin{ejercicio}[Lectura 6, ejercicio 1]\label{ej:lectura-6-1}
Demostrar que $\ell^\infty$ no es separable.
\end{ejercicio}
\begin{solucion}
Sea $E=\{0,1\}^{\N}\subset\ell^\infty$. Este conjunto no es numerable: si sus elementos se listaran como $a^{(1)},a^{(2)},\ldots$, la sucesión $b$ definida por $b_n=1-a^{(n)}_n$ diferiría de la sucesión número $n$ en la coordenada $n$ y no aparecería en la lista. Es el argumento diagonal de Cantor.

Para $a,b\in E$ distintos, existe una coordenada donde difieren en $1$, y ninguna diferencia tiene módulo mayor que $1$; por tanto $\norm{a-b}_\infty=1$. Las bolas $B(a,1/3)$, $a\in E$, son disjuntas dos a dos: un punto común daría por la triangular una distancia entre centros menor que $2/3$.

Si $D=\{d_n:n\in\N\}$ fuera denso, cada bola contendría algún $d_n$. Asignemos a cada $a\in E$ el primer índice de un punto de $D$ que pertenece a $B(a,1/3)$. Bolas disjuntas reciben índices distintos, lo que produciría una inyección del conjunto no numerable $E$ en $\N$. Es imposible; por tanto $\ell^\infty$ no es separable.
\end{solucion}

\begin{ejercicio}[Lectura 6, ejercicio 2]\label{ej:lectura-6-2}
Demostrar que un espacio normado $V$ es separable si y solo si existe un conjunto a lo sumo numerable $L\subset V$ cuyo subespacio generado sea denso.
\end{ejercicio}
\begin{solucion}
Si $D$ es denso y numerable, podemos tomar $L=D$, pues $D\subset\operatorname{span}D$ implica densidad del subespacio generado.

Para el recíproco, sea $\F_0=\Q$ en el caso real y $\F_0=\Q+i\Q$ en el complejo. Sea $D$ el conjunto de todas las combinaciones lineales finitas de elementos de $L$ con coeficientes en $\F_0$, incluyendo el vector cero. Para cada longitud de combinación hay solo un producto finito de conjuntos numerables; al tomar la unión sobre todas las longitudes, $D$ es numerable.

Probemos su densidad. Dado $v\in V$ y $\eps>0$, la densidad de $\operatorname{span}L$ permite escoger $w=\sum_{j=1}^m\alpha_jv_j$ con $\norm{v-w}<\eps/2$. Para cada $j$, elijamos $\beta_j\in\F_0$ tal que
\[
|\alpha_j-\beta_j|<\frac{\eps}{2m(1+\norm{v_j})}.
\]
Los racionales son densos en $\R$, y $\Q+i\Q$ es denso en $\C$ al aproximar por separado las partes real e imaginaria; por ello estas elecciones son posibles. Entonces $u=\sum_{j=1}^m\beta_jv_j\in D$ y
\[
\norm{w-u}\le\sum_{j=1}^m|\alpha_j-\beta_j|\norm{v_j}<\eps/2.
\]
La triangular da $\norm{v-u}<\eps$. Si la combinación es vacía, se toma $u=0$; en particular, $V=\{0\}$ queda cubierto. Así $D$ es denso.
\end{solucion}

\begin{ejercicio}[Lectura 6, ejercicio 3]\label{ej:lectura-6-3}
Si $(A_n)$ es una familia a lo sumo numerable de subespacios separables de un espacio métrico $X$, demostrar que $Y=\bigcup_nA_n$ es separable.
\end{ejercicio}
\begin{solucion}
Para cada $A_n$ no vacío, elijamos $D_n\subset A_n$ denso en su topología relativa y numerable; para el vacío tomemos $D_n=\varnothing$. El conjunto $D=\bigcup_nD_n$ es numerable. Si $y\in Y$ y $r>0$, existe $n$ con $y\in A_n$. Por densidad relativa, $B_X(y,r)\cap A_n$ contiene un punto de $D_n$, y por tanto $B_Y(y,r)$ contiene un punto de $D$. Esto demuestra $\cl D^{\,Y}=Y$. No se necesita que los $A_n$ sean abiertos ni cerrados.
\end{solucion}

\begin{ejercicio}[Lectura 6, ejercicio 4]\label{ej:lectura-6-4}
Demostrar que todo subespacio de un espacio métrico separable es separable.
\end{ejercicio}
\begin{solucion}
Sea $\{B_n\}$ una base numerable de $X$, cuya existencia asegura el teorema~\ref{ch06:separablebase}. Los conjuntos $Y\cap B_n$ forman una base de $Y$: si $U=Y\cap G$ es abierto relativo y $y\in U$, hay $B_n$ con $y\in B_n\subset G$, y entonces $y\in Y\cap B_n\subset U$. Esta base es numerable. Elegimos un punto de cada elemento no vacío de ella; por el lema~\ref{ch06:baseadenso}, esos puntos forman un denso numerable de $Y$.

No basta intersectar con $Y$ un denso numerable prefijado de $X$. Por ejemplo, $\Q$ es denso en $\R$, pero su intersección con $Y=\R\setminus\Q$ es vacía. La prueba con una base elige puntos dentro del subespacio y evita ese error.
\end{solucion}

\begin{ejercicio}[Lectura 6, ejercicio 5]\label{ej:lectura-6-5}
Sean $X_1,\ldots,X_m$ espacios métricos no vacíos. Demostrar que su producto finito, con una métrica producto, es separable si y solo si lo es cada factor.
\end{ejercicio}
\begin{solucion}
Podemos usar $d(x,y)=\max_{1\le j\le m}d_j(x_j,y_j)$. Su bola de radio $r$ es el producto de las bolas de radio $r$ de los factores. Si $D_j$ es denso numerable en $X_j$, el producto $D=\prod_{j=1}^mD_j$ es numerable. Dado $x$ y $r>0$, elegimos $a_j\in D_j$ con $d_j(x_j,a_j)<r$ para cada uno de los finitísimos índices. Entonces $a\in D\cap B(x,r)$ y $D$ es denso.

Recíprocamente, sea $D$ denso numerable en el producto. La proyección $\pi_j(D)$ es numerable. Dado $x_j\in X_j$ y $r>0$, elegimos puntos arbitrarios en los demás factores y formamos $x$ con coordenada $j$ igual a $x_j$. Una aproximación $a\in D$ con $d(a,x)<r$ proporciona $d_j(a_j,x_j)<r$, así que $\pi_j(D)$ es denso en $X_j$.

Las métricas producto finitas dadas por suma, máximo o norma $p$ inducen la misma topología: para números no negativos $t_j$, $\max_jt_j\le(\sum_jt_j^p)^{1/p}\le m^{1/p}\max_jt_j$ y $\max_jt_j\le\sum_jt_j\le m\max_jt_j$. Estas cotas comparan las bolas y trasladan el resultado a cualquiera de esas métricas. La hipótesis de factores no vacíos es necesaria: un factor vacío hace que el producto sea vacío, aunque otro factor sea no separable.
\end{solucion}

\begin{ejercicio}[Lectura 6, ejercicio 6]\label{ej:lectura-6-6}
Demostrar que un conjunto denso en ninguna parte tiene complemento denso, y que un conjunto denso y cerrado tiene complemento denso en ninguna parte.
\end{ejercicio}
\begin{solucion}
Si $\operatorname{int}(\cl A)=\varnothing$, ningún abierto no vacío $U$ está contenido en $\cl A$. Por tanto $U\cap(X\setminus\cl A)\ne\varnothing$, y en particular $U\cap(X\setminus A)\ne\varnothing$. Así el complemento de $A$ es denso. Hemos probado algo más fuerte: el complemento de la clausura es abierto y denso.

Si $A$ es denso y cerrado, $A=\cl A=X$; por tanto su complemento es vacío, cuya clausura tiene interior vacío. La segunda afirmación de la lectura es correcta, aunque trivial con ambas hipótesis. Una formulación adicional útil es que el complemento de un abierto denso es cerrado y denso en ninguna parte: si $G$ es abierto y denso, $F=X\setminus G$ es cerrado y no contiene ningún abierto no vacío, porque todo abierto no vacío intersecta a $G$. Así $\operatorname{int}(\cl F)=\operatorname{int}F=\varnothing$.
\end{solucion}

\begin{ejercicio}[Lectura 6, ejercicio 7]\label{ej:lectura-6-7}
Demostrar que en un espacio métrico no vacío un conjunto no puede ser simultáneamente denso y denso en ninguna parte.
\end{ejercicio}
\begin{solucion}
Si $A$ es denso, $\cl A=X$. Entonces $\operatorname{int}(\cl A)=\operatorname{int}X=X\ne\varnothing$, lo que contradice ser denso en ninguna parte. La condición de espacio no vacío evita la excepción $X=A=\varnothing$.
\end{solucion}

\begin{ejercicio}[Lectura 6, ejercicio 8]\label{ej:lectura-6-8}\label{ch06:intervalos}
Demostrar que todo abierto de $\R$ es unión a lo sumo numerable de intervalos abiertos disjuntos dos a dos.
\end{ejercicio}
\begin{solucion}
Sea $U\subset\R$ abierto. Si es vacío, la unión vacía resuelve el problema. Para $x\in U$, definamos $I_x$ como la unión de todos los intervalos abiertos contenidos en $U$ que contienen a $x$. Hay al menos uno porque $U$ es abierto. La unión es abierta y está contenida en $U$.

Además, $I_x$ es un intervalo: si $a<b$ pertenecen a $I_x$, hay intervalos $J_a,J_b$ de la familia que contienen respectivamente a $a,b$ y ambos contienen a $x$. Como se intersectan, su unión es un intervalo. Para justificar este último hecho, un punto entre dos puntos de $J_a\cup J_b$ queda en uno de los intervalos, o entre ellos; en el segundo caso un punto común a ambos permite dividir el segmento en dos partes contenidas en los respectivos intervalos. Por ello todos los puntos entre $a,b$ pertenecen a $J_a\cup J_b\subset I_x$.

Así $I_x$ es un intervalo abierto, posiblemente una semirrecta o toda $\R$. Es maximal entre los intervalos de $U$ que contienen a $x$, porque su definición ya reúne a todos ellos. Si $I_x$ e $I_y$ se intersectan, su unión es un intervalo abierto contenido en $U$; contiene tanto a $x$ como a $y$. Por maximalidad está contenida tanto en $I_x$ como en $I_y$, y por tanto $I_x=I_y$. Los intervalos distintos de esta colección son disjuntos y cubren $U$ porque $x\in I_x$.

Cada intervalo no vacío contiene un racional. Fijemos una enumeración de $\Q$ y asignemos a cada intervalo el primer racional de la lista que contiene. Dos intervalos disjuntos reciben racionales distintos, de modo que la familia se inyecta en $\Q$ y es a lo sumo numerable. Esto concluye la descomposición, incluidos los intervalos no acotados.
\end{solucion}

\begin{ejercicio}[Lectura 7, ejercicio 1]\label{ej:lectura-7-1}
Demuestre que la composición de dos funciones Lipschitz es Lipschitz y determine una constante para la composición.
\end{ejercicio}
\begin{solucion}
Si $f:X\to Y$ tiene constante $L$ y $g:Y\to Z$ tiene constante $M$, entonces
\[
 \sigma(g(f(x)),g(f(y)))\leq M\rho(f(x),f(y))\leq MLd(x,y).
\]
Así, $ML$ es una constante válida, incluso cuando alguno de los factores es cero. No necesariamente es la mejor: si $g$ es constante en $f(X)$, la composición tiene constante óptima cero aunque $L,M$ sean positivos.
\end{solucion}

\begin{ejercicio}[Lectura 7, ejercicio 2]\label{ej:lectura-7-2}
Sean $f:X\to Y$ y $g:Y\to Z$ continuas. Determine si $g\circ f$ debe ser uniformemente continua cuando: \begin{enumerate}[label=(\roman*)]
\item $f$ es uniformemente continua;
\item $g$ es uniformemente continua.
\end{enumerate}
\end{ejercicio}
\begin{solucion}
En ambos casos la respuesta es negativa. En (i) tomamos $X=Y=Z=\R$, $f(x)=x$ y $g(x)=x^2$. La identidad es $1$-Lipschitz y la composición es el cuadrado, que falla por el ejemplo~\ref{ch07:cuadrado}. En (ii) intercambiamos los papeles: $f(x)=x^2$ y $g(x)=x$. Ahora $g$ es $1$-Lipschitz y la composición vuelve a ser $x^2$.
\end{solucion}

\begin{ejercicio}[Lectura 7, ejercicio 3]\label{ej:lectura-7-3}
Sea $A$ un cerrado no vacío de un espacio métrico $X$. Pruebe que $x\mapsto\dist(x,A)$ es uniformemente continua.
\end{ejercicio}
\begin{solucion}
Para $a\in A$, la desigualdad triangular da $d(x,a)\leq d(x,y)+d(y,a)$. Tomando ínfimos en $a$ obtenemos
\[
 \dist(x,A)\leq d(x,y)+\dist(y,A).
\]
Al intercambiar $x,y$ concluimos
\[
 |\dist(x,A)-\dist(y,A)|\leq d(x,y).
\]
La función es $1$-Lipschitz. La prueba no exige que $A$ sea cerrado; esa hipótesis solamente permite identificar $\dist(x,A)=0$ con $x\in A$. Se agrega la condición $A\ne\varnothing$ porque la distancia al vacío no es una función real finita.
\end{solucion}

\begin{ejercicio}[Lectura 7, ejercicio 4]\label{ej:lectura-7-4}
Dé una función uniformemente continua entre espacios métricos que no sea abierta.
\end{ejercicio}
\begin{solucion}
La función constante $f:\R\to\R$, $f(x)=0$, es $0$-Lipschitz. Sin embargo, la imagen del abierto $(-1,1)$ es $\{0\}$, que no es abierto en $\R$. La continuidad controla preimágenes de abiertos; ser abierta controla imágenes y es una condición distinta.
\end{solucion}

\begin{ejercicio}[Lectura 7, ejercicio 5]\label{ej:lectura-7-5}\label{ch07:ej-extension}
Sea $A$ abierto en un espacio métrico $X$ y sea $f:X\to\R$ uniformemente continua en $A$. Demuestre que $f|_A$ se extiende continuamente a $\cl A$. Dé un contraejemplo si se reemplaza $\R$ por un espacio métrico arbitrario.
\end{ejercicio}
\begin{solucion}
Probaremos la afirmación más fuerte en que $A$ es cualquier subconjunto y el espacio de llegada $Y$ es completo. Si $A=\varnothing$, la extensión tiene dominio vacío y no hay nada que definir. Para $x\in\cl A$ elegimos $a_n\in A$ con $d(a_n,x)<1/n$. Esta sucesión converge y es de Cauchy. Por el teorema~\ref{ch07:cauchy-imagen}, $(f(a_n))$ es de Cauchy; la completitud de $Y$ proporciona un límite. Definimos
\[
 g(x)=\lim_{n\to\infty}f(a_n).
\]
La elección no afecta al valor: si $b_n\to x$ en $A$, entonces $d(a_n,b_n)\leq d(a_n,x)+d(x,b_n)\to0$; la uniformidad implica $\rho(f(a_n),f(b_n))\to0$. Tomando límites, los dos valores propuestos coinciden. Si $x\in A$, podemos elegir la sucesión constante $a_n=x$, por lo que $g(x)=f(x)$.

Además, $g$ es uniformemente continua. Dado $\eps>0$, elijamos $\delta>0$ de modo que $d(a,b)<\delta$ implique $\rho(f(a),f(b))<\eps/2$ para $a,b\in A$. Si $x,y\in\cl A$ y $d(x,y)<\delta/3$, elegimos $a_n\to x$, $b_n\to y$. Para $n$ suficientemente grande, $d(a_n,b_n)<\delta$; al tomar límites queda $\rho(g(x),g(y))\leq\eps/2<\eps$. La extensión es única porque dos aplicaciones continuas que coinciden en un conjunto denso coinciden en su clausura: evaluamos ambas sobre una sucesión $a_n\to x$.

Para el contraejemplo tomamos $X=\R$, $A=(0,1)$ y $Y=(0,1)$ con la distancia usual. Definimos $f(x)=x$ en $A$ y $f(x)=1/2$ fuera de $A$. Su restricción a $A$ es $1$-Lipschitz. Una extensión continua $g:[0,1]\to(0,1)$ tendría $g(0)=\lim_{n\to\infty}g(1/(n+2))=0$, imposible en el espacio de llegada. Lo que falla es su completitud.
\end{solucion}

\begin{ejercicio}[Lectura 7, ejercicio 6]\label{ej:lectura-7-6}
Para subconjuntos no vacíos $A,B$ de $X$, defina $d(A,B)=\inf\{d(a,b):a\in A,b\in B\}$. Pruebe que $f:X\to Y$ es uniformemente continua si y sólo si
\[
 d(A,B)=0\Longrightarrow\rho(f(A),f(B))=0
\]
para todos esos conjuntos.
\end{ejercicio}
\begin{solucion}
Si $f$ es uniformemente continua, dado $\eps>0$ elegimos el correspondiente $\delta>0$. La igualdad $d(A,B)=0$ permite encontrar $a\in A,b\in B$ con $d(a,b)<\delta$. Entonces $\rho(f(A),f(B))\leq\rho(f(a),f(b))<\eps$. Un número no negativo menor que todo $\eps>0$ es cero.

Para el recíproco hay que cuidar las \emph{distancias cruzadas}: saber que $\rho(f(x_n),f(y_n))\geq\eps_0$ no basta por sí solo para separar los conjuntos de imágenes. Si $f$ no es uniformemente continua, existen $\eps_0>0$ y pares $x_n,y_n$ tales que
\[
 d(x_n,y_n)<1/n,\qquad \rho(p_n,q_n)\geq\eps_0,
 \quad p_n=f(x_n),\ q_n=f(y_n).
\]
Pongamos $r=\eps_0/8$. Si alguna bola $B(z,r)$ contiene infinitos $p_n$, retenemos esos índices $I$. Para $i,j\in I$, tenemos $\rho(q_j,z)\geq\eps_0-r$ y, por tanto, $\rho(p_i,q_j)>\eps_0-2r$. Si alguna bola contiene infinitos $q_n$, el argumento con los papeles intercambiados da la misma separación.

Queda el caso en que toda bola de radio $r$ contiene sólo un número finito de términos de cada sucesión, contando sus índices. Elegimos inductivamente índices crecientes $n_k$. Una vez elegidos los primeros $k$, descartamos todos los índices $n$ para los que $p_n\in B(q_{n_j},r)$ o $q_n\in B(p_{n_j},r)$ para algún $j\leq k$. Sólo descartamos una unión finita de conjuntos finitos; por ello queda un índice mayor que $n_k$. La sucesión retenida satisface $\rho(p_{n_i},q_{n_j})\geq r$ si $i\ne j$, y la misma desigualdad vale cuando $i=j$ porque $\eps_0\geq r$.

En cualquiera de los casos obtenemos un conjunto infinito de índices $I$ y una constante $c>0$ tales que $\rho(p_i,q_j)\geq c$ para todo $i,j\in I$. Para $A=\{x_i:i\in I\}$ y $B=\{y_i:i\in I\}$ se cumple $d(A,B)=0$, pues $I$ no está acotado y $d(x_i,y_i)<1/i$, pero $\rho(f(A),f(B))\geq c$. Esto contradice la propiedad supuesta.
\end{solucion}

\begin{ejercicio}[Lectura 7, ejercicio 7]\label{ej:lectura-7-7}
Suponga que $f:X\to Y$ satisface $\rho(f(x),f(y))\geq\alpha d(x,y)$ para alguna $\alpha>0$. Pruebe que es inyectiva, que su inversa sobre la imagen es Lipschitz y que $f$ es abierta si $f(X)$ es abierto en $Y$.
\end{ejercicio}
\begin{solucion}
Si $f(x)=f(y)$, la desigualdad da $0\geq\alpha d(x,y)$ y por tanto $x=y$. Para $u=f(x),v=f(y)$ obtenemos
\[
 d(f^{-1}(u),f^{-1}(v))\leq\alpha^{-1}\rho(u,v).
\]
Así, $f^{-1}:f(X)\to X$ es continua. Si $U\subset X$ es abierto, $f(U)=(f^{-1})^{-1}(U)$ es abierto en el subespacio $f(X)$. Cuando $f(X)$ es abierto en $Y$, todo abierto suyo es abierto en $Y$, y $f$ es una aplicación abierta. La hipótesis no afirma que $f$ sea continua; controla únicamente su inversa.
\end{solucion}

\chapter{Espacios completos y completación}

\begin{ejercicio}[Lectura 8, ejercicio 1]\label{ej:lectura-8-1}
Demuestre que $d_p$ es una métrica en $\Q$ y, más precisamente, una ultramétrica.
\end{ejercicio}
\begin{solucion}
La factorización prima garantiza la unicidad de $v_p(r)$ y $v_p(rs)=v_p(r)+v_p(s)$ para racionales no nulos. Sean $r,s\ne0$ y supongamos $v_p(r)=k\leq v_p(s)=l$. Escribiendo $r=p^ka/b$ y $s=p^lc/d$ con denominadores y numeradores no divisibles por $p$, obtenemos
\[
 r+s=p^k\frac{ad+p^{l-k}bc}{bd}.
\]
El denominador no es divisible por $p$, y el numerador es un entero. Si no es cero, su valoración es no negativa; de aquí $v_p(r+s)\geq k$. Si es cero, $|r+s|_p=0$. En ambos casos
\[
 |r+s|_p\leq p^{-k}=\max\{|r|_p,|s|_p\}.
\]
Si alguno de los sumandos es cero, la desigualdad se verifica directamente. También $|r|_p\geq0$, con igualdad exactamente para $r=0$, y $|-r|_p=|r|_p$. Por tanto $d_p$ es positiva, separa puntos y es simétrica. Aplicando la desigualdad anterior a $(r-t)+(t-s)$ obtenemos
\[
 d_p(r,s)\leq\max\{d_p(r,t),d_p(t,s)\}\leq d_p(r,t)+d_p(t,s).
\]
La primera desigualdad es la propiedad ultramétrica y la segunda da la desigualdad triangular usual.
\end{solucion}

\begin{ejercicio}[Lectura 8, ejercicio 2]\label{ej:lectura-8-2}
Demuestre la equivalencia entre el axioma del supremo y la completitud de $\R$ con su distancia usual.
\end{ejercicio}
\begin{solucion}
Trabajamos con las propiedades de cuerpo ordenado arquimediano de los reales; el contenido del ejercicio es comparar las dos formulaciones de su propiedad de completitud.

Supongamos el axioma del supremo y sea $(x_n)$ de Cauchy. Está acotada. Para cada $N$, definamos $a_N=\inf_{n\geq N}x_n$ y $b_N=\sup_{n\geq N}x_n$. El ínfimo existe aplicando el axioma del supremo al conjunto de los opuestos. La sucesión $(a_N)$ crece, $(b_N)$ decrece y $a_N\leq b_N$. Para cualquier $\eps>0$, la condición de Cauchy aplicada con $\eps/2$ implica que, a partir de cierto $N$, todos los términos de la cola están en un intervalo de longitud $\eps$; por ello $b_N-a_N\leq\eps$. Sea $x=\sup_Na_N$. Para todo $N$, se tiene $a_N\leq x\leq b_N$: la segunda desigualdad se debe a que $b_N$ es una cota superior de todos los $a_M$, también de los correspondientes a $M<N$. Si $n\geq N$, tanto $x_n$ como $x$ están en $[a_N,b_N]$, de modo que $|x_n-x|\leq b_N-a_N\to0$.

Recíprocamente, supongamos que toda sucesión real de Cauchy converge. Sea $A\ne\varnothing$ acotado superiormente. Si tiene máximo, ése es su supremo. En otro caso elegimos $a_0\in A$ y una cota superior $b_0>a_0$. La sucesión de intervalos se construye por bisección: si $m_n=(a_n+b_n)/2$ es cota superior de $A$, ponemos $(a_{n+1},b_{n+1})=(a_n,m_n)$; si no lo es, ponemos $(a_{n+1},b_{n+1})=(m_n,b_n)$. Aquí $a_n$ no tiene por qué pertenecer a $A$, pero nunca es cota superior, y $b_n$ siempre lo es. Los intervalos están anidados y su longitud es $(b_0-a_0)2^{-n}$. Esto hace que ambas sucesiones sean de Cauchy y que tengan el mismo límite $s$.

Todo $a\in A$ satisface $a\leq b_n$ para cada $n$, luego $a\leq s$. Si $t<s$, para algún $n$ se tiene $t<a_n$, y como $a_n$ no es cota superior existe $a\in A$ con $a>a_n>t$. Así, ningún número menor que $s$ es cota superior. Hemos probado $s=\sup A$. Si se desean puntos de $A$ como en la sugerencia original, elegimos $c_n\in A$ con $a_n<c_n\leq b_n$; entonces también $c_n\to s$.
\end{solucion}

\begin{ejercicio}[Lectura 8, ejercicio 3]\label{ej:lectura-8-3}
Demuestre que $\ell^\infty$ es completo.
\end{ejercicio}
\begin{solucion}
Sea $(x^{(n)})$ de Cauchy para $\norm{x}_\infty=\sup_k|x_k|$. Cada coordenada es de Cauchy, porque $|x^{(n)}_k-x^{(m)}_k|\leq\norm{x^{(n)}-x^{(m)}}_\infty$. Definamos $x_k=\lim_nx^{(n)}_k$. Existe $M$ con $\norm{x^{(n)}}_\infty\leq M$, y al pasar al límite para cada $k$ resulta $|x_k|\leq M$; así $x\in\ell^\infty$. Dado $\eps>0$, elegimos $N$ con $\norm{x^{(n)}-x^{(m)}}_\infty<\eps/2$ para $n,m\geq N$. Fijamos $n\geq N$ y dejamos $m\to\infty$: para cada $k$, $|x^{(n)}_k-x_k|\leq\eps/2$. Tomar el supremo en $k$ produce $\norm{x^{(n)}-x}_\infty\leq\eps/2<\eps$.
\end{solucion}

\begin{ejercicio}[Lectura 8, ejercicio 4]\label{ej:lectura-8-4}
Pruebe que los subespacios $c$ y $c_0$ de sucesiones convergentes y de sucesiones convergentes a cero son espacios de Banach con la norma uniforme.
\end{ejercicio}
\begin{solucion}
Ambos son subespacios lineales de $\ell^\infty$: las combinaciones lineales preservan la existencia del límite y, para $c_0$, preservan el límite cero. Por el teorema~\ref{ch08:cerrado-completo}, basta probar que son cerrados.

Sea $x^{(n)}\in c$ y $\norm{x^{(n)}-x}_\infty\to0$. Para $\eps>0$ elegimos $n$ con $\norm{x^{(n)}-x}_\infty<\eps/3$. Como la sucesión escalar $k\mapsto x^{(n)}_k$ converge, existe $K$ con $|x^{(n)}_k-x^{(n)}_j|<\eps/3$ para $k,j\geq K$. Por tanto
\[
 |x_k-x_j|\leq|x_k-x^{(n)}_k|+|x^{(n)}_k-x^{(n)}_j|+|x^{(n)}_j-x_j|<\eps.
\]
La sucesión de coordenadas de $x$ es de Cauchy en los escalares y converge; de ahí $x\in c$.

Si cada $x^{(n)}\in c_0$, elegimos $n$ con $\norm{x^{(n)}-x}_\infty<\eps/2$ y luego $K$ con $|x^{(n)}_k|<\eps/2$ para $k\geq K$. Entonces $|x_k|<\eps$ para $k\geq K$, luego $x\in c_0$.
\end{solucion}

\begin{ejercicio}[Lectura 8, ejercicio 5]\label{ej:lectura-8-5}
Para espacios métricos no vacíos, pruebe que el producto de dos espacios es completo si y sólo si cada factor lo es. Extienda el resultado a todo producto finito.
\end{ejercicio}
\begin{solucion}
Usamos la métrica máxima $D((x,y),(x',y'))=\max\{d(x,x'),\rho(y,y')\}$. Una sucesión es de Cauchy para $D$ exactamente cuando sus dos coordenadas lo son: una cota para $D$ acota ambas distancias, y dos umbrales para las coordenadas se reúnen tomando su máximo. Si ambos factores son completos, las coordenadas convergen a $x,y$, y el máximo de sus distancias a esos límites tiende a cero.

Si el producto es completo, fijamos $y_0\in Y$ y consideramos una sucesión de Cauchy $(x_n)$ de $X$. Los pares $(x_n,y_0)$ forman una sucesión de Cauchy que converge a algún $(x,y)$, y la primera coordenada da $x_n\to x$. Fijando $x_0$ se demuestra lo mismo para $Y$. Para $n$ factores, la métrica $\max_{1\leq j\leq n}d_j(x_j,y_j)$ permite repetir exactamente las dos implicaciones: tomamos el máximo de los $n$ umbrales y, para la necesidad, fijamos las otras $n-1$ coordenadas. La métrica suma y la euclídea de las distancias son equivalentes mediante cotas constantes a la métrica máxima; por ello tienen las mismas sucesiones de Cauchy. La no vacuidad evita que un producto vacío oculte un factor incompleto.
\end{solucion}

\begin{ejercicio}[Lectura 8, ejercicio 6]\label{ej:lectura-8-6}
Demuestre que el conjunto de Cantor, con la distancia heredada de $\R$, es completo y no tiene puntos aislados.
\end{ejercicio}
\begin{solucion}
Sea $C_0=[0,1]$ y sea $C_n$ el conjunto que queda después de eliminar los tercios medios abiertos durante las primeras $n$ etapas. Cada $C_n$ es una unión finita de intervalos cerrados de longitud $3^{-n}$, y $C=\bigcap_{n\geq0}C_n$ es cerrado en $\R$. Como $\R$ es completo, también lo es $C$.

Fijemos $x\in C$ y $\eps>0$. Elegimos $n$ con $3^{-n}<\eps$ y un intervalo componente $[a,b]$ de $C_n$ que contenga $x$. Ambos extremos pertenecen a todas las etapas posteriores: nunca se elimina un extremo al retirar un tercio medio abierto. Por tanto $a,b\in C$. Como $a\ne b$, al menos uno de ellos es distinto de $x$ y dista de $x$ a lo sumo $3^{-n}<\eps$. Ningún punto de $C$ es aislado.
\end{solucion}

\begin{ejercicio}[Lectura 8, ejercicio 7]\label{ej:lectura-8-7}
Demuestre que $C$ es incontable y que
\[
 C=\left\{\sum_{n=1}^\infty\frac{a_n}{3^n}:a_n\in\{0,2\}\right\}.
\]
\end{ejercicio}
\begin{solucion}
Para cada palabra $(a_1,\ldots,a_n)$ de ceros y doses, una componente de $C_n$ es
\[
 I(a_1,\ldots,a_n)=\left[\sum_{j=1}^n\frac{a_j}{3^j},\,
 \sum_{j=1}^n\frac{a_j}{3^j}+3^{-n}\right].
\]
Esto se verifica por inducción: escoger el tercio izquierdo agrega el dígito $0$ y escoger el derecho agrega el dígito $2$. Una serie con esos dígitos pertenece al intervalo indicado, pues su cola está entre cero y $\sum_{j>n}2/3^j=3^{-n}$. Pertenece entonces a todo $C_n$ y a $C$.

Recíprocamente, si $x\in C$, en cada etapa está en una componente y la componente de la etapa siguiente está contenida en la anterior. Sus elecciones izquierda/derecha proporcionan dígitos $a_n\in\{0,2\}$. La diferencia entre $x$ y la suma de sus primeros $n$ dígitos está entre cero y $3^{-n}$, por lo que $x$ es la serie correspondiente. No es necesario usar la expansión ternaria convencional en los extremos: por ejemplo, $1/3=0.02222\ldots$ usa sólo dígitos permitidos.

La correspondencia es inyectiva. Si dos sucesiones difieren por primera vez en $k$, la diferencia aportada allí tiene valor absoluto $2/3^k$, mientras que toda la cola posterior puede compensar a lo sumo $\sum_{j>k}2/3^j=1/3^k$. Las sumas difieren al menos en $1/3^k>0$. Finalmente, el conjunto de sucesiones de ceros y doses es incontable por el argumento diagonal: frente a una supuesta enumeración, cambiamos el dígito $n$ de la fila $n$. Su imagen inyectiva $C$ es, pues, incontable.
\end{solucion}

\begin{ejercicio}[Lectura 9, ejercicio 1]\label{ej:lectura-9-1}\label{ch09:unicidad}
Sean $(Y,i)$ y $(Y',i')$ dos completaciones de $X$. Demuestre que existe una isometría sobreyectiva $T:Y\to Y'$ tal que $T\circ i=i'$.
\end{ejercicio}
\begin{solucion}
Si $X=\varnothing$, la densidad fuerza $Y=Y'=\varnothing$, y la afirmación es inmediata. En otro caso, para $y\in Y$ elegimos $x_n\in X$ con $i(x_n)\to y$. Como $i$ es isometría, $(x_n)$ es de Cauchy en $X$, y $(i'(x_n))$ lo es en el espacio completo $Y'$. Definimos
\[
 T(y)=\lim_{n\to\infty}i'(x_n).
\]
Si otra sucesión $z_n$ satisface $i(z_n)\to y$, entonces
\[
 d(x_n,z_n)=\rho(i(x_n),i(z_n))
 \leq\rho(i(x_n),y)+\rho(y,i(z_n))\longrightarrow0.
\]
La distancia entre $i'(x_n)$ e $i'(z_n)$ es la misma; sus límites coinciden y $T$ está bien definida. Para $y,z\in Y$, elegimos aproximaciones $i(x_n)\to y$, $i(z_n)\to z$ y obtenemos
\[
 \rho'(T(y),T(z))=\lim_nd(x_n,z_n)=\rho(y,z).
\]
Aquí la continuidad de la distancia se justifica por $|\rho(a,b)-\rho(c,d)|\leq\rho(a,c)+\rho(b,d)$. Por tanto $T$ es una isometría. La sucesión constante $x_n=x$ muestra $T(i(x))=i'(x)$.

Para probar sobreyectividad, sea $y'\in Y'$. Elegimos $x_n\in X$ con $i'(x_n)\to y'$. Entonces $(i(x_n))$ es de Cauchy y converge en $Y$ a cierto $y$. Por la definición anterior, $T(y)=y'$. Finalmente, la isometría que cumple la identidad requerida es única: es continua y queda determinada en la imagen densa $i(X)$; si $S$ es otra, para $i(x_n)\to y$ tenemos $S(y)=\lim_ni'(x_n)=T(y)$.
\end{solucion}

\chapter{Categoría de Baire y principio de contracción}

\begin{ejercicio}[Lectura 10, ejercicio 1]\label{ej:lectura-10-1}
Determine si es necesario que las bolas del teorema de bolas anidadas sean cerradas.
\end{ejercicio}
\begin{solucion}
Sí, si se pretende esa conclusión bajo las restantes hipótesis. En $\R$, las bolas abiertas $B(1/n,1/n)=(0,2/n)$ están anidadas y sus radios tienden a cero, pero su intersección es vacía. Todo número positivo queda fuera cuando $n>2/x$, y cero no pertenece a ninguna bola. Las bolas cerradas correspondientes tienen intersección $\{0\}$. El papel de ser cerrado es conservar el límite de los centros.
\end{solucion}

\begin{ejercicio}[Lectura 10, ejercicio 2]\label{ej:lectura-10-2}
Enuncie y pruebe el recíproco del teorema de bolas cerradas anidadas.
\end{ejercicio}
\begin{solucion}
El recíproco afirma: si toda sucesión de bolas cerradas anidadas con radios tendientes a cero tiene intersección no vacía, entonces $X$ es completo. Sea $(x_n)$ de Cauchy. Para cada $k\geq1$ elegimos índices estrictamente crecientes $n_k$ de modo que
\[
 d(x_m,x_l)<2^{-k-2}\qquad(m,l\geq n_k).
\]
Podemos aumentar cualquier umbral para garantizar $n_{k+1}>n_k$. Pongamos $B_k=\overline B(x_{n_k},2^{-k})$. Si $z\in B_{k+1}$, entonces
\[
 d(z,x_{n_k})\leq2^{-k-1}+d(x_{n_{k+1}},x_{n_k})
 <2^{-k-1}+2^{-k-2}<2^{-k},
\]
por lo que $B_{k+1}\subset B_k$. La hipótesis proporciona $x\in\bigcap_kB_k$, y $d(x,x_{n_k})\leq2^{-k}\to0$. La sucesión original es de Cauchy y tiene esta subsucesión convergente; por la proposición~\ref{ch08:cauchy-basico}, converge a $x$.
\end{solucion}

\begin{ejercicio}[Lectura 10, ejercicio 3]\label{ej:lectura-10-3}
Pruebe que $X$ es completo si y sólo si toda sucesión de cerrados \emph{no vacíos} $A_{n+1}\subset A_n$ con $\diam(A_n)\to0$ tiene intersección no vacía.
\end{ejercicio}
\begin{solucion}
Se explicita la no vacuidad, necesaria porque la sucesión constante $A_n=\varnothing$ refutaría el enunciado sin ella. Si $X$ es completo, elegimos $a_n\in A_n$. Para $m,k\geq n$, ambos puntos pertenecen a $A_n$ y $d(a_m,a_k)\leq\diam(A_n)\to0$. La sucesión es de Cauchy y converge a $a\in X$. Para cada $n$, una cola está en el cerrado $A_n$, de donde $a\in A_n$. La intersección es incluso un singleton porque dos puntos suyos distarían a lo sumo $\diam(A_n)$ para todo $n$.

Para la vuelta, sea $(x_n)$ de Cauchy y pongamos $A_n=\cl{\{x_k:k\geq n\}}$. Son cerrados no vacíos y anidados. Si las distancias entre términos de la cola son menores que $\eps$, las distancias entre puntos de su clausura son a lo sumo $\eps$: se aproximan dos puntos de la clausura por términos de la cola y se pasa al límite en la desigualdad triangular. Por tanto $\diam(A_n)\to0$. Sea $x$ un punto de la intersección. Para $k\geq n$, tenemos $x,x_k\in A_n$, así que $d(x,x_k)\leq\diam(A_n)\to0$. Luego $x_n\to x$ y $X$ es completo.
\end{solucion}

\begin{ejercicio}[Lectura 10, ejercicio 4]\label{ej:lectura-10-4}
Pruebe que $A$ es denso en ninguna parte si y sólo si $X\setminus\cl A$ es denso en $X$.
\end{ejercicio}
\begin{solucion}
Un conjunto $D$ es denso si y sólo si todo abierto no vacío lo intersecta. Aplicado a $D=X\setminus\cl A$, esto dice que ningún abierto no vacío está contenido en $\cl A$, es decir, $\interior{\cl A}=\varnothing$. La clausura escrita en el ejercicio original es esencial: reemplazar $X\setminus\cl A$ por $X\setminus A$ sólo detectaría que $A$ tiene interior vacío.
\end{solucion}

\begin{ejercicio}[Lectura 10, ejercicio 5]\label{ej:lectura-10-5}
Demuestre que todo subconjunto de un conjunto magro es magro.
\end{ejercicio}
\begin{solucion}
Sea $A=\bigcup_nA_n$ con cada $A_n$ denso en ninguna parte, y sea $B\subset A$. Entonces $B=\bigcup_n(B\cap A_n)$. La inclusión $\cl{B\cap A_n}\subset\cl{A_n}$ implica $\interior{\cl{B\cap A_n}}\subset\interior{\cl{A_n}}=\varnothing$. Cada término es denso en ninguna parte y la unión es magra. El mismo argumento muestra que una unión numerable de conjuntos magros es magra, reenumerando los pares de índices mediante una enumeración de $\N\times\N$.
\end{solucion}

\begin{ejercicio}[Lectura 10, ejercicio 6]\label{ej:lectura-10-6}
Sea $X$ completo no vacío y sea $X=\bigcup_{n\geq1}F_n$, donde los $F_n$ son cerrados. Pruebe que alguno tiene interior no vacío.
\end{ejercicio}
\begin{solucion}
Si todos tuvieran interior vacío, por ser cerrados serían densos en ninguna parte. La igualdad expresaría $X$ como un conjunto magro en sí mismo, contra el corolario~\ref{ch10:segunda}. No se concluye que todos los $F_n$ tengan interior, sino que al menos uno debe tenerlo.
\end{solucion}

\begin{ejercicio}[Lectura 10, ejercicio 7]\label{ej:lectura-10-7}
Demuestre que un espacio de Banach de dimensión infinita posee una base algebraica y que toda base algebraica suya es incontable.
\end{ejercicio}
\begin{solucion}
Por base se entiende aquí \emph{base de Hamel}: cada vector es combinación lineal \emph{finita} de sus elementos. Primero justificamos la existencia mediante el lema de Zorn: un conjunto parcialmente ordenado en el que toda cadena admite una cota superior tiene un elemento maximal. Ordenamos por inclusión los subconjuntos linealmente independientes de $V$. La familia contiene al vacío. La unión de una cadena es independiente, porque cualquier relación lineal utiliza solo finitos vectores, y una de las familias de la cadena contiene a todos ellos. Esa unión es una cota superior. Zorn produce un conjunto independiente maximal $E$. Si existiera $v\notin\spanop E$, el conjunto $E\cup\{v\}$ seguiría siendo independiente: una relación con coeficiente no nulo de $v$ lo expresaría como combinación de $E$, y una relación con coeficiente cero contradice la independencia de $E$. Así $E$ genera $V$ y es una base.

Supongamos que $V$ tiene una base numerable $(v_n)$ y pongamos $E_n=\spanop\{v_1,\ldots,v_n\}$. Entonces $V=\bigcup_nE_n$.

Justificamos dos hechos. Primero, todo subespacio de dimensión finita es cerrado. Para una base $w_1,\ldots,w_m$ de tal subespacio, existe $c>0$ con
\[
 \left\|\sum_{j=1}^m a_jw_j\right\|\geq c\max_j|a_j|.
\]
Si no existiera, habría vectores de coeficientes $a^{(k)}$ con $\max_j|a_j^{(k)}|=1$ y cuyas combinaciones tendieran a cero. Por Bolzano--Weierstrass en dimensión finita, una subsucesión de coeficientes converge a $a$ con $\max_j|a_j|=1$. La desigualdad triangular muestra continuidad de la combinación lineal, así que $\sum_ja_jw_j=0$, contra la independencia. La cota probada implica que una sucesión convergente de combinaciones tiene coeficientes de Cauchy; éstos convergen y su combinación representa el límite. El subespacio es cerrado.

Segundo, un subespacio lineal propio no tiene interior. Si contuviera $B(a,r)$, al restar $a$ contendría $B(0,r)$. Para cualquier $v\in V$, un múltiplo $tv$ con $t>0$ suficientemente pequeño estaría en esa bola; al multiplicar por $1/t$, también $v$ pertenecería al subespacio, que ya no sería propio.

Cada $E_n$ es, pues, cerrado de interior vacío. Esto contradice Baire porque $V$ es completo y no vacío. El resultado no excluye bases de Schauder numerables, en las que se admiten series convergentes; ésa es otra noción de base.
\end{solucion}

\begin{ejercicio}[Lectura 10, ejercicio 8]\label{ej:lectura-10-8}
Sea $(q_n)$ una enumeración de $\Q$, y para $\eps>0$ defina
\[
 A_\eps=\bigcup_{n\geq1}(q_n-\eps2^{-n-1},q_n+\eps2^{-n-1}).
\]
Determine si $E=\bigcap_{\eps>0}A_\eps$ contiene puntos irracionales además de los racionales.
\end{ejercicio}
\begin{solucion}
Cada $A_\eps$ es abierto y contiene a $\Q$, por lo que es denso en $\R$. Como $A_\delta\subset A_\eps$ cuando $\delta\leq\eps$, se cumple
\[
 E=\bigcap_{m\geq1}A_{1/m}.
\]
Una inclusión es inmediata; para la otra, dado $\eps>0$ elegimos $m$ con $1/m\leq\eps$ y usamos la inclusión de los conjuntos. Por Baire, $E$ es denso. Más aún,
\[
 E\setminus\Q=\bigcap_{m\geq1}A_{1/m}\ \cap\ \bigcap_{n\geq1}(\R\setminus\{q_n\})
\]
es una intersección numerable de abiertos densos; por Baire es densa y, en particular, no vacía. Por tanto hay irracionales en $E$, incluso en todo intervalo abierto.

De hecho, $E$ es incontable. Su complemento es magro, por ser unión de los cerrados de interior vacío $\R\setminus A_{1/m}$. Si $E$ fuera numerable, también sería magro, y $\R$ sería unión de dos conjuntos magros, contra Baire. La suma de longitudes de los intervalos de $A_\eps$ es $\sum_n\eps2^{-n}=\eps$. Este dato anticipa una distinción importante para el siguiente volumen: ser grande por categoría y ser grande por longitud son nociones diferentes.
\end{solucion}

\begin{ejercicio}[Lectura 10, ejercicio 9]\label{ej:lectura-10-9}
\begin{enumerate}[label=(\roman*)]
\item Examine la afirmación original: «No existe una función $f:\R\to\R$ discontinua en todo número irracional». Formule y demuestre la versión correcta referente al conjunto exacto de discontinuidades.
\item Para cada conjunto a lo sumo numerable $A\subset\R$, construya una función cuyas discontinuidades sean exactamente los puntos de $A$.
\end{enumerate}
\end{ejercicio}
\begin{solucion}
\textbf{(i)} La afirmación literal es falsa: $\mathbf1_{\Q}$ es discontinua en todos los puntos, y por tanto también en todos los irracionales. La versión correcta es que ninguna función $\R\to\R$ tiene como conjunto de discontinuidades \emph{exactamente} $\R\setminus\Q$.

Para demostrarlo introducimos la oscilación de una función $f:X\to Y$ entre espacios métricos:
\[
 \omega_f(x)=\inf_{r>0}\diam f(B(x,r))\in[0,\infty].
\]
La función es continua en $x$ si y sólo si $\omega_f(x)=0$. En un sentido, la continuidad con tolerancia $\eps/2$ hace que la imagen de una bola tenga diámetro a lo sumo $\eps$. En el otro, si una bola tiene imagen de diámetro menor que $\eps$, entonces $\rho(f(y),f(x))<\eps$ para todos sus puntos $y$.

Para $a>0$, el conjunto $D_a=\{x:\omega_f(x)\geq a\}$ es cerrado. Si $x\notin D_a$, existe $r>0$ con $\diam f(B(x,r))<a$. Cuando $d(x,y)<r/2$, la bola $B(y,r/2)$ está contenida en $B(x,r)$ y $\omega_f(y)<a$. Así el complemento de $D_a$ es abierto. El conjunto de discontinuidades es, por tanto,
\[
 D(f)=\bigcup_{m\geq1}D_{1/m},
\]
una unión numerable de cerrados. Si fuera exactamente el conjunto de irracionales, cada cerrado $D_{1/m}$ carecería de interior, pues cualquier intervalo contiene racionales. Serían todos densos en ninguna parte. Al agregar los singletons racionales, expresaríamos $\R$ como una unión numerable de conjuntos densos en ninguna parte, contradicción con Baire.

\textbf{(ii)} Si $A=\varnothing$, sirve la función cero. Si $A$ es finito o numerable no vacío, enumeremos sus puntos distintos como $a_1,a_2,\ldots$ (deteniéndonos si es finito), y definamos
\[
 f(x)=\begin{cases}2^{-n},&x=a_n,\\0,&x\notin A.\end{cases}
\]
Si $x\notin A$ y $\eps>0$, elijamos $N$ con $2^{-N}<\eps$. Una bola suficientemente pequeña alrededor de $x$ evita los finitos puntos $a_1,\ldots,a_N$ que existan. En ella $|f(y)|<\eps$, por lo que $f$ es continua en $x$. En cambio, todo intervalo contiene puntos fuera de $A$, ya que un intervalo es incontable. Podemos elegir $y_k\notin A$ con $y_k\to a_n$, de modo que $f(y_k)=0\not\to2^{-n}=f(a_n)$. Así, $f$ es discontinua en cada $a_n$ y sólo allí.
\end{solucion}

\begin{ejercicio}[Lectura 10, ejercicio 10]\label{ej:lectura-10-10}
Construya funciones $\R\to\R$ discontinuas en todos los puntos de: \begin{enumerate}[label=(\roman*)]\item $\Q$;\item el conjunto de Cantor.\end{enumerate}
\end{ejercicio}
\begin{solucion}
En (i) podemos obtener incluso exactamente ese conjunto: enumeramos $\Q=\{q_n\}$ y ponemos $f(q_n)=2^{-n}$ y $f(x)=0$ para $x$ irracional. El ejercicio~\ref{ej:lectura-10-9} demuestra que sus discontinuidades son exactamente $\Q$.

En (ii) tomamos $g=\mathbf1_C$. Fuera de $C$, la función vale cero en una vecindad porque $C$ es cerrado. En un punto $x\in C$, existen puntos de $\R\setminus C$ arbitrariamente próximos. Para justificarlo, si un intervalo abierto estuviera contenido en $C$, también estaría contenido en cada $C_n$; como un intervalo conexo no puede atravesar los huecos de $C_n$, tendría que estar dentro de una componente de longitud $3^{-n}$ para todo $n$, imposible para un intervalo de longitud positiva. Por tanto $C$ tiene interior vacío. Podemos elegir $x_k\notin C$ con $x_k\to x$, y $g(x_k)=0$ mientras $g(x)=1$. Las discontinuidades de $g$ son exactamente $C$.
\end{solucion}

\begin{ejercicio}[Lectura 11, ejercicio 1]\label{ej:lectura-11-1}
Si $X$ es completo no vacío y $f:X\to X$ satisface que $f\circ f$ es una contracción, determine si $f$ tiene un punto fijo.
\end{ejercicio}
\begin{solucion}
La aplicación $T=f^2$ tiene un único punto fijo $a$. No necesitamos suponer que $f$ sea continua. Por asociatividad de la composición,
\[
 T(f(a))=f(f(f(a)))=f(T(a))=f(a).
\]
Así $f(a)$ también es punto fijo de $T$; por unicidad, $f(a)=a$. Si $b$ fuera otro punto fijo de $f$, sería punto fijo de $T$, luego $b=a$. En general, el mismo razonamiento funciona si alguna potencia $f^k$ es contracción: $f$ conmuta con sus potencias y todo punto fijo de $f$ es fijo para $f^k$.
\end{solucion}

\begin{ejercicio}[Lectura 11, ejercicio 2]\label{ej:lectura-11-2}
Usando el ejercicio~\ref{ej:lectura-11-1}, pruebe que $x=\cos x$ tiene una única solución en $[0,1]$.
\end{ejercicio}
\begin{solucion}
La aplicación $f(x)=\cos x$ lleva $[0,1]$ a $[\cos1,1]\subset[0,1]$. Su cuadrado por composición es $f^2(x)=\cos(\cos x)$ y
\[
 |(f^2)'(x)|=|\sin(\cos x)\sin x|\leq\sin1<1.
\]
El valor medio prueba que $f^2$ es contracción en el intervalo completo; el ejercicio~\ref{ej:lectura-11-1} da un único punto fijo de $f$. También se podría aplicar Banach directamente porque $|f'(x)|=\sin x\leq\sin1<1$. La formulación mediante $f^2$ ilustra que la propiedad contractiva puede verificarse tras varias iteraciones.
\end{solucion}

\begin{ejercicio}[Lectura 11, ejercicio 3]\label{ej:lectura-11-3}
Pruebe que $f(x)=x^2+x^{-1}$ no tiene puntos fijos en $[1,\infty)$.
\end{ejercicio}
\begin{solucion}
Para $x\geq1$, $x^2\geq x$ y $x^{-1}>0$, de donde $f(x)>x$. Por tanto $f(x)=x$ es imposible. La aplicación sí lleva el intervalo a sí mismo, pero no es una contracción: su derivada $2x-x^{-2}$ es no acotada. La invariancia de un espacio completo, sin contracción, no garantiza un punto fijo.
\end{solucion}

\begin{ejercicio}[Lectura 11, ejercicio 4]\label{ej:lectura-11-4}
Sean $0<r<1$ y $\alpha\in\R$. En $C([0,r])$ defina
\[
 (\Phi f)(x)=\alpha+\int_0^x f(t)\,dt.
\]
Pruebe que $\Phi$ es contracción para la norma uniforme y encuentre su punto fijo.
\end{ejercicio}
\begin{solucion}
La integral de una función continua tiene una primitiva continua, por lo que $\Phi$ toma valores en el espacio indicado. Además,
\[
 |(\Phi f)(x)-(\Phi g)(x)|
 \leq\int_0^x|f(t)-g(t)|\,dt
 \leq r\norm{f-g}_\infty.
\]
Tomando el supremo resulta $\norm{\Phi f-\Phi g}_\infty\leq r\norm{f-g}_\infty$, con $r<1$.

Justificamos la completitud del dominio sin anticipar resultados: una sucesión de Cauchy en norma uniforme tiene un límite puntual porque $\R$ es completo. Pasando al límite en una de las dos funciones de la estimación uniforme de Cauchy, se obtiene convergencia uniforme al límite. Ese límite es continuo: dado $x$ y $\eps>0$, elegimos una función de la sucesión a distancia uniforme menor que $\eps/3$ del límite, y usamos su continuidad en $x$ con tolerancia $\eps/3$; la desigualdad triangular suma tres errores menores que $\eps/3$. La función límite es acotada por una función de la sucesión más un error uniforme finito. Esto prueba que $C([0,r])$ es completo.

Banach proporciona un único punto fijo. La función $f(x)=\alpha e^x$ lo es, pues
\[
 \alpha+\int_0^x\alpha e^t\,dt=\alpha+\alpha(e^x-1)=\alpha e^x.
\]
También se descubre la solución diferenciando la ecuación integral: un punto fijo satisface $f'=f$ y $f(0)=\alpha$; multiplicar por $e^{-x}$ muestra que la derivada de $e^{-x}f(x)$ es cero.
\end{solucion}

\begin{ejercicio}[Lectura 11, ejercicio 5]\label{ej:lectura-11-5}\label{ch11:hiper-completo}
Demuestre que $(\mathcal F(X),H)$ es completo si y sólo si $X$ es completo.
\end{ejercicio}
\begin{solucion}
Supondremos primero $X$ completo y no vacío. Sea $(A_n)$ de Cauchy para $H$. Las funciones $u_n(x)=\dist(x,A_n)$ satisfacen
\[
 |u_n(x)-u_m(x)|\leq H(A_n,A_m)
\]
para todo $x$. Por completitud de $\R$, existe $u(x)=\lim_nu_n(x)$, y la estimación de Cauchy, al pasar $m\to\infty$, demuestra que la convergencia es uniforme en todo $X$. No es necesario que las funciones $u_n$ sean acotadas individualmente. Cada una es $1$-Lipschitz; al tomar límites en $|u_n(x)-u_n(y)|\leq d(x,y)$ vemos que $u$ también lo es. Por tanto
\[
 A=\{x\in X:u(x)=0\}
\]
es cerrado. Demostraremos que es no vacío y acotado, y que $H(A_n,A)\to0$.

\textbf{Construcción de un punto límite.} Seleccionemos índices crecientes $n_k$ con $H(A_{n_k},A_{n_{k+1}})<2^{-k-1}$. Elegimos $a_1\in A_{n_1}$ y sucesivamente $a_{k+1}\in A_{n_{k+1}}$ con $d(a_k,a_{k+1})<2^{-k}$. Tal elección es posible por la definición de distancia a un conjunto y la cota estricta anterior. Como $\sum_k2^{-k}<\infty$, las colas de la sucesión satisfacen $d(a_l,a_k)\leq\sum_{j=k}^{l-1}2^{-j}$; es de Cauchy y converge a algún $a\in X$. La convergencia uniforme $u_{n_k}\to u$ y $u_{n_k}(a_k)=0$ implican
\[
 |u(a)|\leq |u(a)-u(a_k)|+|u(a_k)-u_{n_k}(a_k)|\longrightarrow0.
\]
Así $a\in A$, y $A$ no es vacío.

\textbf{Acotación.} Para algún índice $N$ y alguna constante finita $M$, la convergencia uniforme da $|u_N(x)-u(x)|\leq M$ para todo $x$. Si $x\in A$, entonces $\dist(x,A_N)\leq M$. Como $A_N$ es acotado, elegimos $a_0\in A_N$ y $R$ con $d(a_0,b)\leq R$ para todo $b\in A_N$. Para cualquier $x\in A$ existe $b\in A_N$ con $d(x,b)<M+1$, así que $d(x,a_0)<M+1+R$. Esto prueba que $A\in\mathcal F(X)$.

\textbf{Convergencia en ambas direcciones.} Dado $\eps>0$, elegimos $N$ tal que $H(A_n,A_m)<\eps$ para $n,m\geq N$. Al dejar $m\to\infty$ en la desigualdad para las funciones de distancia, obtenemos $|u_n(x)-u(x)|\leq\eps$ para todo $x$ y $n\geq N$. En particular,
\[
 \sup_{a\in A}\dist(a,A_n)\leq\eps.
\]
Para la dirección restante fijemos $n\geq N$ y $b\in A_n$. Construiremos una cadena que empieza en $b$. Para cada $k\geq0$ elijamos umbrales crecientes $N_k$, con $N_0=N$, tales que $H(A_l,A_m)<\eps2^{-k}$ para $l,m\geq N_k$. Pongamos $n_0=n$, y para $k\geq1$ escojamos índices crecientes $n_k\geq N_k$. Entonces
\[
 H(A_{n_k},A_{n_{k+1}})<\eps2^{-k}\quad(k\geq0).
\]
Partiendo de $b_0=b$, elegimos $b_{k+1}\in A_{n_{k+1}}$ con $d(b_k,b_{k+1})<2\eps2^{-k}$. La suma de estas cotas es $4\eps$, de modo que $(b_k)$ es de Cauchy y converge a un punto $b_\infty$ con $d(b,b_\infty)\leq4\eps$. El mismo argumento con $u_{n_k}$ usado antes demuestra $b_\infty\in A$. Así $\dist(b,A)\leq4\eps$, uniformemente para $b\in A_n$. Concluimos $H(A_n,A)\leq4\eps$ para $n\geq N$. Como $\eps$ es arbitrario, hay convergencia.

Recíprocamente, supongamos completo $\mathcal F(X)$ y sea $(x_n)$ de Cauchy en $X$. La identidad $H(\{x_n\},\{x_m\})=d(x_n,x_m)$ muestra que los singletons convergen en $\mathcal F(X)$ a algún $A\ne\varnothing$. Para $a,b\in A$,
\[
 d(a,b)\leq d(a,x_n)+d(x_n,b)\leq2H(A,\{x_n\})\longrightarrow0.
\]
Así $A$ tiene un solo punto, digamos $x$, y $d(x_n,x)=H(\{x_n\},A)\to0$. Por tanto $X$ es completo. Si $X=\varnothing$, también $\mathcal F(X)=\varnothing$ y ambos espacios son completos por la definición; la equivalencia sigue siendo válida.
\end{solucion}

\begin{ejercicio}[Lectura 11, ejercicio 6]\label{ej:lectura-11-6}
Calcule por iteración la solución del sistema
\[
 5x^2-y^2=0,\qquad y-\tfrac14(\sin x+\cos y)=0
\]
que se encuentra en la región positiva próxima a $(1/4,1/4)$.
\end{ejercicio}
\begin{solucion}
Para la rama positiva $y=\sqrt5\,x$, tomamos
\[
 T(x,y)=\left(\frac{y}{\sqrt5},\frac{\sin x+\cos y}{4}\right)
\]
en el rectángulo completo $R=[0,1/2]^2$. La primera coordenada está entre cero y $1/(2\sqrt5)<1/2$. La segunda está entre $\cos(1/2)/4>0$ y $(\sin(1/2)+1)/4<1/2$. Por tanto $T(R)\subset R$. Las sumas de valores absolutos de sus derivadas por fila son $1/\sqrt5<1/2$ y $(|\cos x|+|\sin y|)/4\leq(1+\sin(1/2))/4<1/2$. Podemos usar $q=1/2$ como constante de contracción.

La iteración simultánea, empezando en $(x_0,y_0)=(0.25,0.25)$, es
\[
 x_{n+1}=y_n/\sqrt5,\qquad y_{n+1}=(\sin x_n+\cos y_n)/4.
\]
En ambas fórmulas se usan los valores de la iteración anterior. Algunos valores calculados son:
\begin{center}
\begin{tabular}{rrr}
\toprule
$n$&$x_n$&$y_n$\\\midrule
0&0.250000000000&0.250000000000\\
1&0.111803398875&0.304079095241\\
2&0.135988305499&0.266423427535\\
3&0.119148178954&0.275072068197\\
5&0.120889895207&0.271597989694\\
10&0.121245551908&0.271102661689\\
15&0.121241889930&0.271105174792\\
20&0.121241911634&0.271105155669\\\bottomrule
\end{tabular}
\end{center}
La cota a priori da
\[
 \norm{(x_n,y_n)-(x_*,y_*)}_\infty
 \leq 2^{1-n}\norm{(x_1,y_1)-(x_0,y_0)}_\infty
 <0.277\,2^{-n}.
\]
En la iteración $20$, el error garantizado es menor que $2.65\cdot10^{-7}$, sin confundir los decimales mostrados con todos los decimales certificados. Así la solución es aproximadamente $(0.121242,0.271105)$, única dentro del rectángulo para la rama elegida. Sustituir el punto fijo en las dos fórmulas demuestra que resuelve el sistema original.
\end{solucion}

\chapter{Compacidad en espacios métricos}

\begin{ejercicio}[Lectura 13, ejercicio 1]\label{ej:lectura-13-1}
Una familia no vacía $\mathcal F$ tiene la propiedad de intersección finita cuando toda intersección de una cantidad finita de sus miembros es no vacía. Probar que $X$ es compacto si y solo si toda familia de cerrados con esta propiedad tiene intersección no vacía.
\end{ejercicio}
\begin{solucion}
Si $X$ es compacto y $\bigcap\mathcal F=\varnothing$, los complementos $X\setminus F$ forman una cubierta abierta. Una subcubierta finita implicaría $F_1\cap\cdots\cap F_m=\varnothing$, contradicción. En sentido inverso, si una cubierta abierta no tuviera subcubierta finita, los complementos de sus miembros serían cerrados con intersecciones finitas no vacías. Su intersección total, sin embargo, sería vacía. Si $X=\varnothing$, ambas propiedades se cumplen; no existe en él una familia de cerrados no vacíos con la propiedad indicada.
\end{solucion}

\begin{ejercicio}[Lectura 13, ejercicio 2]\label{ej:lectura-13-2}
Probar que una unión finita de subconjuntos compactos es compacta.
\end{ejercicio}
\begin{solucion}
Sea $K=K_1\cup\cdots\cup K_m$ y sea $\mathcal U$ una cubierta por abiertos del espacio ambiente. Para cada $i$, la compacidad de $K_i$ proporciona una subfamilia finita $\mathcal U_i\subset\mathcal U$ que lo cubre. La unión de las $m$ subfamilias es finita y cubre $K$. No se exige que los compactos sean disjuntos.
\end{solucion}

\begin{ejercicio}[Lectura 13, ejercicio 3]\label{ej:lectura-13-3}
En un compacto $X$, una sucesión decreciente de cerrados no vacíos $F_n$ tiene intersección no vacía. Deducir que $X$ es completo.
\end{ejercicio}
\begin{solucion}
La intersección de finitos $F_n$ coincide con el que tenga el índice mayor, así que la familia tiene la propiedad de intersección finita. Aplicamos el ejercicio \ref{ej:lectura-13-1}. Si $(x_n)$ es de Cauchy, consideramos
\[
F_n=\cl{\{x_k:k\ge n\}}.
\]
Son cerrados, no vacíos y decrecientes; elegimos $x\in\bigcap_nF_n$. Dado $\eps>0$, existe $N$ con $d(x_j,x_k)<\eps/2$ cuando $j,k\ge N$. Como $x\in F_N$, hay $k\ge N$ con $d(x,x_k)<\eps/2$. Entonces, para todo $j\ge N$, $d(x_j,x)<\eps$. Por tanto $x_n\to x$ y $X$ es completo.
\end{solucion}

\begin{ejercicio}[Lectura 13, ejercicio 4]\label{ej:lectura-13-4}
Probar que la bola unitaria cerrada de $\ell^\infty$ es cerrada y acotada, pero no compacta.
\end{ejercicio}
\begin{solucion}
La norma es continua porque $|\norm{x}_\infty-\norm{y}_\infty|\le\norm{x-y}_\infty$. Así la bola $K=\{x:\norm{x}_\infty\le1\}$ es cerrada y tiene diámetro a lo sumo $2$. Sean $e_n$ las sucesiones con un único $1$, en la coordenada $n$. Para $n\ne m$, $\norm{e_n-e_m}_\infty=1$. Las bolas abiertas de radio $1/3$ centradas en todos los puntos de $K$ forman una cubierta relativa. Cada bola contiene a lo sumo un $e_n$, pues dos puntos de la misma bola están a distancia menor que $2/3$. Ninguna subfamilia finita cubre todos los $e_n$. La misma separación ocurre entre cualesquiera dos sucesiones distintas de ceros y unos, como sugiere la lectura.
\end{solucion}

\begin{ejercicio}[Lectura 13, ejercicio 5]\label{ej:lectura-13-5}
Si $X$ es compacto y $f:X\to Y$ es continua e inyectiva, probar que $f^{-1}:f(X)\to X$ es continua.
\end{ejercicio}
\begin{solucion}
Para todo cerrado $F\subset X$, el conjunto $F$ es compacto, su imagen $f(F)$ es compacta y, por ser $Y$ métrico, es cerrada en $Y$ y en $f(X)$. Pero $f(F)$ es precisamente la preimagen de $F$ por $f^{-1}$. La caracterización de continuidad mediante cerrados prueba el resultado. La compacidad convierte aquí una biyección continua en un homeomorfismo.
\end{solucion}

\begin{ejercicio}[Lectura 13, ejercicio 6]\label{ej:lectura-13-6}
Sean $X$ compacto, $f:X\to Y$ continua y $A\subset X$ tal que $f(X\setminus A)\subset Y\setminus f(A)$. Probar: (i) $f|_A:A\to f(A)$ es cerrada; (ii) si además es inyectiva, es un homeomorfismo.
\end{ejercicio}
\begin{solucion}
La hipótesis dice que ningún punto exterior a $A$ comparte imagen con un punto de $A$; equivalentemente, $f^{-1}(f(A))=A$. Para (i), sea $F$ cerrado en $A$, de modo que $F=A\cap C$ con $C$ cerrado en $X$. Afirmamos
\[
f(F)=f(C)\cap f(A).
\]
Una inclusión es inmediata por las dos pertenencias. Si $y=f(c)\in f(C)\cap f(A)$, la hipótesis impide $c\notin A$; luego $c\in C\cap A=F$. Ahora $C$ es compacto y $f(C)$ cerrado en $Y$, de modo que $f(F)$ es cerrado en $f(A)$. Para (ii), la restricción es continua y biyectiva sobre su imagen. Ser cerrada significa que su inversa tiene preimagen cerrada para cada cerrado de $A$, por lo que la inversa es continua. No hemos supuesto que $A$ sea cerrado ni compacto.
\end{solucion}

\begin{ejercicio}[Lectura 13, ejercicio 7]\label{ej:lectura-13-7}
Si $X,Y$ son compactos, probar que $X\times Y$, con la métrica máximo, es compacto.
\end{ejercicio}
\begin{solucion}
Los rectángulos $B_X(x,r)\times B_Y(y,s)$ son abiertos y forman una base: contienen y son contenidos localmente en bolas de la métrica máximo. Sea $\mathcal U$ una cubierta abierta del producto. Para cada $(x,y)$ elegimos un rectángulo $V_{xy}\times W_{xy}$ que contenga $(x,y)$ y esté contenido en algún $U_{xy}\in\mathcal U$. Fijado $x$, los $W_{xy}$ cubren $Y$; elegimos $y_1,\ldots,y_{m_x}$ y ponemos $V_x=\bigcap_jV_{xy_j}$. Este es un abierto que contiene $x$, y
\[
V_x\times Y\subset\bigcup_{j=1}^{m_x}U_{xy_j}.
\]
Los $V_x$ cubren $X$; una cantidad finita $V_{x_1},\ldots,V_{x_k}$ basta. Los conjuntos $U_{x_i y_j}$ seleccionados para estos $x_i$ constituyen una subcubierta finita del producto. Si uno de los factores es vacío, el producto vacío ya es compacto.
\end{solucion}

\begin{ejercicio}[Lectura 13, ejercicio 8]\label{ej:lectura-13-8}
Probar que toda cubierta abierta $\mathcal U$ de un compacto métrico admite un número de Lebesgue: existe $r>0$ tal que todo subconjunto de diámetro menor que $r$ está contenido en algún miembro de $\mathcal U$.
\end{ejercicio}
\begin{solucion}
Suponemos $X\ne\varnothing$; para el espacio vacío el resultado se interpreta para subconjuntos no vacíos, o para cubiertas no vacías. Para cada $x$ elegimos $U_x\in\mathcal U$ y $r_x>0$ con $B(x,r_x)\subset U_x$. Una cantidad finita $B(x_i,r_i/2)$ cubre $X$; fijemos $r=\frac12\min_i r_i$. Si $A\ne\varnothing$ y $\diam A<r$, tomamos $a\in A$ e $i$ con $d(a,x_i)<r_i/2$. Para todo $b\in A$,
\[
d(b,x_i)\le d(b,a)+d(a,x_i)<r+r_i/2\le r_i.
\]
Luego $A\subset U_{x_i}$. El conjunto vacío está contenido en cualquier miembro de la cubierta. Un número de Lebesgue controla conjuntos completos de diámetro pequeño, no solamente bolas con un centro preestablecido.\index{numero de Lebesgue@número de Lebesgue}
\end{solucion}

\begin{ejercicio}[Lectura 14, ejercicio 1]\label{ej:lectura-14-1}
Si cada $(X_n,d_n)$ es compacto, probar que $\prod_{n\ge1}X_n$ es compacto para la métrica
\[
D(x,y)=\sum_{n=1}^\infty 2^{-n}\frac{d_n(x_n,y_n)}{1+d_n(x_n,y_n)}.
\]
\end{ejercicio}
\begin{solucion}
Si algún factor es vacío, el producto es vacío. En el otro caso, la serie converge porque cada sumando es menor que $2^{-n}$. La función $h(t)=t/(1+t)$ es creciente y subaditiva: para $a,b\ge0$,
\[
\frac{a+b}{1+a+b}=\frac a{1+a+b}+\frac b{1+a+b}
\le\frac a{1+a}+\frac b{1+b}.
\]
Por ello cada $h\circ d_n$ satisface la desigualdad triangular y también la suma $D$. La positividad, simetría y separación de puntos se comprueban término a término.

Sea $(x^{(k)})$ una sucesión del producto. Extraemos una subsucesión para la cual converge la primera coordenada; de esta, otra para la segunda; continuamos inductivamente. Escribamos $k_j^{(n)}$ los índices de la subsucesión elegida en el paso $n$, ordenados crecientemente, con cada lista contenida en la anterior. Los índices diagonales $k_n^{(n)}$ son estrictamente crecientes: el $(n+1)$-ésimo índice de la lista siguiente es al menos el $(n+1)$-ésimo de la anterior, y este supera su $n$-ésimo. Para cada coordenada fija $i$, la cola de la diagonal pertenece a la subsucesión del paso $i$ y converge a un punto $x_i\in X_i$.

Sea $x=(x_i)$. Dado $\eps>0$, elegimos $N$ con $\sum_{i>N}2^{-i}<\eps/2$. En las primeras $N$ coordenadas la convergencia permite hacer que la suma finita de distancias transformadas sea menor que $\eps/2$. La cola ya está controlada independientemente del índice. Así $D(x^{(k_n^{(n)})},x)<\eps$ para $n$ grande. El producto es secuencialmente compacto y el teorema \ref{ch14:equivalencias} concluye la prueba. Esta estimación prueba también que la convergencia en $D$ equivale a la convergencia en cada coordenada: la otra dirección se obtiene acotando cada sumando por $D$.
\end{solucion}

\begin{ejercicio}[Lectura 14, ejercicio 2]\label{ej:lectura-14-2}
Sea $X$ compacto y $U\subset X\times X$ abierto que contiene la diagonal $\Delta_X$. Probar que existe $r>0$ tal que $d(x,y)<r$ implica $(x,y)\in U$.
\end{ejercicio}
\begin{solucion}
Si no existiera tal $r$, elegiríamos $x_n,y_n$ con $d(x_n,y_n)<1/n$ y $(x_n,y_n)\notin U$. Una subsucesión $x_{n_j}$ converge a $x$; la desigualdad $d(y_{n_j},x)\le1/n_j+d(x_{n_j},x)$ muestra que $y_{n_j}\to x$. Entonces $(x_{n_j},y_{n_j})\to(x,x)\in U$. Como $U$ es abierto, la sucesión pertenece eventualmente a $U$, contradicción. El argumento convierte cercanía a la diagonal en una cota uniforme.
\end{solucion}

\begin{ejercicio}[Lectura 14, ejercicio 3]\label{ej:lectura-14-3}
Probar que todo espacio métrico no compacto contiene un subconjunto infinito numerable, cerrado y discreto.
\end{ejercicio}
\begin{solucion}
Existe una sucesión sin subsucesiones convergentes. Ningún valor aparece infinitas veces; su recorrido es infinito. Elegimos inductivamente una subsucesión de valores distintos y llamamos $A$ a su recorrido. Si $A$ tuviera un punto de acumulación, el lema \ref{ch14:acumulacion} permitiría extraer, respetando el orden de los índices, una subsucesión convergente, contradicción. Por tanto $A$ no tiene puntos de acumulación. Todo punto de $\cl A\setminus A$ sería uno de acumulación; así $A$ es cerrado. Para cada $a\in A$, que $a$ no sea de acumulación significa que existe $r_a>0$ con $B(a,r_a)\cap A=\{a\}$. Esta es precisamente la discreción de la topología relativa.
\end{solucion}

\begin{ejercicio}[Lectura 14, ejercicio 4]\label{ej:lectura-14-4}
Probar que $X$ es compacto si y solo si toda función continua $X\to\R$ es acotada.
\end{ejercicio}
\begin{solucion}
La dirección directa se sigue de que las imágenes compactas en $\R$ son acotadas. Para la inversa, si $X$ no es compacto, el ejercicio \ref{ej:lectura-14-3} proporciona $A=\{a_n:n\ge1\}$ cerrado y discreto. Construiremos directamente una función continua en todo $X$ que tome el valor $n$ en $a_n$. La discreción permite separar los centros; la cerradura impedirá que los soportes se acumulen fuera de $A$.

Elegimos
\[
0<r_n<\min\{1/n,\tfrac13\dist(a_n,A\setminus\{a_n\})\},\qquad
\varphi_n(x)=\max\{0,1-d(x,a_n)/r_n\}.
\]
La distancia del centro al resto es positiva por discreción. El \emph{soporte} de una función $\varphi$ es $\supp\varphi=\cl{\{x:\varphi(x)\ne0\}}$. La familia de soportes $\supp\varphi_n\subset\overline B(a_n,r_n)$ es localmente finita: cada punto tiene una vecindad que toca solo finitos de ellos. En efecto, si $x\notin A$, $\delta=\dist(x,A)>0$; para $n$ grande, $r_n<\delta/3$, y esas bolas no tocan $B(x,\delta/3)$. Si $x=a_k$, poniendo $\delta=\dist(a_k,A\setminus\{a_k\})>0$, toda bola con $n\ne k$ queda a distancia mayor que $2\delta/3$ de $a_k$, ya que $r_n<d(a_n,a_k)/3$. Por tanto
\[
f(x)=\sum_{n=1}^\infty n\varphi_n(x)
\]
es localmente una suma finita de funciones continuas, luego es continua. Además $f(a_n)=n$. Es una función continua no acotada, contradiciendo la hipótesis.

Una vez demostrado el teorema de Tietze \ref{ch12:tietze}, dispondremos de otra prueba de este paso: la función $a_n\mapsto n$ es continua en el cerrado discreto $A$, y la versión real no necesariamente acotada de dicho teorema permite extenderla continuamente a $X$. La construcción anterior demuestra el resultado sin depender de ese teorema posterior.
\end{solucion}

\begin{ejercicio}[Lectura 14, ejercicio 5]\label{ej:lectura-14-5}
Sea $1\le p<\infty$. Probar que un subconjunto cerrado y acotado $A\subset\ell^p$ es compacto si y solo si
\[
\forall\eps>0\ \exists N\ \forall x\in A:\quad
\sum_{k=N}^\infty |x_k|^p<\eps.
\]
Deducir que $K=\{x\in\ell^2:0\le x_n\le1/n\}$ es compacto.
\end{ejercicio}
\begin{solucion}
\emph{Necesidad.} Sea $A$ compacto. Fijamos $\eps>0$ y elegimos $\delta>0$ con $2^p\delta^p<\eps$. Finitas bolas de radio $\delta$, centradas en $a^{(1)},\ldots,a^{(m)}\in A$, cubren $A$. Cada centro pertenece a $\ell^p$, por lo que para un $N$ común a todos ellos sus colas tienen suma menor que $\delta^p$. Dado $x\in A$, elegimos un centro a distancia menor que $\delta$ y usamos $|u+v|^p\le2^{p-1}(|u|^p+|v|^p)$:
\[
\sum_{k=N}^\infty|x_k|^p
\le2^{p-1}\left(\norm{x-a^{(j)}}_p^p+
\sum_{k=N}^\infty|a_k^{(j)}|^p\right)<2^p\delta^p<\eps.
\]

\emph{Suficiencia.} Sea $(x^{(j)})$ una sucesión en $A$ y $M$ una cota para sus normas. Como cada coordenada está acotada por $M$, el procedimiento diagonal del ejercicio \ref{ej:lectura-14-1}, usando Bolzano--Weierstrass en $\R$ o $\C$, produce una subsucesión que converge coordenada a coordenada a $x=(x_k)$. Para todo $m$,
\[
\sum_{k=1}^m|x_k|^p=\lim_j\sum_{k=1}^m|x_k^{(j)}|^p\le M^p.
\]
Las sumas parciales son crecientes y acotadas, por lo que $x\in\ell^p$. Dado $\eta>0$, elegimos $N$ de modo que las colas de todos los elementos de $A$ sean menores que $\eta$. Pasando al límite en cada suma finita de la cola y luego tomando el supremo de esas sumas obtenemos $\sum_{k=N}^\infty|x_k|^p\le\eta$. Así
\[
\norm{x^{(j)}-x}_p^p
\le\sum_{k<N}|x_k^{(j)}-x_k|^p+2^p\eta.
\]
Primero hacemos $\eta$ tan pequeño como necesitemos y después hacemos pequeña la suma finita mediante la convergencia coordenada. Se obtiene convergencia en norma. Puesto que $A$ es cerrado, $x\in A$, y la compacidad secuencial prueba la compacidad.

Para $K$, cada condición $0\le x_n\le1/n$ es cerrada porque la evaluación de coordenada es continua. Además $\norm{x}_2^2\le\sum_n1/n^2<\infty$, y las colas están uniformemente acotadas por $\sum_{k\ge N}1/k^2\to0$. El criterio se aplica. La condición de colas es la parte de la compacidad que la mera acotación en dimensión infinita no detecta.\index{cubo de Hilbert}
\end{solucion}

\begin{ejercicio}[Lectura 14, ejercicio 6]\label{ej:lectura-14-6}
Sea $W$ una vecindad compacta de $x$ y $f:W\to[0,\infty)$ continua, con $f(x)=0$ y $f(y)>0$ para $y\in W\setminus\{x\}$. Probar que para toda vecindad $V$ de $x$ existe $r>0$ tal que $\{y\in W:f(y)<r\}\subset V$.
\end{ejercicio}
\begin{solucion}
El subnivel debe tomarse en $W$, que es el dominio de $f$; escribir $y\in X$ sin extender $f$ no define un conjunto. Elegimos un abierto $O$ con $x\in O\subset V$. El conjunto $W\setminus O$ es compacto. Si es vacío, cualquier $r>0$ funciona. Si no, $f$ alcanza en él un mínimo $m$. Ninguno de sus puntos es $x$, luego $m>0$. Tomamos $r=m$. Un punto de $W$ con $f(y)<r$ no puede estar en $W\setminus O$, y por ello pertenece a $O\subset V$. El argumento explica por qué la positividad punto a punto se vuelve una cota positiva uniforme lejos del único cero.
\end{solucion}

\begin{ejercicio}[Lectura 15, ejercicio 1]\label{ej:lectura-15-1}
Si $A\subset X$ es totalmente acotado, probar que su cerradura también lo es.
\end{ejercicio}
\begin{solucion}
Fijado $\eps>0$, cubrimos $A$ con finitas bolas $B(a_i,\eps/2)$. Por ser la unión finita, su cerradura está contenida en la unión de las cerraduras de esas bolas, y cada una de estas está contenida en $\overline B(a_i,\eps/2)\subset B(a_i,\eps)$. En detalle, si un punto estuviera fuera de todas las cerraduras, la intersección de las finitas vecindades que evitan cada bola sería una vecindad que evita su unión. Por tanto las bolas $B(a_i,\eps)$ cubren $\cl A$.
\end{solucion}

\begin{ejercicio}[Lectura 15, ejercicio 2]\label{ej:lectura-15-2}
Probar que el recorrido de toda sucesión de Cauchy es totalmente acotado.
\end{ejercicio}
\begin{solucion}
Para $\eps>0$, elegimos $N$ con $d(x_n,x_m)<\eps$ si $n,m\ge N$. La cola está en $B(x_N,\eps)$, y los términos anteriores quedan cubiertos por $B(x_1,\eps),\ldots,B(x_{N-1},\eps)$. Esta es una cubierta finita del recorrido. No necesitamos que la sucesión converja ni que el ambiente sea completo.
\end{solucion}

\begin{ejercicio}[Lectura 15, ejercicio 3]\label{ej:lectura-15-3}
Probar que $X$ es localmente compacto si y solo si cada $x$ admite $r_x>0$ con la bola cerrada $\overline B(x,r_x)$ compacta.
\end{ejercicio}
\begin{solucion}
Si $K$ es una vecindad compacta de $x$, existe $s>0$ con $B(x,s)\subset K$. Tomamos $r_x=s/2$; la bola cerrada de ese radio es un cerrado de $K$, luego compacta. Para la inversa, $\overline B(x,r_x)$ contiene la bola abierta $B(x,r_x)$ y por tanto es una vecindad compacta. Se trata de la bola cerrada, no de afirmar que la bola abierta sea compacta.
\end{solucion}

\begin{ejercicio}[Lectura 15, ejercicio 4]\label{ej:lectura-15-4}
Probar que $\Q$, con la métrica inducida de $\R$, no es localmente compacto.
\end{ejercicio}
\begin{solucion}
Supongamos que $q\in\Q$ tiene una vecindad compacta $K\subset\Q$. Entonces $\Q\cap(q-r,q+r)\subset K$ para algún $r>0$. Elegimos un irracional $\alpha\in(q-r,q+r)$ y racionales $q_n$ en ese intervalo que converjan en $\R$ a $\alpha$. Todos pertenecen a $K$. Una subsucesión debería converger en $\Q$ a algún $a\in K$; la inclusión continua $\Q\hookrightarrow\R$ daría convergencia a $a$ en $\R$, contradiciendo la unicidad del límite $\alpha\notin\Q$. Así una subparte arbitraria de un espacio localmente compacto puede dejar de ser localmente compacta. Las hipótesis abierto o cerrado de la proposición \ref{ch15:subespacios} son suficientes; no se afirma que sean individualmente necesarias para cada subespacio.
\end{solucion}

\begin{ejercicio}[Lectura 15, ejercicio 5]\label{ej:lectura-15-5}
Sean $X$ localmente compacto, $A\subset X$ compacto y $U\subset X$ abierto con $A\subset U$. Probar:
\begin{enumerate}[label=(\roman*)]
\item existen $V$ abierto y $B$ compacto con $A\subset V\subset B\subset U$;
\item existe $f:X\to[0,1]$ continua, igual a $1$ en $A$, cuyo soporte es compacto y está contenido en $U$.
\end{enumerate}
\end{ejercicio}
\begin{solucion}
Para (i), cada $a\in A$ admite una bola cerrada compacta de radio $s_a>0$. Como $U$ es abierto, podemos disminuir el radio hasta elegir $r_a>0$ con $\overline B(a,2r_a)\subset U$ y con esa bola todavía compacta, por ser cerrada en la bola original. Extraemos una subcubierta finita $B(a_i,r_i)$ de $A$ y ponemos
\[
V=\bigcup_iB(a_i,r_i),\qquad B=\bigcup_i\overline B(a_i,r_i).
\]
Entonces $\cl V\subset B\subset U$ y $B$ es compacto.

Para (ii), si $A=\varnothing$, elegimos $f=0$. Para $A\ne\varnothing$ y $X\setminus V\ne\varnothing$, definimos
\[
f(x)=\frac{\dist(x,X\setminus V)}{\dist(x,X\setminus V)+\dist(x,A)}.
\]
Ambas distancias son continuas. El denominador nunca es cero: si $\dist(x,A)=0$, el cerrado $A$ contiene $x$, y entonces $x\in V$ tiene distancia positiva al cerrado $X\setminus V$. La función toma valores en $[0,1]$, vale $1$ en $A$ y vale $0$ fuera de $V$. Por tanto $\supp f=\cl{\{x:f(x)\ne0\}}\subset\cl V\subset B\subset U$; es cerrado en el compacto $B$, luego compacto. Si $V=X$, las inclusiones de (i) implican $B=X$ compacto y podemos elegir $f=1$. Esta construcción explica por qué se debe reservar espacio entre $A$ y el exterior de $U$.\index{funcion de corte@función de corte}\index{soporte}
\end{solucion}

\begin{ejercicio}[Lectura 15, ejercicio 6]\label{ej:lectura-15-6}
Si el compacto $A$ está cubierto por abiertos $U_1,\ldots,U_n$ de un espacio localmente compacto $X$, construir funciones continuas $f_k:X\to[0,1]$ de soporte compacto contenido en $U_k$, tales que $\sum_k f_k\le1$ en $X$ y $\sum_kf_k=1$ en $A$.
\end{ejercicio}
\begin{solucion}
Si $A$ es vacío, tomamos todas las funciones nulas. Para cada $a\in A$, elegimos $k(a)$ con $a\in U_{k(a)}$ y un abierto $V_a$ que contenga $a$ y tenga cerradura compacta contenida en $U_{k(a)}$, mediante el inciso (i) del ejercicio \ref{ej:lectura-15-5} aplicado a $\{a\}$. El inciso (ii) de ese ejercicio, aplicado al compacto $\cl V_a$, produce $h_a:X\to[0,1]$ igual a $1$ en $\cl V_a$, con soporte compacto contenido en $U_{k(a)}$. Elegimos $a_1,\ldots,a_m$ de modo que los $V_{a_j}$ cubran $A$.

Agrupamos por el abierto al que están subordinadas:
\[
g_k=\sum_{j:k(a_j)=k}h_{a_j},\qquad S=\sum_{k=1}^n g_k,\qquad
f_k=\frac{g_k}{\max\{1,S\}}.
\]
Una suma sin términos se interpreta como cero. El denominador es continuo y al menos $1$. Como $0\le g_k\le S\le\max\{1,S\}$, tenemos $0\le f_k\le1$. Su soporte está contenido en la unión finita de los soportes de los $h_{a_j}$ correspondientes; esta unión es compacta y está contenida en $U_k$. Finalmente,
\[
\sum_k f_k=\frac S{\max\{1,S\}}\le1.
\]
Si $x\in A$, algún $V_{a_j}$ contiene $x$, así que $h_{a_j}(x)=1$ y $S(x)\ge1$. Allí la suma vale exactamente $1$. Normalizar por $S$ sin el máximo no serviría en los puntos donde $S=0$.\index{particion de la unidad@partición de la unidad}
\end{solucion}

\chapter{Conexidad, componentes y caminos}

\begin{ejercicio}[Lectura 17, ejercicio 1: precisión del enunciado]\label{ej:lectura-17-1}
Probar que $X$ es conexo si y solo si todo subconjunto no vacío y \emph{propio} $A\subsetneq X$ tiene frontera no vacía.
\end{ejercicio}
\begin{solucion}
La palabra propio es necesaria, pues $\partial X=\varnothing$ en cualquier espacio. Recordemos $\partial A=\cl A\setminus\interior A$. Si esta frontera es vacía, de $\interior A\subset A\subset\cl A$ se deduce $\interior A=A=\cl A$; por tanto $A$ es abierto y cerrado. Recíprocamente, si $A$ es abierto y cerrado, esas tres expresiones coinciden y su frontera es vacía. Así la condición del ejercicio excluye exactamente los abiertos y cerrados no triviales, que caracterizan la desconexión.
\end{solucion}

\begin{ejercicio}[Lectura 17, ejercicio 2]\label{ej:lectura-17-2}
Probar: (i) los subconjuntos conexos de $\R$ son exactamente los intervalos; (ii) una función real continua sobre un espacio conexo toma todos los valores intermedios entre dos valores de su imagen; (iii) si toma un valor negativo y uno positivo, tiene un cero.
\end{ejercicio}
\begin{solucion}
Para (i), los intervalos son conexos por el lema \ref{ch17:intervalos}. Si $A\subset\R$ no es un intervalo, hay $a,b\in A$ y $c\notin A$ con $a<c<b$. Entonces $A\cap(-\infty,c)$ y $A\cap(c,\infty)$ son dos abiertos relativos disjuntos no vacíos que cubren $A$, de modo que $A$ no es conexo.

Para (ii), si $X$ es conexo y $f:X\to\R$ continua, el teorema \ref{ch17:imagen} muestra que $f(X)$ es conexo, y por (i) es un intervalo. Si $f(x)=a\le b=f(y)$, todo $t\in[a,b]$ pertenece a $f(X)$, lo que significa que existe $z\in X$ con $f(z)=t$. Para (iii), aplicamos (ii) a $t=0$, que está entre el valor negativo y el positivo. Los extremos de $X$ no necesitan existir: el argumento solo exige dos puntos donde se conozcan los signos.\index{teorema del valor intermedio}\index{teorema de Bolzano}
\end{solucion}

\begin{ejercicio}[Lectura 17, ejercicio 3]\label{ej:lectura-17-3}
Si $A$ es conexo y $A\subset B\subset\cl A$, probar que $B$ es conexo. Deducir que la cerradura de un conexo es conexa.
\end{ejercicio}
\begin{solucion}
Sea $f:B\to\{0,1\}$ continua. Su restricción a $A$ es constante, digamos de valor $0$. Si $b\in B$ tuviera $f(b)=1$, la preimagen $f^{-1}(\{1\})$ sería una vecindad abierta relativa de $b$ que no toca $A$. Pero $b\in\cl A$ y $A\subset B$, de modo que toda vecindad relativa de $b$ en $B$ toca $A$, contradicción. Así $f$ es constante y $B$ es conexo. Para $A=\varnothing$, también $B=\varnothing$ y el resultado es válido. Tomar $B=\cl A$ da la conclusión final.
\end{solucion}

\begin{ejercicio}[Lectura 17, ejercicio 4]\label{ej:lectura-17-4}
Probar que todo espacio normado es conexo.
\end{ejercicio}
\begin{solucion}
Para $x\in V$, el segmento $[0,x]=\{tx:0\le t\le1\}$ es imagen continua de $[0,1]$: $\norm{tx-sx}=|t-s|\norm{x}$. Es por tanto conexo. Todos estos segmentos contienen $0$, y su unión es $V$. La proposición \ref{ch17:union} prueba que $V$ es conexo. No se requiere completitud; el argumento vale también en un espacio normado complejo usando el parámetro real $t$.
\end{solucion}

\begin{ejercicio}[Lectura 17, ejercicio 5]\label{ej:lectura-17-5}
Sea $X$ conexo, $X=A\cup B$ y $A\cap B$ conexo. Probar que $A,B$ son conexos si (i) ambos son cerrados o (ii) ambos son abiertos. En (iii), decidir qué ocurre si $A$ es abierto y $B$ cerrado.
\end{ejercicio}
\begin{solucion}
Si $A\cap B=\varnothing$ y ambos son abiertos, forman una separación salvo que uno sea vacío. Si ambos son cerrados, cada uno es complemento del otro y por tanto también abierto; la misma conclusión se aplica. El miembro no vacío es entonces $X$ y ambos son conexos.

Supongamos ahora $A\cap B\ne\varnothing$. Si $A$ fuera disconexo, existiría una función continua no constante $g:A\to\{0,1\}$. Su restricción al conexo $A\cap B$ sería constante, de valor $c$. Definimos $f=g$ en $A$ y $f=c$ en $B$; las definiciones coinciden en la intersección. Si $A,B$ son abiertos, las preimágenes de abiertos por $f$ son uniones de conjuntos abiertos en $X$, así que $f$ es continua. Si $A,B$ son cerrados, para cada cerrado $F$ del codominio la preimagen es la unión de un cerrado relativo de $A$ y uno de $B$; ambos son cerrados en $X$, luego la unión lo es y $f$ también es continua. En cualquiera de los casos $f$ sería continua no constante sobre $X$, contradicción. Así $A$ es conexo. Intercambiando los papeles de $A$ y $B$, el mismo argumento de pegado prueba que $B$ es conexo.

Para (iii) no se garantiza la conexidad de ninguno de los dos por separado. En $X=\R$, tomemos $A=\R\setminus\{0\}$ y $B=[-1,0]$. El primero es abierto y disconexo, el segundo cerrado, su unión es $\R$ y su intersección $[-1,0)$ es conexa. Para que falle la conexidad de $B$, tomemos $X=[0,2]$, $A=[0,2)$, abierto relativo, y $B=\{1,2\}$, cerrado relativo. Cubren $X$ y su intersección es $\{1\}$, pero $B$ es disconexo. La condición uniforme abierto/abierto o cerrado/cerrado es la que permite el pegado anterior.
\end{solucion}

\begin{ejercicio}[Lectura 17, ejercicio 6]\label{ej:lectura-17-6}\label{ch17:ex6}
Sean $X$ conexo, $A\subset X$ conexo y $B\subset X\setminus A$ abierto y cerrado en $X\setminus A$. Probar que $A\cup B$ es conexo.
\end{ejercicio}
\begin{solucion}
Si $A$ es vacío, $B$ es abierto y cerrado en $X$ y por tanto es vacío o todo $X$, ambos conexos. Si $A\ne\varnothing$, sea $f:A\cup B\to\{0,1\}$ continua. En $A$ es constante; cambiando los nombres de $0$ y $1$ podemos suponer que vale $0$. Pongamos $D=f^{-1}(\{1\})\subset B$. Este conjunto es abierto y cerrado en $A\cup B$.

Demostraremos que $D$ es abierto y cerrado en $X$. Como $B$ es cerrado en $X\setminus A$, tenemos $\cl B\subset A\cup B$. Por ello $\cl D\subset A\cup B$ y, puesto que $D$ es cerrado en este subespacio, $\cl D=D$: es cerrado en $X$. Para la apertura, si $x\in D$, existe un abierto $O$ de $X$ con $x\in O$ y $O\cap(A\cup B)\subset D$. Como $B$ es abierto en $X\setminus A$, existe un abierto $P$ de $X$ con $x\in P$ y $P\cap(X\setminus A)\subset B$. Todo punto de $O\cap P$ que estuviera en $A$ o en $B$ estaría en $D$ por la primera condición; todo punto que no estuviera en $A$ estaría en $B$ por la segunda. Así $O\cap P\subset D$, y $D$ es abierto. La conexidad de $X$ fuerza $D=\varnothing$ o $D=X$; la segunda posibilidad es imposible porque $A\ne\varnothing$ y $D\cap A=\varnothing$. Entonces $f$ es constante y $A\cup B$ es conexo.

La sugerencia de la lectura es una versión de pegado que también podemos justificar. Si $f$ es continua por separado en $E$ y $F$, y $E\setminus F$ es abierto y cerrado en la diferencia simétrica $E\mathbin{\triangle}F$, entonces $f$ es continua en $E\cup F$, siempre que las definiciones coincidan en la intersección. Un punto de $E\setminus F$ tiene una vecindad relativa en la unión contenida en $E$, por la apertura en la diferencia simétrica; la cerradura da la afirmación correspondiente para $F\setminus E$. En un punto de $E\cap F$, intersectamos las dos vecindades proporcionadas por la continuidad de las restricciones. En los tres casos obtenemos la continuidad en la unión.
\end{solucion}

\begin{ejercicio}[Lectura 17, ejercicio 7]\label{ej:lectura-17-7}
Construir: (i) un espacio métrico conexo con una bola disconexa; (ii) un conjunto conexo cuyo interior sea disconexo.
\end{ejercicio}
\begin{solucion}
Para (i), en $\R^2$ con la distancia euclídea consideremos
\[
X=(\{-1,1\}\times[0,2])\cup([-1,1]\times\{2\}).
\]
Es la unión de tres segmentos conexos concatenados, luego es conexo. La bola relativa de centro $(-1,0)$ y radio $\sqrt5$ contiene el brazo izquierdo completo, el tramo superior con $-1\le x<0$, y el tramo derecho $\{1\}\times[0,1)$. No contiene ningún punto con primera coordenada entre $0$ y $1$. Sus intersecciones con los semiplanos $x<1/2$ y $x>1/2$ son abiertos relativos disjuntos no vacíos que la cubren. Por tanto la bola es disconexa.

Para (ii), tomemos la unión $K$ de los discos cerrados de radio $1$ centrados en $(-1,0)$ y $(1,0)$. Cada disco es conexo por convexidad, y ambos se tocan en $(0,0)$; su unión es conexa. Su interior es exactamente la unión de los dos discos abiertos, que son disjuntos. Los puntos de frontera distintos del punto de tangencia no son interiores porque tienen vecindades pequeñas que no tocan el otro disco y contienen puntos exteriores al primero. El punto de tangencia tampoco es interior, pues $(0,t)\notin K$ para todo $t\ne0$. Así el interior tiene dos partes abiertas no vacías separadas por el eje vertical y es disconexo.
\end{solucion}

\begin{ejercicio}[Lectura 18, ejercicio 1]\label{ej:lectura-18-1}
Si un espacio métrico tiene una cantidad finita de componentes, probar que cada componente es abierta.
\end{ejercicio}
\begin{solucion}
Sean $C_1,\ldots,C_m$ las componentes. Cada una es cerrada, y para un índice $i$,
\[
X\setminus C_i=\bigcup_{j\ne i}C_j
\]
es una unión finita de cerrados, por tanto cerrada. Su complemento $C_i$ es abierto. La finitud es esencial para este argumento: en $\Q$ hay infinitas componentes unipuntuales y ninguna es abierta.
\end{solucion}

\begin{ejercicio}[Lectura 18, ejercicio 2]\label{ej:lectura-18-2}
Probar que todo espacio ultramétrico es totalmente disconexo.
\end{ejercicio}
\begin{solucion}
Recordemos que una ultramétrica satisface $d(x,z)\le\max\{d(x,y),d(y,z)\}$. Toda bola abierta $B(a,r)$ es también cerrada. En efecto, si $y\notin B(a,r)$, entonces $d(a,y)\ge r$. Para $z\in B(y,d(a,y))$, si $d(a,z)<d(a,y)$, la desigualdad ultramétrica daría
\[
d(a,y)\le\max\{d(a,z),d(z,y)\}<d(a,y),
\]
una contradicción. Así $d(a,z)\ge d(a,y)\ge r$, y la bola alrededor de $y$ queda en el complemento de $B(a,r)$. El complemento es abierto.

Dados $x\ne y$, la bola $B(x,d(x,y))$ contiene a $x$ y no a $y$, y es abierta y cerrada. Sus intersecciones con cualquier conjunto que contenga ambos puntos lo separan. Por tanto ningún conexo puede contener dos puntos distintos, que es la desconexión total.
\end{solucion}

\begin{ejercicio}[Lectura 18, ejercicio 3]\label{ej:lectura-18-3}
Un espacio es totalmente separado si para cualesquiera puntos \emph{distintos} $x,y$ hay una separación $U,V$ del espacio con $x\in U$, $y\in V$. Probar que Cantor es totalmente separado.
\end{ejercicio}
\begin{solucion}
Usamos la construcción por eliminación de tercios del ejercicio \ref{ej:lectura-8-6}. En la etapa $n$, el conjunto $C_n$ es unión de $2^n$ intervalos cerrados, disjuntos entre sí, de longitud $3^{-n}$; dos intervalos distintos están separados por un intervalo abierto de longitud positiva. Además $C\subset C_n$ y $C=\bigcap_n C_n$.

Dados $x,y\in C$ distintos, elegimos $n$ con $3^{-n}<|x-y|$. No pueden estar en el mismo intervalo de $C_n$, porque el diámetro de cada uno es $3^{-n}$. Sea $I$ el intervalo de esa etapa que contiene a $x$ y sea $J$ la unión de los demás. Definimos
\[
U=C\cap I,\qquad V=C\cap J.
\]
Ambos son cerrados relativos en $C$: $I$ es cerrado y $J$ es una unión finita de cerrados. Son disjuntos, su unión es $C$ y contienen respectivamente a $x$ y a $y$. Al ser complementarios, ambos son también abiertos relativos. Así constituyen la separación requerida. La precisión ``distintos'' es necesaria: un mismo punto no puede pertenecer a dos abiertos disjuntos.\index{separacion total@separación total}

Esta prueba también tendrá una interpretación por coordenadas cuando establezcamos el homeomorfismo $C\cong\{0,1\}^{\N}$ en la proposición \ref{ch16:cantor}. Dos puntos distintos corresponden entonces a sucesiones que difieren en alguna coordenada $k$. Los conjuntos $\{b:b_k=0\}$ y $\{b:b_k=1\}$ son abiertos en la topología producto, pues la proyección de coordenada es continua y los singletons de $\{0,1\}$ son abiertos. Son complementarios y no vacíos; transportados a $C$, separan los dos puntos. Si los dígitos aparecen en el orden contrario, se intercambian los conjuntos. Esta segunda descripción no se necesita para la prueba mediante intervalos.
\end{solucion}

\begin{ejercicio}[Lectura 18, ejercicio 4]\label{ej:lectura-18-4}
Si $A\subset X$ es abierto y cerrado, probar que es unión de componentes de $X$.
\end{ejercicio}
\begin{solucion}
Si una componente $C$ toca $A$, los conjuntos $C\cap A$ y $C\cap(X\setminus A)$ son abiertos y cerrados relativos, disjuntos, y cubren $C$. La conexidad de $C$ y el hecho de que $C\cap A$ sea no vacío obligan a que el otro sea vacío. Por tanto $C\subset A$. Cada punto de $A$ pertenece a una componente, y acabamos de probar que esa componente entera está en $A$. Así $A$ es la unión de las componentes que toca. No se afirma que toda unión de componentes sea abierta: en $\Q$ un singleton es una unión de componentes pero no es abierto.
\end{solucion}

\begin{ejercicio}[Lectura 18, ejercicio 5]\label{ej:lectura-18-5}
Probar que $A=\{(x,x\sin(1/x)):x\ne0\}\cup\{(0,0)\}$ es arcoconexo.
\end{ejercicio}
\begin{solucion}
Definimos $h:\R\to\R$ por $h(x)=x\sin(1/x)$ para $x\ne0$ y $h(0)=0$. Fuera del cero es continua, y en cero la estimación $|h(x)|\le|x|$ prueba continuidad. El mapa $F(x)=(x,h(x))$ es continuo e inyectivo, y su inversa sobre $A$ es la primera proyección, que es continua. Por tanto $A$ es homeomorfo a $\R$. Explícitamente, para dos puntos distintos $F(a),F(b)$, el mapa
\[
\gamma(t)=F((1-t)a+tb),\qquad0\le t\le1,
\]
es continuo e inyectivo porque $a\ne b$ y la primera coordenada es estrictamente monótona. Tiene los extremos requeridos y su imagen está en $A$. El factor $x$ amortigua las oscilaciones y permite atravesar el origen, a diferencia del ejemplo \ref{ch18:seno}.
\end{solucion}

\begin{ejercicio}[Lectura 18, ejercicio 6]\label{ej:lectura-18-6}
Sean $X$ conexo, $A\subset X$ conexo y $B$ una componente de $X\setminus A$. Probar que $X\setminus B$ es conexo.
\end{ejercicio}
\begin{solucion}
Si $A=\varnothing$, la única componente del espacio no vacío conexo $X$ es $X$, así que $X\setminus B=\varnothing$ es conexo. Supongamos $A\ne\varnothing$ y que $Y=X\setminus B$ admite una separación $U,V$. Puesto que $A\subset Y$ es conexo, debe estar enteramente en uno de sus miembros; renombramos para que $A\subset U$. El conjunto $V$ es no vacío, abierto y cerrado en $Y=X\setminus B$.

Aplicamos el ejercicio \ref{ej:lectura-17-6} al espacio conexo $X$, al conjunto conexo $B$ y al conjunto $V$, que es abierto y cerrado en su complemento. Obtenemos que $B\cup V$ es conexo. Como $V\cap A=\varnothing$ y $B\subset X\setminus A$, esta unión es un subconjunto conexo de $X\setminus A$. Contiene estrictamente a $B$, porque $V$ es no vacío y disjunto de $B$. Esto contradice la maximalidad de $B$ como componente de $X\setminus A$. Por tanto $Y$ no admite separación.
\end{solucion}

\begin{ejercicio}[Lectura 18, ejercicio 7: convención necesaria]\label{ej:lectura-18-7}
Definamos $x\approx y$ cuando $x=y$ o existe un arco de extremos $x,y$. Probar que es una relación de equivalencia y que sus clases, las componentes arcoconexas, están contenidas en componentes conexas.
\end{ejercicio}
\begin{solucion}
Sin la alternativa $x=y$, la relación escrita con arcos inyectivos no sería reflexiva: una aplicación inyectiva de $[0,1]$ no puede tener los mismos valores en $0$ y $1$. Con la convención corregida, la reflexividad forma parte de la definición. Para simetría, si $p$ es un arco de $x$ a $y$, $t\mapsto p(1-t)$ es un arco de $y$ a $x$. Para transitividad, los casos con puntos iguales son inmediatos por sustitución; si hay arcos de $x$ a $y$ y de $y$ a $z$, el lema \ref{ch18:recorte} produce un arco de $x$ a $z$ cuando son distintos, y si coinciden se aplica la primera alternativa. Así la relación es de equivalencia.

Además, si $p$ une $x$ y $y$, todo punto de su imagen está en la clase de $x$, porque la restricción de $p$ al intervalo que va de $0$ al parámetro de ese punto, reparametrizada linealmente, es un arco, salvo el punto inicial que ya está incluido por igualdad. En consecuencia, las clases son arcoconexas como subespacios, no solo pares de puntos conectables mediante arcos que podrían salir de la clase. Por la proposición \ref{ch18:arcoimplica}, cada clase es conexa; la proposición \ref{ch18:componentes} la sitúa dentro de una componente conexa. También se obtiene directamente de que la imagen de cualquier arco es un conjunto conexo que contiene sus dos extremos.\index{componente arcoconexa}
\end{solucion}

\chapter{Compactificaciones y espacios universales}

\begin{ejercicio}[Lectura 16, ejercicio 1]\label{ej:lectura-16-1}
Probar que $x\mapsto e^{2i\arctan x}$ encaja $\R$ densamente en $S^1$. Probar que $S^1$ y $[0,1]$ no son homeomorfos.
\end{ejercicio}
\begin{solucion}
La función $2\arctan$ es un homeomorfismo de $\R$ sobre $(-\pi,\pi)$, y $t\mapsto e^{it}$ es un homeomorfismo de ese intervalo sobre $S^1\setminus\{-1\}$. Su inversa es el argumento en $(-\pi,\pi)$; es continuo en la circunferencia cortada porque localmente puede elegirse un intervalo angular de longitud menor que $2\pi$. Así la composición es un encaje. Su imagen es densa: el único punto que falta, $-1$, es límite de $e^{2i\arctan n}$.

Al quitar cualquier punto de $S^1$ queda un espacio homeomorfo a $\R$, por una rotación seguida de la parametrización anterior; por tanto queda conexo. Al quitar $1/2$ de $[0,1]$, quedan los dos abiertos relativos no vacíos $[0,1/2)$ y $(1/2,1]$, que son disjuntos. Un homeomorfismo induciría otro entre los espacios con los puntos correspondientes eliminados, y preservaría conexidad. Esto es imposible. La conexidad de los intervalos quedó demostrada en el lema \ref{ch17:intervalos} a partir de la propiedad del supremo.
\end{solucion}

\begin{ejercicio}[Lectura 16, ejercicio 2]\label{ej:lectura-16-2}
Probar que $S^2$, el toro $S^1\times S^1$ y $[0,1]^2$ son compactificaciones de $\R^2$ y son dos a dos no homeomorfos.
\end{ejercicio}
\begin{solucion}
\emph{Construcción de las compactificaciones.} Los tres espacios son compactos: la esfera es cerrada y acotada en $\R^3$, y los otros dos son productos finitos de compactos. La proyección estereográfica inversa
\[
\sigma(u,v)=\left(\frac{2u}{1+u^2+v^2},
\frac{2v}{1+u^2+v^2},\frac{u^2+v^2-1}{1+u^2+v^2}\right)
\]
es un homeomorfismo de $\R^2$ sobre $S^2\setminus\{(0,0,1)\}$. Sustituir la fórmula verifica que la suma de cuadrados es uno, y la inversa es $(x,y,z)\mapsto(x/(1-z),y/(1-z))$. Ambos mapas son continuos en sus dominios. El punto norte es límite de $\sigma(n,0)$, de modo que la imagen es densa. El producto de dos copias del encaje del ejercicio \ref{ej:lectura-16-1} encaja $\R^2$ densamente en el toro. El producto del mapa $u:\R\to(0,1)$ del ejemplo sobre la recta extendida encaja $\R^2$ densamente en $[0,1]^2$. La densidad en los productos se verifica en rectángulos abiertos, que forman una base.

El producto $S^1\times S^1$ está naturalmente en $\R^4$, no literalmente en $\R^3$. Su modelo usual en $\R^3$ es la imagen de
\[
(e^{i\alpha},e^{i\beta})\longmapsto
((2+\cos\beta)\cos\alpha,(2+\cos\beta)\sin\alpha,\sin\beta).
\]
Es continuo e inyectivo: el radio cilíndrico $\rho=\sqrt{x^2+y^2}$ determina $\cos\beta=\rho-2$, la tercera coordenada determina $\sin\beta$, y $(x/\rho,y/\rho)$ determina $e^{i\alpha}$. Como el dominio es compacto, es un homeomorfismo sobre su imagen.

\emph{Herramienta elemental para distinguirlos.} Un \emph{lazo} en un espacio es una aplicación continua $\gamma:[0,1]\to X$ con $\gamma(0)=\gamma(1)$. Decimos que se contrae si existe un mapa continuo $F:[0,1]^2\to X$ con $F(t,0)=\gamma(t)$, $F(t,1)$ constante en $t$ y $F(0,s)=F(1,s)$ para todo $s$. Esta propiedad es preservada por homeomorfismos. Demostraremos los hechos que necesitamos, sin suponer teoría del grupo fundamental. Usaremos la conexidad del intervalo, demostrada mediante el supremo en el lema \ref{ch17:intervalos}, y las nociones de caminos del capítulo \ref{cap:nuevo06}.\index{lazo}\index{homotopia@homotopía}

Todo camino $\gamma:[0,1]\to S^1$ admite un argumento continuo $\vartheta:[0,1]\to\R$ con $\gamma(t)=e^{i\vartheta(t)}$. Para construirlo, por continuidad uniforme elegimos una partición tal que, en cada subintervalo $[t_j,t_{j+1}]$, $|\gamma(t)-\gamma(t_j)|<1$. El cociente $\gamma(t)/\gamma(t_j)$ evita $-1$ y tiene un argumento principal continuo en $(-\pi,\pi)$. Elegida $\vartheta(0)$, definimos en ese subintervalo
\[
\vartheta(t)=\vartheta(t_j)+\operatorname{Arg}(\gamma(t)/\gamma(t_j)).
\]
Las definiciones se pegan en los extremos. Dos argumentos continuos difieren en un múltiplo constante de $2\pi$, pues su diferencia es continua y toma valores en el conjunto discreto $2\pi\Z$.

Para un lazo, el entero
\[
\operatorname{ind}(\gamma)=\frac{\vartheta(1)-\vartheta(0)}{2\pi}
\]
no depende del argumento elegido. Si dos lazos $\alpha,\beta$ satisfacen $\sup_t|\alpha(t)-\beta(t)|<2$, entonces $q(t)=\operatorname{Arg}(\beta(t)/\alpha(t))$ es continuo y $q(0)=q(1)$. Un argumento de $\beta$ se obtiene sumando $q$ a uno de $\alpha$, así que tienen el mismo índice. En una deformación continua $F(t,s)$, la continuidad uniforme permite dividir el intervalo de $s$ en finitos pasos tales que lazos sucesivos estén uniformemente a distancia menor que $1$. Por lo anterior el índice es constante durante la deformación. El lazo $e^{2\pi it}$ tiene índice $1$; un lazo constante tiene índice $0$. En consecuencia, el primero no se contrae en $S^1$.\index{indice de un lazo@índice de un lazo}

\emph{El toro tiene un lazo que no se contrae.} Consideremos $\gamma(t)=(e^{2\pi it},1)$. Si se contrajera en el toro, al componer la contracción con la primera proyección obtendríamos una contracción de $e^{2\pi it}$ en $S^1$, que acabamos de excluir.

\emph{Todos los lazos de la esfera se contraen.} Este paso necesita justificación: una curva continua puede llenar una región, así que no debemos suponer sin prueba que omite algún punto de la esfera. Sea $\gamma$ un lazo en $S^2$. Elegimos una partición tan fina que en cada intervalo $[t_j,t_{j+1}]$ se cumpla $\norm{\gamma(t)-\gamma(t_j)}<1/4$. Unimos los valores consecutivos $a_j=\gamma(t_j)$ y $a_{j+1}$ mediante el arco
\[
\beta(t)=\frac{(1-r)a_j+ra_{j+1}}{\norm{(1-r)a_j+ra_{j+1}}},\qquad
r=\frac{t-t_j}{t_{j+1}-t_j}.
\]
El denominador es positivo porque los extremos están a distancia menor que $1/4$, de modo que no son antipodales. Si $v=(1-r)a_j+ra_{j+1}$, entonces $\norm{v-a_j}<1/4$ y
\[
\norm{v/\norm v-a_j}\le |1-\norm v|+\norm{v-a_j}<1/2.
\]
Por tanto $\norm{\beta(t)-\gamma(t)}<3/4$. Los dos valores nunca son antipodales, y
\[
F(t,s)=\frac{(1-s)\gamma(t)+s\beta(t)}{\norm{(1-s)\gamma(t)+s\beta(t)}}
\]
define una deformación continua de lazos entre $\gamma$ y $\beta$.

La imagen de $\beta$ está en una unión finita de planos propios que pasan por el origen, intersectados con $S^2$, porque cada arco está en el plano generado por sus extremos; si ambos son iguales elegimos cualquier plano que contenga ese punto. Hay un punto $q\in S^2$ fuera de todos esos planos. Para verlo sin un argumento de medida, escribamos sus ecuaciones como $A_jx+B_jy+C_jz=0$, con algún coeficiente no cero. El vector $(1,t,t^2)$ pertenece al plano $j$ solo si $A_j+B_jt+C_jt^2=0$, ecuación con finitas raíces. Elegimos $t$ fuera de la unión finita de raíces y normalizamos el vector. Entonces $q$ no pertenece a la imagen de $\beta$.

Una rotación seguida de la proyección estereográfica identifica $S^2\setminus\{q\}$ con $\R^2$. Allí todo lazo se contrae por interpolación lineal con un punto fijo. Transportamos esa contracción a la esfera y la concatenamos con $F$, usando la primera mitad del parámetro para $F$ y la segunda para la contracción. Así $\gamma$ se contrae. La esfera y el toro no son homeomorfos.

\emph{El cuadrado y el toro tampoco son homeomorfos.} Todo lazo del cuadrado se contrae mediante $F(t,s)=(1-s)\gamma(t)+s(1/2,1/2)$, pues el cuadrado es convexo; el toro tiene el lazo no contraíble ya construido.

\emph{La esfera y el cuadrado no son homeomorfos.} Al quitar cualquier punto de la esfera queda un espacio homeomorfo a $\R^2$, donde todos los lazos se contraen. En cambio, al quitar $a=(1/2,1/2)$ del cuadrado, el lazo
\[
\gamma(t)=a+\tfrac14(\cos2\pi t,\sin2\pi t)
\]
no se contrae: componer una supuesta contracción con $x\mapsto(x-a)/\norm{x-a}$ daría una contracción del lazo de índice uno en $S^1$. Si existiera un homeomorfismo de la esfera al cuadrado, al eliminar la preimagen de $a$ produciría un homeomorfismo entre estos dos espacios perforados. La contradicción termina las tres comparaciones.
\end{solucion}

\begin{ejercicio}[Lectura 16, ejercicio 3: examen y corrección del enunciado]\label{ej:lectura-16-3}
El enunciado propone que el conjunto de Cantor sea una compactificación de todo espacio discreto infinito separable. Decidir si es válido con la definición de encaje y construir una compactificación correcta.
\end{ejercicio}
\begin{solucion}
Un espacio discreto separable es numerable, porque todo subconjunto denso debe contener cada singleton abierto. Sin embargo, \emph{no puede encajarse densamente en Cantor}. En efecto, si $D$ es un subespacio denso y discreto de un espacio métrico $K$, para cada $d\in D$ existe un abierto $U\subset K$ con $U\cap D=\{d\}$. Si $U$ contuviera un punto $y\ne d$, el abierto no vacío $U\setminus\{d\}$ tendría que tocar $D$ por densidad, contradicción. Así $U=\{d\}$ y $d$ es un punto aislado de $K$. Pero Cantor no tiene puntos aislados: manteniendo un prefijo binario y cambiando algún dígito posterior obtenemos puntos distintos arbitrariamente próximos.

La sugerencia de enumerar un subconjunto denso de vértices proporciona una inyección continua desde el discreto numerable, pero su inversa sobre la imagen no es continua. Esta es precisamente la dificultad de la definición original. Una versión correcta es la siguiente: si $X=\{x_1,x_2,\ldots\}$ es discreto, el mapa $x_n\mapsto1/n$ lo encaja densamente en
\[
K=\{0\}\cup\{1/n:n\ge1\}.
\]
Cada $1/n$ es aislado en $K$, de modo que la imagen tiene la topología discreta; $0$ es su único punto límite añadido. El conjunto $K$ es cerrado y acotado en $\R$, luego compacto.

La afirmación adicional de que todo espacio ultramétrico tiene a Cantor como compactificación tampoco es válida en esa generalidad: los discretos anteriores son ultramétricos, y los no separables ni siquiera admiten compactificación métrica. Sí puede afirmarse, por el teorema \ref{ch16:encaje}, que todo espacio ultramétrico \emph{separable} admite alguna compactificación métrica.
\end{solucion}

\begin{ejercicio}[Lectura 16, ejercicio 4]\label{ej:lectura-16-4}
Probar que $T:K\to H$, $T(x)_n=nx_n$, es un homeomorfismo, donde $K=\{x\in\ell^2:0\le x_n\le1/n\}$.
\end{ejercicio}
\begin{solucion}
El mapa toma valores en $H$ y es inyectivo. Para $y\in H$, la sucesión $x_n=y_n/n$ satisface $\sum_n|x_n|^2\le\sum_n1/n^2<\infty$ y pertenece a $K$, de modo que es sobreyectivo. Si $x^{(j)}\to x$ en $\ell^2$, entonces $|x_n^{(j)}-x_n|\le\norm{x^{(j)}-x}_2\to0$ para cada $n$; por tanto $nx_n^{(j)}\to nx_n$ y la convergencia coordenada implica convergencia en $H$. Así $T$ es continua. Como $K$ es compacto y $H$ es métrico, la biyección continua $T$ es un homeomorfismo. Obsérvese que no estamos afirmando continuidad de este multiplicador en todo $\ell^2$ con valores en $\ell^2$: aquí el codominio tiene la topología producto.
\end{solucion}

\begin{ejercicio}[Lectura 16, ejercicio 5]\label{ej:lectura-16-5}
Probar que todo espacio métrico separable es homeomorfo a un subespacio del toro infinito $\mathbb T^\infty=(S^1)^{\N}$.
\end{ejercicio}
\begin{solucion}
Por el teorema \ref{ch16:encaje}, basta encajar $H$ en $\mathbb T^\infty$. El mapa $h:[0,1]\to S^1$, $h(t)=e^{i\pi t}$, es continuo e inyectivo; por compacidad es un homeomorfismo sobre el semicírculo superior. Aplicamos $h$ coordenada a coordenada:
\[
\Phi((x_n))=(e^{i\pi x_n})_{n\ge1}.
\]
Este mapa es continuo e inyectivo en las topologías producto. Como $H$ es compacto y el producto numerable de circunferencias es métrico, $\Phi$ es un encaje. La composición $\Phi\circ J$ da el encaje de $X$. Se utiliza un semicírculo, no $e^{2\pi it}$ en todo $[0,1]$, pues este último identificaría los dos extremos.
\end{solucion}

\chapter{Geometría normada y dimensión finita}

\begin{ejercicio}[Ampliación: una recta como cociente]\label{ej:ampl-cocientes-1}
En $\R^2$ con la norma euclidiana sea $M=\{(t,t):t\in\R\}$. Calcule la norma de $[(a,b)]$, encuentre el representante de norma mínima y construya una isometría lineal de $\R^2/M$ sobre $\R$. Este ejercicio y su desarrollo son propios de esta edición.
\end{ejercicio}
\begin{solucion}
Para cualquier $t$,
\[
 (a-t)^2+(b-t)^2
 =2\left(t-\frac{a+b}{2}\right)^2+\frac{(a-b)^2}{2}.
\]
El mínimo se alcanza únicamente en $t=(a+b)/2$, de modo que
\[
 \norm{[(a,b)]}=\frac{|a-b|}{\sqrt2},\qquad
 (a,b)-\left(\frac{a+b}{2},\frac{a+b}{2}\right)
 =\left(\frac{a-b}{2},\frac{b-a}{2}\right).
\]
La aplicación $T([(a,b)])=(a-b)/\sqrt2$ está bien definida porque restar un vector de $M$ no cambia $a-b$. Es lineal y preserva la norma calculada. Es sobreyectiva, pues $T([(\sqrt2s,0)])=s$; su núcleo es cero porque $a-b=0$ equivale a $(a,b)\in M$. Esto demuestra la isometría, no sólo una coincidencia de dimensiones.
\end{solucion}

\begin{ejercicio}[Ampliación: el límite como coordenada del cociente]\label{ej:ampl-cocientes-2}
Sean $c$ y $c_0$ los espacios de sucesiones escalares convergentes y convergentes a cero, con la norma uniforme. Demuestre que $c_0$ es cerrado en $c$ y que $[x]\mapsto\lim_nx_n$ es un isomorfismo isométrico de $c/c_0$ sobre $\F$. Adaptación del ejercicio de la presentación \cite{max-seminormado}; la solución se desarrolla aquí.
\end{ejercicio}
\begin{solucion}
Si $x^{(k)}\in c_0$ y $\norm{x^{(k)}-x}_\infty\to0$, dado $\eps>0$ elegimos $k$ con error menor que $\eps/2$ y luego $N$ con $|x^{(k)}_n|<\eps/2$ para $n\ge N$. Entonces $|x_n|<\eps$ para esos índices, así que $x\in c_0$. Esto demuestra cerradura.

Escribamos $L=\lim_nx_n$. Para todo $y\in c_0$, la sucesión $x_n-y_n$ converge a $L$, por lo que $|L|\le\norm{x-y}_\infty$. Tomar ínfimos da $|L|\le\norm{[x]}$. Por otro lado, $y=x-L\mathbf1$ pertenece a $c_0$, y $x-y=L\mathbf1$ tiene norma $|L|$. Así $\norm{[x]}=|L|$. Dos representantes difieren en $c_0$ exactamente cuando tienen el mismo límite; la aplicación indicada está bien definida, es lineal e inyectiva. Las sucesiones constantes producen todos los escalares, luego también es sobreyectiva. En este caso el representante constante alcanza el ínfimo.
\end{solucion}

\begin{ejercicio}[Ampliación: un cociente que pierde toda distancia]\label{ej:ampl-cocientes-3}
En $\ell^2$ sea $c_{00}$ el subespacio de sucesiones de soporte finito. Pruebe que $\ell^2/c_{00}$ no es el espacio vectorial cero, aunque su seminorma cociente vale cero en todas las clases. Explique por qué esto no contradice el teorema~\ref{ampl:cociente-banach}. Adaptación del contraejemplo propuesto en \cite{max-seminormado}.
\end{ejercicio}
\begin{solucion}
La sucesión $x=(1/n)_{n\ge1}$ pertenece a $\ell^2$ y no a $c_{00}$, así que $[x]\ne[0]$. Para cualquier $u\in\ell^2$, sus truncamientos $u^{(N)}\in c_{00}$ satisfacen
\[
 \norm{u-u^{(N)}}_2^2=\sum_{n>N}|u_n|^2\longrightarrow0.
\]
De aquí $\inf_{v\in c_{00}}\norm{u-v}_2=0$, y la seminorma de toda clase es cero. El subespacio $c_{00}$ es denso y propio, por lo que no es cerrado. Falta exactamente la hipótesis que haría de la seminorma una norma. Si se divide por $\cl{c_{00}}=\ell^2$, el cociente sí es el espacio cero: identificar todas las distancias nulas elimina las clases algebraicamente distintas que la seminorma era incapaz de separar.
\end{solucion}

\begin{ejercicio}[Ampliación: observar sólo los extremos]\label{ej:ampl-cocientes-4}
En $C([0,1],\F)$ defina $p(f)=|f(0)|+|f(1)|$. Describa su núcleo y pruebe que el cociente por ese núcleo es isométrico a $\F^2$ con la norma $\norm{(a,b)}_1=|a|+|b|$. Determine si el cociente es completo aunque $p$ no sea una norma. Ejercicio propio de esta edición.
\end{ejercicio}
\begin{solucion}
La desigualdad triangular y la homogeneidad de $p$ se verifican en cada extremo. Su núcleo es
\[
 N_p=\{f\in C([0,1]):f(0)=f(1)=0\}.
\]
Contiene funciones no nulas, por ejemplo $f(t)=t(1-t)$; por ello $p$ no es una norma. La aplicación $T([f])=(f(0),f(1))$ está bien definida, es lineal y tiene núcleo cero en el cociente. Para cualquier $(a,b)$, la función afín $f(t)=(1-t)a+tb$ tiene esos extremos, de modo que $T$ es sobreyectiva. La proposición~\ref{ampl:nucleo-seminorma} da
\[
 \norm{[f]}=p(f)=\norm{T([f])}_1.
\]
Finalmente, una sucesión de Cauchy en $\F^2$ para $\norm{\cdot}_1$ es de Cauchy en cada coordenada; los dos límites escalares producen el límite en esa norma. La isometría transporta la completitud al cociente. Los valores interiores de las funciones no son necesarios para recuperar el límite de sus clases.
\end{solucion}

\begin{ejercicio}[Ampliación: norma cociente sin mínimo]\label{ej:ampl-cocientes-5}
Sea $\varphi:V\to\F$ lineal y no nula, y suponga
$C=\sup_{\norm v\le1}|\varphi(v)|<\infty$. Pruebe que $M=\ker\varphi$ es cerrado y
\[
 \norm{[x]}_{V/M}=\frac{|\varphi(x)|}{C}.
\]
Después considere $V=\ell^1$ y $\varphi(x)=\sum_{n\ge1}\frac{n}{n+1}x_n$. Demuestre que $C=1$ y que la clase cuyos representantes satisfacen $\varphi(x)=1$ tiene norma $1$, pero no contiene ningún vector de norma $1$. Ejercicio propio de esta edición.
\end{ejercicio}
\begin{solucion}
La homogeneidad da $|\varphi(v)|\le C\norm v$ para todo $v$, y $C>0$ porque $\varphi\ne0$. La misma desigualdad aplicada a diferencias demuestra continuidad, luego su núcleo es cerrado. Si $m\in M$,
$|\varphi(x)|=|\varphi(x-m)|\le C\norm{x-m}$; el ínfimo proporciona la cota inferior.

Si $\varphi(x)=0$, la clase es cero y la fórmula está probada. En otro caso elegimos $u_k$ con $\norm{u_k}\le1$ y $|\varphi(u_k)|\to C$, suprimiendo los primeros términos si es necesario para que $\varphi(u_k)\ne0$. Los vectores
\[
 v_k=\frac{\varphi(x)}{\varphi(u_k)}u_k
\]
satisfacen $\varphi(v_k)=\varphi(x)$, así que representan $[x]$. Sus normas son a lo sumo $|\varphi(x)|/|\varphi(u_k)|$. Al pasar al límite resulta la cota superior y, por tanto, la igualdad.

En $\ell^1$ la serie que define $\varphi$ converge absolutamente y $|\varphi(x)|\le\norm x_1$. Como $\varphi(e_n)=n/(n+1)\to1$, el supremo $C$ es $1$. Para cada $x\ne0$,
\[
 |\varphi(x)|\le\sum_n\frac{n}{n+1}|x_n|
 =\norm x_1-\sum_n\frac{|x_n|}{n+1}<\norm x_1.
\]
La última suma es estrictamente positiva porque algún $x_n$ no es cero. Todo representante con $\varphi(x)=1$ tiene, pues, norma mayor que $1$. Sin embargo, $v_n=((n+1)/n)e_n$ satisface $\varphi(v_n)=1$ y $\norm{v_n}_1\downarrow1$. La clase existe, su ínfimo es $1$ y no lo alcanza ningún representante.
\end{solucion}

\begin{ejercicio}[Lectura 19, ejercicio 1]\label{ej:lectura-19-1}
Para $u,v\ne0$, ponga $u'=u/\norm u$ y $v'=v/\norm v$. Demuestre que
\[\norm{u'-v'}\leq\frac{2\norm{u-v}}{\norm u}.\]
\end{ejercicio}
\begin{solucion}
Introducimos $v/\norm u$, de modo que una diferencia mida el cambio de vector y la otra el cambio de escala:
\begin{align*}
 \norm{u'-v'}
 &\leq\frac{\norm{u-v}}{\norm u}
 +\norm v\left|\frac1{\norm u}-\frac1{\norm v}\right|\\
 &=\frac{\norm{u-v}+|\norm v-\norm u|}{\norm u}
 \leq\frac{2\norm{u-v}}{\norm u}.
\end{align*}
La última desigualdad es la desigualdad triangular inversa.
\end{solucion}

\begin{ejercicio}[Lectura 19, ejercicio 2]\label{ej:lectura-19-2}
Para $A,B\subseteq V$, sea $A+B=\{a+b:a\in A,b\in B\}$. Pruebe: (i) si $A$ es abierto, $A+B$ es abierto; (ii) si $A,B$ son compactos, $A+B$ es compacto; (iii) si $A$ es cerrado y $B$ compacto, $A+B$ es cerrado. Dé además dos conjuntos cerrados cuya suma no sea cerrada.
\end{ejercicio}
\begin{solucion}
(i) Cada traslación $x\mapsto x+b$ es un homeomorfismo, pues ella y su inversa preservan distancias. Así $A+B=\bigcup_{b\in B}(A+b)$ es una unión de abiertos.

(ii) El producto finito $A\times B$ es compacto y la suma es continua por la proposición~\ref{ch19:operaciones}. Su imagen es exactamente $A+B$, luego es compacta.

(iii) Si $a_n+b_n\to x$, la compacidad de $B$ da una subsucesión $b_{n_j}\to b\in B$. Entonces $a_{n_j}=(a_{n_j}+b_{n_j})-b_{n_j}\to x-b$. La cerradura de $A$ implica $x-b\in A$, y así $x\in A+B$. En un espacio métrico este criterio de sucesiones equivale a que el conjunto sea cerrado. Si alguno de los conjuntos es vacío, todas las conclusiones son inmediatas.

Para el contraejemplo tome, en $\R^2$,
\[
 A=\{(x,y):xy=1\},\qquad B=\R\times\{0\}.
\]
Ambos son cerrados, como conjuntos de ceros de funciones continuas. Su suma es $\R\times(\R\setminus\{0\})$: para obtener $(s,t)$ con $t\ne0$ se suma $(1/t,t)$ a $(s-1/t,0)$, y ningún sumando puede producir segunda coordenada cero. Esta suma no es cerrada, pues $(0,1/n)\to(0,0)$.
\end{solucion}

\begin{ejercicio}[Lectura 19, ejercicio 3]\label{ej:lectura-19-3}
Demuestre que todo subespacio propio $M$ de $V$ tiene interior vacío.
\end{ejercicio}
\begin{solucion}
Si $B(x,r)\subseteq M$, entonces $x\in M$ y, restándolo, $B(0,r)\subseteq M$. Para cualquier $v\ne0$, el vector $(r/(2\norm v))v$ pertenece a esa bola; multiplicarlo por $2\norm v/r$ muestra que $v\in M$. También $0\in M$. Esto da $M=V$, contradicción. No se necesita que $M$ sea cerrado.
\end{solucion}

\begin{ejercicio}[Lectura 19, ejercicio 4]\label{ej:lectura-19-4}
Pruebe que $V$ es separable si y solo si $\overline B(0,1)$ es separable.
\end{ejercicio}
\begin{solucion}
Si $V$ es separable, todo subespacio métrico es separable. Para recordar el argumento, tome una base numerable de bolas de $V$ y elija un punto de $\overline B(0,1)$ en cada miembro de la base que lo interseque. Los puntos elegidos forman un conjunto numerable y denso en la bola: toda vecindad relativa contiene una de esas bolas básicas alrededor de alguno de sus puntos.

Recíprocamente, si $D$ es numerable y denso en $\overline B(0,1)$, el conjunto $E=\bigcup_{m=1}^{\infty}mD$ es numerable. Dados $v\in V$ y $\eps>0$, elija un entero $m\geq\max\{1,\norm v\}$. Hay $d\in D$ con $\norm{v/m-d}<\eps/m$, y entonces $md\in E$ satisface $\norm{v-md}<\eps$.
\end{solucion}

\begin{ejercicio}[Lectura 19, ejercicio 5]\label{ej:lectura-19-5}
Si $(v_n)$ es de Cauchy, $v_n\ne0$, y no converge a $0$, pruebe que $(v_n/\norm{v_n})$ es de Cauchy. Deduzca que $V$ es de Banach si y solo si $\overline B(0,1)$ es completa.
\end{ejercicio}
\begin{solucion}
Por la desigualdad triangular inversa, $(\norm{v_n})$ es una sucesión real de Cauchy. Sea $a$ su límite. Si $a=0$, tendríamos $v_n\to0$, contrario a la hipótesis; luego $a>0$. Para $n$ suficientemente grande, $\norm{v_n}\geq a/2$. El ejercicio~\ref{ej:lectura-19-1} da entonces
\[
 \norm{\frac{v_n}{\norm{v_n}}-\frac{v_m}{\norm{v_m}}}\leq\frac4a\norm{v_n-v_m},
\]
que tiende uniformemente a cero cuando $m,n\to\infty$.

Si $V$ es completo, su bola cerrada es un subconjunto cerrado y por ello completo. Para el recíproco puede usarse una observación aún más directa: toda sucesión de Cauchy $(x_n)$ es acotada. Elija $R>0$ con $\norm{x_n}\leq R$ para todo $n$. Entonces $(x_n/R)$ es de Cauchy en la bola cerrada unitaria, converge a un $y$ de esa bola y $x_n\to Ry$ en $V$. La primera parte explica además cómo controlar las direcciones de una sucesión que permanece alejada del origen.
\end{solucion}

\begin{ejercicio}[Lectura 19, ejercicio 6]\label{ej:lectura-19-6}
Sean $M_1,M_2\leq V$ tales que $\cl{M_1}\cap M_2=\{0\}$ y $M_1+M_2$ es cerrado. Demuestre que $M_1$ es cerrado.
\end{ejercicio}
\begin{solucion}
Como $M_1\subseteq M_1+M_2$ y este último es cerrado, $\cl{M_1}\subseteq M_1+M_2$. Para $x\in\cl{M_1}$ escriba $x=a+b$, con $a\in M_1$ y $b\in M_2$. La proposición~\ref{ch19:cierre-subespacio} implica que $b=x-a\in\cl{M_1}$. Así $b\in\cl{M_1}\cap M_2=\{0\}$, y $x=a\in M_1$. Hemos probado $\cl{M_1}\subseteq M_1$. La barra sobre $M_1$ en la hipótesis, visible en el PDF original, es indispensable para este argumento.
\end{solucion}

\begin{ejercicio}[Lectura 19, ejercicio 7]\label{ej:lectura-19-7}
Si $v_n\to v$, demuestre que las medias de Cesàro $u_n=(v_1+\cdots+v_n)/n$ convergen a $v$.
\end{ejercicio}
\begin{solucion}
Dado $\eps>0$, elija $N$ con $\norm{v_j-v}<\eps/2$ para $j\geq N$. Para $n\geq N$,
\[
 \norm{u_n-v}\leq
 \frac1n\sum_{j=1}^{N-1}\norm{v_j-v}
 +\frac{n-N+1}{n}\frac\eps2.
\]
La primera suma tiene un número fijo de términos, por lo que, dividida entre $n$, es menor que $\eps/2$ para $n$ suficientemente grande. El segundo término no supera $\eps/2$. Así $u_n\to v$. El argumento separa una parte inicial, cuyo peso desaparece, de una cola cuyos términos ya están cerca del límite.
\end{solucion}

\begin{ejercicio}[Lectura 22, ejercicio 4]\label{ej:lectura-22-4}
Si $M$ tiene dimensión finita y $v\notin M$, demuestre que existe $m_0\in M$ con $\dist(v,M)=\norm{v-m_0}$. Deduzca que existe $u$ con $\norm u=\dist(u,M)=1$.
\end{ejercicio}
\begin{solucion}
Sea $d=\dist(v,M)$. El subespacio es cerrado, por lo que $d>0$. Elija $m_j\in M$ con $\norm{v-m_j}\to d$. Como $0\in M$, $d\leq\norm v$. Podemos descartar términos iniciales para que $\norm{v-m_j}\leq\norm v+1$. Entonces
\[
 \norm{m_j}\leq\norm v+\norm{v-m_j}\leq2\norm v+1.
\]
La bola correspondiente de $M$ es compacta por Heine--Borel. Una subsucesión converge a $m_0\in M$, y la continuidad de la norma da $\norm{v-m_0}=d$.

Defina $u=(v-m_0)/d$. Su norma es $1$. Como trasladar por un elemento de $M$ no cambia la distancia a $M$, y multiplicar por $1/d$ la multiplica por $1/d$,
\[
 \dist(u,M)=\frac{\dist(v-m_0,M)}d=\frac{\dist(v,M)}d=1.
\]
En dimensión finita la distancia se alcanza. El lema de Riesz general solo garantiza distancias arbitrariamente próximas a $1$, porque un subespacio cerrado infinito dimensional puede no tener mejores aproximantes.
\end{solucion}

\begin{ejercicio}[Lectura 22, ejercicio 5]\label{ej:lectura-22-5}
Sea $c_{00}=\spanop\{e_n:n\in\N\}$, el espacio de sucesiones de soporte finito. Demuestre que no es completo con ninguna norma.
\end{ejercicio}
\begin{solucion}
\enlargethispage{\baselineskip}
Fije una norma arbitraria y sea $M_n=\spanop\{e_1,\ldots,e_n\}$. Cada $M_n$ es cerrado porque tiene dimensión finita. Es propio y, por el ejercicio~\ref{ej:lectura-19-3}, tiene interior vacío. Todo vector de $c_{00}$ pertenece a alguno de los $M_n$ por tener soporte finito; así, $c_{00}=\bigcup_{n=1}^{\infty}M_n$. Si la norma elegida fuera completa, esta unión contradiría Baire. Como no se impuso ninguna condición especial a la norma, ninguna puede hacer completo a $c_{00}$.
\end{solucion}

\chapter{Operadores acotados y compactos}

\begin{ejercicio}[Lectura 20, ejercicio 1]\label{ej:lectura-20-1}
Si $T\in\Bcal(V_1,V_2)$ y $S\in\Bcal(V_2,V_3)$, demuestre que $ST\in\Bcal(V_1,V_3)$ y $\norm{ST}\leq\norm S\norm T$.
\end{ejercicio}
\begin{solucion}
La composición es lineal porque
\[
 ST(\alpha x+\beta y)=S(\alpha Tx+\beta Ty)=\alpha STx+\beta STy.
\]
Para todo $v$,
\[
 \norm{STv}\leq\norm S\norm{Tv}\leq\norm S\norm T\norm v.
\]
El teorema~\ref{ch20:acotado} da continuidad y la caracterización del ínfimo de constantes da la desigualdad de normas.
\end{solucion}

\begin{ejercicio}[Lectura 20, ejercicio 2]\label{ej:lectura-20-2}
Pruebe que un operador lineal es continuo si y solo si manda sucesiones acotadas en sucesiones acotadas.
\end{ejercicio}
\begin{solucion}
Si $T$ es continuo y $\sup_n\norm{v_n}=M<\infty$, entonces $\norm{Tv_n}\leq\norm T M$. Para el recíproco, si $T$ no fuera continuo, no existiría una constante que lo acote. Para cada $n$ podríamos elegir $x_n\ne0$ con $\norm{Tx_n}>n\norm{x_n}$. Los vectores $v_n=x_n/\norm{x_n}$ forman una sucesión acotada, pero $\norm{Tv_n}>n$, contradicción.
\end{solucion}

\begin{ejercicio}[Lectura 20, ejercicio 3]\label{ej:lectura-20-3}
Si $T:V\to W$ es lineal continuo y $F\leq W$ tiene dimensión finita, demuestre que $T^{-1}(F)$ es un subespacio cerrado. Deduzca que $\ker T$ es cerrado.
\end{ejercicio}
\begin{solucion}
Si $x,y\in T^{-1}(F)$, entonces $Tx,Ty\in F$. Para $\alpha,\beta\in\F$, $T(\alpha x+\beta y)=\alpha Tx+\beta Ty\in F$, así que la preimagen es un subespacio. Por el lema~\ref{ch19:coordenadas}, $F$ es cerrado. La preimagen de un cerrado por una función continua es cerrada. Aplicamos esto a $F=\{0\}$ para obtener la afirmación sobre el núcleo.
\end{solucion}

\begin{ejercicio}[Lectura 20, ejercicio 4]\label{ej:lectura-20-4}
Demuestre que un operador lineal es continuo si y solo si es continuo en algún punto de su dominio.
\end{ejercicio}
\begin{solucion}
Una función continua en todo el dominio lo es en cualquier punto. Si $T$ es continuo en $v_0$, dado $\eps>0$ existe $\delta>0$ tal que $\norm{x-v_0}<\delta$ implica $\norm{Tx-Tv_0}<\eps$. Ponga $x=v_0+h$. Como $Tx-Tv_0=Th$, se obtiene $\norm h<\delta\Rightarrow\norm{Th}<\eps$. Así $T$ es continuo en cero, y el teorema~\ref{ch20:acotado} termina la prueba.
\end{solucion}

\begin{ejercicio}[Lectura 20, ejercicio 5]\label{ej:lectura-20-5}
Si $T_n\to T$ en $\Bcal(V_1,V_2)$ y $S_n\to S$ en $\Bcal(V_2,V_3)$, demuestre que $S_nT_n\to ST$ en norma.
\end{ejercicio}
\begin{solucion}
La sucesión convergente $(S_n)$ es acotada; sea $M=\sup_n\norm{S_n}<\infty$. Usamos la identidad
\[
 S_nT_n-ST=S_n(T_n-T)+(S_n-S)T.
\]
La desigualdad del ejercicio~\ref{ej:lectura-20-1} da
\[
 \norm{S_nT_n-ST}\leq M\norm{T_n-T}+\norm{S_n-S}\norm T\longrightarrow0.
\]
La acotación de $S_n$ es necesaria para controlar el producto del primer término; no basta anunciar un ``álgebra de límites'' sin identificar esa cota.
\end{solucion}

\begin{ejercicio}[Lectura 20, ejercicio 6]\label{ej:lectura-20-6}
En $\R^2$ con norma uniforme, sea $T$ el operador de matriz $\begin{pmatrix}a&b\\c&d\end{pmatrix}$. Demuestre que
\[\norm T=\max\{|a|+|b|,|c|+|d|\}.\]
\end{ejercicio}
\begin{solucion}
Para $\norm{(x,y)}_\infty\leq1$,
\[
 |ax+by|\leq|a|+|b|,\qquad |cx+dy|\leq|c|+|d|.
\]
Esto prueba la cota superior. Si el máximo lo da la primera fila, elija $x=\operatorname{sgn}(a)$ e $y=\operatorname{sgn}(b)$, tomando el valor $1$ cuando el coeficiente sea cero. Entonces $\norm{(x,y)}_\infty=1$ y $ax+by=|a|+|b|$. Si el máximo lo da la segunda fila, elija de la misma manera los signos de $c,d$, para que $cx+dy=|c|+|d|$. En ambos casos la norma de la salida alcanza la cota.
\end{solucion}

\begin{ejercicio}[Lectura 20, ejercicio 7]\label{ej:lectura-20-7}
Sea $(T_n)$ una sucesión acotada en $\Bcal(V,W)$ que converge puntualmente a $T$. Si $A\subseteq V$ es totalmente acotado, pruebe que la convergencia es uniforme en $A$.
\end{ejercicio}
\begin{solucion}
Sea $M=\sup_n\norm{T_n}$. El lema~\ref{ch20:limite-puntual} da $\norm T\leq M$. Si $A$ es vacío o $M=0$, la conclusión es inmediata. Suponga $M>0$. Dado $\eps>0$, la acotación total proporciona puntos $a_1,\ldots,a_k\in V$ tales que cada $v\in A$ está a distancia menor que $\delta=\eps/(6M)$ de alguno de ellos. Para cada $i$, $T_na_i\to Ta_i$. Como hay solo un número finito de centros, existe un mismo $N$ tal que
\[
 \norm{T_na_i-Ta_i}<\eps/3\quad(n\geq N,\ 1\leq i\leq k).
\]
Para $v\in A$, elija un centro con $\norm{v-a_i}<\delta$. Entonces
\begin{align*}
 \norm{T_nv-Tv}
 &\leq\norm{T_n(v-a_i)}+\norm{T_na_i-Ta_i}+\norm{T(a_i-v)}\\
 &<M\delta+\eps/3+M\delta=2\eps/3.
\end{align*}
Por tanto $\sup_{v\in A}\norm{T_nv-Tv}\leq2\eps/3<\eps$ para $n\geq N$. La cota finita de centros convierte un umbral dependiente del punto en un umbral común.
\end{solucion}

\begin{ejercicio}[Lectura 22, ejercicio 1]\label{ej:lectura-22-1}
Demuestre que toda transformación lineal sobreyectiva de un espacio normado sobre un espacio normado de dimensión finita es abierta.
\end{ejercicio}
\begin{solucion}
Sea $T:V\to W$ sobreyectiva. Si $W=\{0\}$, la imagen de un abierto es $\varnothing$ o $\{0\}$, ambos abiertos en $W$. Supongamos $\dim W=k>0$. Elija una base $w_1,\ldots,w_k$ y preimágenes $v_j$ con $Tv_j=w_j$. Defina
\[
 R\left(\sum_j a_jw_j\right)=\sum_j a_jv_j.
\]
Es una aplicación lineal con $TR=I_W$. Como su dominio es finito dimensional, es continua; sea $C=\norm R>0$. Si $U\subseteq V$ es abierto y $y=Tx\in T(U)$ con $x\in U$, existe $r>0$ con $B(x,r)\subseteq U$. Para $z\in W$ con $\norm{z-y}<r/C$,
\[
 \norm{R(z-y)}\leq C\norm{z-y}<r.
\]
Por ello $x+R(z-y)\in U$, y su imagen es $y+(z-y)=z$. Así $B(y,r/C)\subseteq T(U)$ y $T(U)$ es abierto.

No se supuso que $T$ fuera continua. La apertura procede de la existencia de una sección lineal continua $R$ con dominio finito dimensional; no es una aplicación del teorema de la aplicación abierta para Banach.
\end{solucion}

\begin{ejercicio}[Lectura 22, ejercicio 2]\label{ej:lectura-22-2}
Sea $V$ finito dimensional con base $v_1,\ldots,v_k$, donde $k\geq1$. Pruebe que existe $C>0$ tal que
\[
 \norm T\leq C\max_j\norm{Tv_j}
\]
para todo operador lineal $T$ de $V$. Deduzca que si $T_n$ converge puntualmente a cero, entonces $\norm{T_n}\to0$. Compruebe por separado esta última conclusión cuando $V=\{0\}$.
\end{ejercicio}
\begin{solucion}
Por el lema de coordenadas existe $c>0$ tal que $\norm{\sum_j a_jv_j}\geq c\max_j|a_j|$. Para $v=\sum_j a_jv_j$,
\[
 \norm{Tv}\leq\sum_j|a_j|\norm{Tv_j}
 \leq\frac kc\norm v\max_j\norm{Tv_j}.
\]
Tomando el supremo sobre $\norm v\leq1$, sirve $C=k/c$. El argumento funciona también si el codominio de $T$ es cualquier otro espacio normado. Si $T_n$ converge puntualmente a cero, cada una de las $k$ sucesiones $T_nv_j$ converge a cero. El máximo de sus normas tiende a cero porque el conjunto de índices es finito. La desigualdad recién probada concluye que $\norm{T_n}\to0$. Para $V=\{0\}$ todos los operadores son cero y la conclusión es directa.
\end{solucion}

\begin{ejercicio}[Lectura 22, ejercicio 3]\label{ej:lectura-22-3}
Sea $\mathcal A_n=\{T\in\Bcal(V,W):\dim T(V)\leq n\}$. Demuestre que $\mathcal A_n$ es cerrado en norma de operador.
\end{ejercicio}
\begin{solucion}
Suponga $T_j\in\mathcal A_n$ y $T_j\to T$ en norma. Si $\dim T(V)>n$, existen $v_1,\ldots,v_{n+1}$ tales que $Tv_1,\ldots,Tv_{n+1}$ son linealmente independientes. Por la proposición~\ref{ch22:independencia}, existe $\delta>0$ que conserva esa independencia bajo perturbaciones menores que $\delta$. Para cada $i$,
\[
 \norm{T_jv_i-Tv_i}\leq\norm{T_j-T}\norm{v_i}\to0.
\]
Hay un mismo $j$ suficientemente grande para que las $n+1$ diferencias sean menores que $\delta$. Entonces $T_jv_1,\ldots,T_jv_{n+1}$ son independientes, contradiciendo $\dim T_j(V)\leq n$. Así $T\in\mathcal A_n$. El criterio de sucesiones en el espacio normado $\Bcal(V,W)$ prueba la cerradura.
\end{solucion}

\begin{ejercicio}[Lectura 22, ejercicio 6]\label{ej:lectura-22-6}
Sea $V\subseteq\ell^1$ de dimensión finita. Demuestre que para todo $\eps>0$ existe $N\in\N$ tal que, para todo $x\in V$ y todo $\lambda\in\F$,
\[
 \norm{x+\lambda e_N}_1\geq|\lambda|+(1-\eps)\norm x_1.
\]
\end{ejercicio}
\begin{solucion}
Para cada $n$ considere la funcional de coordenada $\varphi_n(x)=x_n$ restringida a $V$. Es continua y $|x_n|\leq\norm x_1$. Como toda sucesión absolutamente sumable tiene términos que tienden a cero, $\varphi_n(x)\to0$ para cada $x\in V$. La versión del ejercicio~\ref{ej:lectura-22-2} con codominio $\F$ muestra que $\norm{\varphi_n}\to0$. Elija $N$ con $\norm{\varphi_N}\leq\eps/2$. Entonces $|x_N|\leq(\eps/2)\norm x_1$ para todos los $x\in V$ a la vez. Finalmente,
\begin{align*}
 \norm{x+\lambda e_N}_1
 &=\sum_{n\ne N}|x_n|+|x_N+\lambda|\\
 &\geq\norm x_1-|x_N|+|\lambda|-|x_N|\\
 &\geq|\lambda|+(1-\eps)\norm x_1.
\end{align*}
El índice $N$ depende de $V$ y $\eps$, pero no de $x$ ni de $\lambda$. Esa uniformidad es exactamente lo que aporta la dimensión finita.
\end{solucion}

\begin{ejercicio}[Lectura 23, ejercicio 1]\label{ej:lectura-23-1}
Si $f:V\to W$ es uniformemente continua en todo el espacio normado $V$ y $A\subseteq V$ es acotado, demuestre que $f(A)$ es acotado.
\end{ejercicio}
\begin{solucion}
Por continuidad uniforme existe $\delta>0$ tal que $\norm{x-y}<\delta$ implica $\norm{f(x)-f(y)}<1$. Elija $R>0$ con $\norm a\leq R$ para $a\in A$ y un entero $N$ con $R/N<\delta$. Para cada $a\in A$, divida el segmento del origen a $a$ en los puntos $a_j=ja/N$, $0\leq j\leq N$. Cada salto tiene norma $\norm{a_j-a_{j-1}}=\norm a/N<\delta$. Por tanto
\[
 \norm{f(a)-f(0)}\leq\sum_{j=1}^N\norm{f(a_j)-f(a_{j-1})}<N.
\]
Así $\norm{f(a)}\leq\norm{f(0)}+N$ para todos los $a\in A$. Es una sola cota porque $N$ se eligió a partir de $R$, no de cada punto. La prueba utiliza que el segmento completo está en el dominio; para funciones definidas solamente en un conjunto métrico arbitrario la afirmación no tendría este argumento.
\end{solucion}

\begin{ejercicio}[Lectura 23, ejercicio 2]\label{ej:lectura-23-2}
Si $W$ es de Banach, pruebe que un operador lineal $T:V\to W$ es compacto si y solo si manda conjuntos acotados en conjuntos totalmente acotados.
\end{ejercicio}
\begin{solucion}
La proposición~\ref{ch23:totalmente-acotado} ya da la demostración completa. Podemos expresar el núcleo del argumento mediante sucesiones. Si la imagen de algún acotado $A$ no fuera totalmente acotada, existiría $\eps>0$ y podríamos elegir inductivamente $y_n=Tx_n$ con $x_n\in A$ y $\norm{y_n-y_m}\geq\eps$ para $n\ne m$: cada nuevo punto se elige fuera de las bolas de radio $\eps$ centradas en los anteriores. Ninguna subsucesión de $(y_n)$ sería de Cauchy, contradiciendo compacidad. En sentido inverso, una sucesión en un conjunto totalmente acotado tiene una subsucesión de Cauchy, obtenida escogiendo sucesivamente bolas de radios tendiendo a cero que contengan infinitos términos. La completitud de $W$ transforma esa subsucesión de Cauchy en una convergente en $W$.
\end{solucion}

\begin{ejercicio}[Lectura 23, ejercicio 3]\label{ej:lectura-23-3}
Sean $S:V\to W$ y $T:W\to X$ operadores continuos. Si uno de ellos es compacto, demuestre que $TS$ es compacto.
\end{ejercicio}
\begin{solucion}
Sea $(v_n)$ acotada en $V$. Si $S$ es compacto, existe una subsucesión $Sv_{n_j}\to w$ en $W$. La continuidad de $T$ da $TSv_{n_j}\to Tw$ en $X$.

Si $T$ es compacto, la continuidad lineal de $S$ da $\norm{Sv_n}\leq\norm S\sup_m\norm{v_m}$, de modo que $(Sv_n)$ es acotada. La compacidad de $T$ aplicada a esta sucesión produce una subsucesión convergente de $(TSv_n)$. En ambos casos se verifica la definición. No se requiere que los espacios sean completos.
\end{solucion}

\begin{ejercicio}[Lectura 23, ejercicio 4]\label{ej:lectura-23-4}
Sea $T:V\to V$ compacto y $V$ infinito dimensional. Pruebe que no existe un operador continuo $S$ con $ST=TS=I$.
\end{ejercicio}
\begin{solucion}
Si existiera, el ejercicio~\ref{ej:lectura-23-3} haría compacto a $ST=I$. Por el teorema~\ref{ch23:identidad}, eso obligaría a que $V$ tuviera dimensión finita, contradicción. De hecho, basta la igualdad $ST=I$ para llegar a la contradicción. El resultado se refiere a una inversa continua: una biyección algebraica no garantiza automáticamente una inversa continua en espacios normados arbitrarios.
\end{solucion}

\chapter{Dualidad, Hahn--Banach y geometría convexa}

\begin{ejercicio}[Lectura 21, ejercicio 1]\label{ej:lectura-21-1}
Demuestre que el espacio de todas las funcionales lineales de $\F^n$ es isomorfo a $\F^n$.
\end{ejercicio}
\begin{solucion}
Una funcional queda determinada por $a_j=\varphi(e_j)$, porque $\varphi(x)=\sum_{j=1}^n a_jx_j$. Así $J(a)(x)=\sum_j a_jx_j$ es una biyección lineal entre $\F^n$ y su dual algebraico. No introducimos conjugados en esta fórmula porque queremos que $J$ sea lineal también cuando $\F=\C$.

Con la norma uniforme en el dominio de las funcionales,
\[
 |J(a)(x)|\leq\sum_j|a_j|\quad\text{si }\norm x_\infty\leq1.
\]
Elija $x_j=\overline{a_j}/|a_j|$ cuando $a_j\ne0$ y $x_j=1$ cuando $a_j=0$; en el caso real son simplemente signos. Entonces $a_jx_j=|a_j|$, y obtenemos $\norm{J(a)}=\sum_j|a_j|$. Todas las funcionales son continuas y $J$ es una isometría si el espacio de coeficientes lleva norma $\ell^1$. Como $\norm a_\infty\leq\sum_j|a_j|\leq n\norm a_\infty$, también es un isomorfismo para la norma uniforme, y el lema~\ref{ch19:coordenadas} permite cualquier otra norma en dimensión finita.
\end{solucion}

\begin{ejercicio}[Lectura 21, ejercicio 2]\label{ej:lectura-21-2}
Sea $\varphi\in V^*$ y $M=\ker\varphi$. Pruebe que
\[|\varphi(v)|=\norm\varphi\,\dist(v,M)\qquad(v\in V).\]
\end{ejercicio}
\begin{solucion}
Si $\varphi=0$, entonces $M=V$ y ambos lados son cero. Suponga $\varphi\ne0$. Para $m\in M$,
\[
 |\varphi(v)|=|\varphi(v-m)|\leq\norm\varphi\norm{v-m}.
\]
Tomar el ínfimo da $|\varphi(v)|\leq\norm\varphi\dist(v,M)$. Si $\varphi(v)=0$, $v\in M$ y ya hay igualdad. En otro caso, por la definición de norma existen $u_j$ con $\norm{u_j}\leq1$ y $|\varphi(u_j)|\to\norm\varphi>0$. Descartamos términos iniciales si hace falta para que $\varphi(u_j)\ne0$. El vector
\[
 m_j=v-\frac{\varphi(v)}{\varphi(u_j)}u_j
\]
pertenece a $M$. Por tanto
\[
 \dist(v,M)\leq\norm{v-m_j}
 \leq\frac{|\varphi(v)|}{|\varphi(u_j)|}.
\]
Al pasar al límite obtenemos la desigualdad inversa. No hemos supuesto que la norma de $\varphi$ se alcance.
\end{solucion}

\begin{ejercicio}[Lectura 21, ejercicio 3]\label{ej:lectura-21-3}
Sea $A\subseteq V$ no vacío y $U=\overline B(0,1)$. Pruebe
\[
 \cl{A+U}=\{v\in V:\dist(v,A)\leq1\}.
\]
Deduzca que $A+U$ es cerrado si y solo si todo $v$ con $\dist(v,A)=1$ admite $a\in A$ con $\norm{v-a}=1$.
\end{ejercicio}
\begin{solucion}
La función distancia a un conjunto no vacío es continua: la desigualdad triangular da $|\dist(x,A)-\dist(y,A)|\leq\norm{x-y}$. El conjunto del lado derecho es por ello cerrado. Contiene a $A+U$, así que también contiene su cierre.

Si $\dist(v,A)\leq1$, elija $a_n\in A$ con $\norm{v-a_n}<1+1/n$. Defina
\[
 u_n=\frac{v-a_n}{\max\{1,\norm{v-a_n}\}},\qquad z_n=a_n+u_n.
\]
Tenemos $\norm{u_n}\leq1$ y $z_n\in A+U$. Si $\norm{v-a_n}\leq1$, entonces $z_n=v$. Si es mayor que $1$, la diferencia entre $v-a_n$ y su normalización tiene norma $\norm{v-a_n}-1<1/n$. Por tanto $z_n\to v$, y queda demostrada la igualdad.

Cuando $\dist(v,A)<1$, la definición de ínfimo ya proporciona $a\in A$ con $\norm{v-a}<1$, de modo que $v\in A+U$. Solo pueden faltar en la suma puntos de distancia exactamente $1$. Para esos puntos, pertenecer a $A+U$ equivale a que exista un $a\in A$ con $\norm{v-a}\leq1$; como la distancia es $1$, tal desigualdad debe ser igualdad. Esto prueba el criterio. Si $A=\varnothing$, la suma es vacía y cerrada, y la identidad se interpreta con $\dist(v,\varnothing)=+\infty$.
\end{solucion}

\begin{ejercicio}[Lectura 21, ejercicio 4]\label{ej:lectura-21-4}
Sea $\norm\varphi=1$, $M=\ker\varphi$ y $U=\overline B(0,1)$. Demuestre
\[
 \cl{M+U}=\{v:|\varphi(v)|\leq1\}.
\]
Pruebe que $M+U$ es cerrado si y solo si existe $u\in U$ con $|\varphi(u)|=1$.
\end{ejercicio}
\begin{solucion}
La primera igualdad se sigue de los ejercicios~\ref{ej:lectura-21-2} y~\ref{ej:lectura-21-3} porque $\dist(v,M)=|\varphi(v)|$. Si $M+U$ es cerrado, elija $v$ con $\varphi(v)=1$, posible porque una funcional no nula es sobreyectiva sobre $\F$. La igualdad del cierre da $v=m+u$, con $m\in M$ y $u\in U$. De ahí $\varphi(u)=1$.

Recíprocamente, suponga $u_0\in U$ y $|\varphi(u_0)|=1$. Si $|\varphi(v)|\leq1$, defina
\[
 u=\frac{\varphi(v)}{\varphi(u_0)}u_0.
\]
Entonces $\norm u\leq1$ y $\varphi(v-u)=0$. Así $v=(v-u)+u\in M+U$. La suma coincide con el cerrado $\{v:|\varphi(v)|\leq1\}$.

La conclusión original escribe $|v|=1$ donde debe aparecer $|\varphi(v)|=1$: la condición relevante es que la funcional alcance su norma. No sucede siempre. En $c_0$ con norma uniforme,
\[
 \varphi(x)=\sum_{n=1}^{\infty}2^{-n}x_n
\]
tiene norma $1$, porque la desigualdad triangular da la cota y las sucesiones con sus primeros $N$ términos iguales a $1$ dan valores $1-2^{-N}$. Si $\norm x_\infty\leq1$ y $|\varphi(x)|=1$, entonces $\sum_n2^{-n}|x_n|=1$, lo que obliga a $|x_n|=1$ para todo $n$: cualquier déficit positivo en un término haría la suma menor que $1$. Eso contradice $x_n\to0$. En este ejemplo $M+U$ no es cerrado.
\end{solucion}

\begin{ejercicio}[Lectura 21, ejercicio 5]\label{ej:lectura-21-5}
Sea $V$ de Banach y $T\in\Bcal(V,V)$ con $\sum_{n=0}^{\infty}\norm{T^n}<\infty$. Demuestre que $I-T$ es un isomorfismo y calcule su inversa.
\end{ejercicio}
\begin{solucion}
El espacio $\Bcal(V,V)$ es completo. Por el teorema~\ref{ch19:series}, la serie $S=\sum_{n=0}^{\infty}T^n$ converge en norma de operador. Sus sumas parciales satisfacen
\[
 (I-T)\sum_{n=0}^NT^n=I-T^{N+1}
 =\left(\sum_{n=0}^NT^n\right)(I-T).
\]
Como la serie de normas converge, $\norm{T^{N+1}}\to0$. La continuidad de la composición, probada en el ejercicio~\ref{ej:lectura-20-5}, permite pasar al límite y obtener $(I-T)S=S(I-T)=I$. Por tanto $(I-T)^{-1}=S$ es continuo. En particular, si $\norm T<1$, la hipótesis se cumple porque $\norm{T^n}\leq\norm T^n$.
\end{solucion}

\begin{ejercicio}[Lectura 21, ejercicio 6]\label{ej:lectura-21-6}
Sean $V,W$ de Banach, con $V\ne\{0\}$, y $T:V\to W$ un isomorfismo. Si $S\in\Bcal(V,W)$ satisface $\norm{S-T}<1/\norm{T^{-1}}$, demuestre que $S$ es un isomorfismo. Deduzca que los isomorfismos forman un abierto de $\Bcal(V,W)$, incluyendo por separado el caso $V=W=\{0\}$.
\end{ejercicio}
\begin{solucion}
Supongamos $V\ne\{0\}$; cuando $V=W=\{0\}$ solo existe un operador y la conclusión de apertura es inmediata. Defina $E=T^{-1}(S-T)$. Entonces $\norm E\leq\norm{T^{-1}}\norm{S-T}<1$ y
\[
 S=T(I+E).
\]
El ejercicio~\ref{ej:lectura-21-5} aplicado a $-E$ da que $I+E$ es invertible. Así
\[
 S^{-1}=(I+E)^{-1}T^{-1},\qquad
 \norm{S^{-1}}\leq\frac{\norm{T^{-1}}}{1-\norm E}.
\]
Cada isomorfismo $T$ tiene, por tanto, una bola abierta de radio $1/\norm{T^{-1}}$ formada por isomorfismos. Esa es precisamente la definición de que el conjunto sea abierto. Si no existen isomorfismos, el conjunto vacío también es abierto.
\end{solucion}

\begin{ejercicio}[Lectura 21, ejercicio 7]\label{ej:lectura-21-7}
Sea $T$ un isomorfismo de $V$ sobre un subespacio de $W$. Pruebe que, si $\norm{T'-T}$ es suficientemente pequeña, $T'$ es un isomorfismo sobre su imagen.
\end{ejercicio}
\begin{solucion}
El teorema~\ref{ch21:cota-inferior} proporciona $a>0$ con $\norm{Tv}\geq a\norm v$. Si $\delta=\norm{T'-T}<a$, entonces
\[
 \norm{T'v}\geq\norm{Tv}-\norm{(T'-T)v}\geq(a-\delta)\norm v.
\]
También $\norm{T'v}\leq\norm{T'}\norm v$. El mismo teorema muestra que $T'$ es un isomorfismo de $V$ sobre $T'(V)$, con inversa de norma a lo sumo $1/(a-\delta)$. No se afirma que las imágenes de $T$ y $T'$ sean iguales.
\end{solucion}

\begin{ejercicio}[Lectura 21, ejercicio 8]\label{ej:lectura-21-8}
Pruebe que la evaluación canónica
\[
 J:V\longrightarrow V^{**},\qquad J(v)(\varphi)=\varphi(v),
\]
es una isometría lineal.
\end{ejercicio}
\begin{solucion}
Para cada $v$, la aplicación $J(v)$ es lineal en $\varphi$, y
\[
 |J(v)(\varphi)|=|\varphi(v)|\leq\norm\varphi\norm v.
\]
Así $J(v)\in V^{**}$ y $\norm{J(v)}\leq\norm v$. Además, para $\alpha,\beta\in\F$,
\[
 J(\alpha v+\beta w)(\varphi)
 =\varphi(\alpha v+\beta w)
 =\alpha J(v)(\varphi)+\beta J(w)(\varphi),
\]
de modo que $J$ es lineal. Si $v\ne0$, el corolario~\ref{ch21:funcional-normante} proporciona una funcional $\varphi$ de norma $1$ con $\varphi(v)=\norm v$. Por tanto
\[
 \norm{J(v)}\geq|J(v)(\varphi)|=\norm v.
\]
Para $v=0$ la igualdad es directa. En consecuencia $\norm{J(v)}=\norm v$ para todo $v$, y $J$ es inyectiva. Su sobreyectividad es una propiedad adicional, llamada reflexividad, que no se deduce de esta prueba.
\end{solucion}

\begin{ejercicio}[Lectura 24, ejercicio 1]\label{ej:lectura-24-1}
Demuestre que la intersección de una familia no vacía de conjuntos convexos es convexa.
\end{ejercicio}
\begin{solucion}
Sea $C=\bigcap_{A\in\mathcal A}A$. Si $x,y\in C$ y $0\leq t\leq1$, entonces $x,y\in A$ para cada $A\in\mathcal A$. Su convexidad implica $(1-t)x+ty\in A$ para cada miembro de la familia, y por ello el punto pertenece a $C$. Si $C$ es vacío, la condición no tiene pares de puntos que comprobar y también se cumple.
\end{solucion}

\begin{ejercicio}[Lectura 24, ejercicio 2]\label{ej:lectura-24-2}
Si $A,B\subseteq V$ y $C\subseteq W$ son convexos, $\alpha\in\R$ y $T:V\to W$ es lineal, pruebe que $\alpha A$, $A+B$, $T(A)$ y $T^{-1}(C)$ son convexos.
\end{ejercicio}
\begin{solucion}
Para $a_1,a_2\in A$,
\[
 (1-t)\alpha a_1+t\alpha a_2=\alpha((1-t)a_1+ta_2)\in\alpha A.
\]
Esto vale también para $\alpha=0$ y para $\alpha<0$, porque los coeficientes interiores siguen siendo $1-t,t$. Para dos puntos $a_1+b_1,a_2+b_2\in A+B$,
\[
 (1-t)(a_1+b_1)+t(a_2+b_2)
 =((1-t)a_1+ta_2)+((1-t)b_1+tb_2)\in A+B.
\]
Para la imagen, $(1-t)T(a_1)+tT(a_2)=T((1-t)a_1+ta_2)\in T(A)$. Para la preimagen, si $Tx,Ty\in C$, la linealidad da $T((1-t)x+ty)=(1-t)Tx+tTy\in C$, y por tanto $(1-t)x+ty\in T^{-1}(C)$. Los casos de conjuntos vacíos satisfacen la definición de forma vacua.
\end{solucion}

\begin{ejercicio}[Lectura 24, ejercicio 3]\label{ej:lectura-24-3}
Sean $A_1,\ldots,A_n$ conjuntos convexos no vacíos. Demuestre que $\operatorname{co}(A_1\cup\cdots\cup A_n)$ está formado por las combinaciones $\sum_{j=1}^n\alpha_jx_j$, con $x_j\in A_j$, $\alpha_j\geq0$ y $\sum_j\alpha_j=1$.
\end{ejercicio}
\begin{solucion}
Toda combinación de la forma indicada es, por definición, una combinación convexa de puntos de la unión, así que pertenece a su envolvente. Para la otra inclusión, tome una combinación convexa finita de puntos de la unión y asigne cada punto a uno de los $A_j$ que lo contenga. Agrupe los términos según esa asignación. Sea $\alpha_j$ la suma de sus coeficientes. Si $\alpha_j>0$, el grupo se escribe $\alpha_jx_j$, donde $x_j$ es una combinación convexa de puntos de $A_j$ y pertenece a $A_j$ por convexidad. Si $\alpha_j=0$, elija cualquier $x_j\in A_j$; su contribución será cero. La suma de los $\alpha_j$ es $1$, y hemos obtenido la representación buscada.

La condición de no vacío permite elegir $x_j$ incluso cuando su coeficiente sea cero. Si alguno de los conjuntos originales es vacío, debe omitirse de la lista; exigir literalmente un punto de un conjunto vacío haría falsa la representación.
\end{solucion}

\begin{ejercicio}[Lectura 24, ejercicio 4: revisión y solución]\label{ej:lectura-24-4}
La lectura propone, para conjuntos cerrados $A,B$, las afirmaciones: (i) si ambos son compactos, $\operatorname{co}(A\cup B)$ es compacto; (ii) si $A$ es compacto y $B$ acotado, $\operatorname{co}(A\cup B)$ es cerrado. Examine ambas afirmaciones y pruebe sus versiones correctas. Demuestre además que la envolvente convexa de $\{(0,0)\}\cup\{(1,y):y\in\R\}$ no es cerrada.
\end{ejercicio}
\begin{solucion}
\emph{Primero, el defecto en las hipótesis.} Tal como están escritas, ambas afirmaciones son falsas en espacios normados generales: falta convexidad. En $\ell^2$, tome
\[
 A=\{0\}\cup\{e_n/n:n\in\N\},\qquad B=\{0\}.
\]
El conjunto $A$ es compacto porque $e_n/n\to0$: dada una cubierta abierta, el miembro que contiene a $0$ cubre toda la cola y solo quedan finitos puntos que cubrir. También $B$ es compacto. Toda combinación convexa finita de puntos de $A$ tiene soporte finito. Sin embargo,
\[
 z_N=\sum_{n=1}^N2^{-n}\frac{e_n}{n}+2^{-N}0
 \in\operatorname{co}(A)
\]
y $z_N$ converge en $\ell^2$ a $z=\sum_{n=1}^{\infty}2^{-n}e_n/n$, que tiene infinitas coordenadas no nulas y no pertenece a $\operatorname{co}(A)$. La convergencia se verifica, por ejemplo, acotando la norma de la cola por $\sum_{n>N}2^{-n}/n\to0$. Así $\operatorname{co}(A\cup B)$ no es cerrado y, en consecuencia, tampoco compacto. Este único ejemplo contradice los dos incisos originales.

\emph{Versión corregida del inciso (i).} Suponga ahora que $A,B$ son no vacíos, compactos y convexos. Por el ejercicio~\ref{ej:lectura-24-3},
\[
 \operatorname{co}(A\cup B)
 =\{(1-t)a+tb:a\in A,\ b\in B,\ 0\leq t\leq1\}.
\]
Es la imagen del compacto $A\times B\times[0,1]$ por la aplicación continua $(a,b,t)\mapsto(1-t)a+tb$. Por ello es compacto.

\emph{Versión corregida del inciso (ii).} Suponga $A$ compacto convexo no vacío y $B$ cerrado, acotado, convexo y no vacío. Si $z_n=(1-t_n)a_n+t_nb_n$ converge a $z$, la compacidad de $A\times[0,1]$ permite pasar a una subsucesión con $a_n\to a\in A$ y $t_n\to t\in[0,1]$. Si $t=0$, la acotación de $B$ da $t_nb_n\to0$, y entonces $z=a\in A$. Si $t>0$, para índices grandes $t_n>0$ y
\[
 b_n=\frac{z_n-(1-t_n)a_n}{t_n}
 \longrightarrow b=\frac{z-(1-t)a}{t}.
\]
Como $B$ es cerrado, $b\in B$, y $z=(1-t)a+tb$ pertenece a la envolvente. Por el criterio de sucesiones, esta es cerrada. Si uno de los conjuntos es vacío, la envolvente se reduce al otro conjunto convexo, y las conclusiones correspondientes son directas.

\emph{El ejemplo de la recta vertical.} Sea $A=\{(0,0)\}$ y $B=\{(1,y):y\in\R\}$. Ambos son cerrados convexos y $A$ es compacto, pero $B$ no es acotado. Sus combinaciones son $(1-t)(0,0)+t(1,y)=(t,ty)$. De ahí
\[
 \operatorname{co}(A\cup B)
 =\bigl((0,1]\times\R\bigr)\cup\{(0,0)\}.
\]
Los puntos $(1/n,1)$ pertenecen a la envolvente y convergen a $(0,1)$, que no pertenece. La acotación de $B$ en la versión corregida del inciso (ii) no puede eliminarse.
\end{solucion}

\begin{ejercicio}[Lectura 24, ejercicio 5]\label{ej:lectura-24-5}
Si $A$ es convexo, demuestre que $C=\bigcup_{s\geq0}sA$ es un cono convexo.
\end{ejercicio}
\begin{solucion}
Si $A$ es vacío, $C$ es vacío y satisface la condición de estabilidad de manera vacua. Suponga $A\ne\varnothing$, por lo que $0\in C$. Sean $u=sa$, $v=tb$, con $s,t\geq0$ y $a,b\in A$, y sean $\alpha,\beta\geq0$. Ponga $q=\alpha s+\beta t$. Si $q=0$, ambos productos son cero y $\alpha u+\beta v=0\in C$. Si $q>0$,
\[
 \alpha u+\beta v
 =q\left(\frac{\alpha s}{q}a+\frac{\beta t}{q}b\right).
\]
Los dos coeficientes interiores son no negativos y suman $1$, por lo que el paréntesis pertenece a $A$. Así el resultado pertenece a $qA\subseteq C$.
\end{solucion}

\begin{ejercicio}[Lectura 24, ejercicio 6]\label{ej:lectura-24-6}
Demuestre que la envolvente convexa cerrada de un compacto no vacío $A$ es el cierre de su envolvente convexa. Determine si la compacidad es necesaria.
\end{ejercicio}
\begin{solucion}
No es necesaria: la igualdad vale para cualquier subconjunto de un espacio normado, incluso vacío. Por la proposición~\ref{ch24:interior-cierre}, $\cl{\operatorname{co}(A)}$ es cerrado y convexo, y contiene a $A$. Por eso contiene a la envolvente convexa cerrada, que es el menor conjunto con esas propiedades. En sentido inverso, cualquier cerrado convexo $D$ que contenga a $A$ contiene $\operatorname{co}(A)$ por su convexidad y luego contiene su cierre por ser cerrado. Al intersectar todos esos $D$, se obtiene que la envolvente convexa cerrada contiene $\cl{\operatorname{co}(A)}$. Las dos inclusiones dan igualdad. Esta afirmación es de cerradura y convexidad, no una afirmación de compacidad de la envolvente.
\end{solucion}

\begin{ejercicio}[Lectura 24, ejercicio 7]\label{ej:lectura-24-7}
Sean $A,B$ convexos y disjuntos, y sea $x\in V$. Pruebe que
\[
 \operatorname{co}(A\cup\{x\})\cap B=\varnothing
 \quad\text{o}\quad
 \operatorname{co}(B\cup\{x\})\cap A=\varnothing.
\]
\end{ejercicio}
\begin{solucion}
Si $A$ o $B$ es vacío, una de las intersecciones es vacía directamente. Suponga ambos no vacíos y, buscando una contradicción, que ambas intersecciones contienen puntos. Por la convexidad y el ejercicio~\ref{ej:lectura-24-3} existen $a,a_0\in A$, $b,b_0\in B$ y $s,t\in[0,1]$ tales que
\[
 a=(1-s)b_0+sx,\qquad b=(1-t)a_0+tx.
\]
Los números $s,t$ son positivos: si $s=0$, $a=b_0\in A\cap B$, y si $t=0$, $b=a_0\in A\cap B$. Multiplique la primera igualdad por $t$ y la segunda por $s$, y elimine el término $stx$. Resulta
\[
 ta+s(1-t)a_0=sb+t(1-s)b_0.
\]
La suma de los coeficientes de cada lado es la misma, $q=t+s-st>0$.
Al dividir entre $q$, el lado izquierdo es una combinación convexa de puntos de $A$ y el derecho una combinación convexa de puntos de $B$. Por tanto su valor común pertenece a $A\cap B$, contradicción. Al menos una de las dos intersecciones debe ser vacía.
\end{solucion}

\chapter{Espacios de funciones y convergencia uniforme}

\begin{ejercicio}[Lectura 26, ejercicio 1]\label{ej:lectura-26-1}\label{ch26:ej1}
Sea $V$ Banach y $f_n:S\to V$. Demuestre que $(f_n)$ converge uniformemente si y solo si para todo $\eps>0$ existe $N$ tal que $\norm{f_n(x)-f_m(x)}<\eps$ para todos $n,m\ge N$ y $x\in S$.
\end{ejercicio}
\begin{solucion}
Si $f_n\rightrightarrows f$, elegimos $N$ con $\norm{f_n(x)-f(x)}<\eps/2$ para todo $n\ge N$ y $x$. Para $n,m\ge N$, la desigualdad triangular da el resultado.

Recíprocamente, para cada $x$ la hipótesis hace de $(f_n(x))$ una sucesión de Cauchy en $V$. La completitud define $f(x)=\lim_nf_n(x)$. Dado $\eps>0$, elegimos $N$ usando $\eps/2$ en la hipótesis. Fijamos $n\ge N$ y $x$ y hacemos $m\to\infty$; obtenemos $\norm{f_n(x)-f(x)}\le\eps/2<\eps$. Como el índice $N$ es independiente de $x$, la convergencia es uniforme. No necesitamos que cada $f_n$ sea acotada: solo controlamos las diferencias de los términos suficientemente avanzados.
\end{solucion}

\begin{ejercicio}[Lectura 26, ejercicio 2]\label{ej:lectura-26-2}
Sean $(f_n)$ y $(g_n)$ sucesiones de funciones escalares sobre $S$, acotadas en $B(S,\F)$, que convergen uniformemente. Demuestre que $(f_ng_n)$ converge uniformemente. Dé un contraejemplo al suprimir el control de acotación.
\end{ejercicio}
\begin{solucion}
Escribamos $f_n\rightrightarrows f$ y $g_n\rightrightarrows g$. Existen $A,B<\infty$ con $\norm{f_n}_\infty\le A$ y $\norm{g_n}_\infty\le B$. Al pasar al límite en cada punto, $\norm f_\infty\le A$ y $\norm g_\infty\le B$. La identidad
\[
 f_ng_n-fg=f_n(g_n-g)+(f_n-f)g
\]
proporciona
\[
 \norm{f_ng_n-fg}_\infty
 \le A\norm{g_n-g}_\infty+B\norm{f_n-f}_\infty\longrightarrow0.
\]
Si las funciones no están acotadas, tomemos $S=\R$, $f_n(x)=x$ y $g_n(x)=1/n$. Ambas sucesiones convergen uniformemente, respectivamente a $f(x)=x$ y a cero. Sin embargo, $f_n(x)g_n(x)=x/n$ no converge uniformemente a cero porque su error no está acotado en $\R$ para ningún $n$. Es suficiente que falle el control de uno de los factores.

Si todas las funciones pertenecen a $B(S,\F)$ y la convergencia es uniforme, la acotación de toda la sucesión se deduce automáticamente: una cola está a distancia menor que uno del límite, y los primeros términos son un conjunto finito de funciones acotadas. La hipótesis expresa, por tanto, el contexto de funciones acotadas en el cual el producto es continuo.
\end{solucion}

\begin{ejercicio}[Lectura 26, ejercicio 3: enunciado corregido]\label{ej:lectura-26-3}
Sean $V,W$ Banach, $f_n\rightrightarrows f$ y $\norm{f_n}_\infty\le\alpha$. La lectura propone que la continuidad de $g$ en la bola de radio $\alpha$ basta para que $g\circ f_n\rightrightarrows g\circ f$. Examine la afirmación y pruebe una versión válida.
\end{ejercicio}
\begin{solucion}
La condición adecuada es que $g$ sea \emph{uniformemente continua} en un conjunto que contenga todas las imágenes $f_n(S)$ y $f(S)$. Con esa hipótesis, dado $\eps>0$ existe $\delta>0$ tal que
$\norm{u-v}<\delta$ implica $\norm{g(u)-g(v)}<\eps$. La convergencia uniforme da un índice $N$ con $\norm{f_n(x)-f(x)}<\delta$ para todo $n\ge N$ y $x$. Al aplicar $g$ obtenemos el resultado. En particular, basta continuidad de $g$ sobre un conjunto compacto que contenga todas esas imágenes.

Si la bola del enunciado es abierta, incluso puede ocurrir que $f$ tome valores en su frontera y $g\circ f$ no esté definido. La afirmación falla también cuando todas las composiciones sí están definidas: en $S=(0,1)$, $V=W=\R$, tomemos
\[
 f(x)=x,\qquad f_n(x)=(1-1/n)x,\qquad g(t)=\frac1{1-t}\quad(-1<t<1).
\]
Tenemos $\norm{f_n-f}_\infty=1/n$ y $\norm{f_n}_\infty\le1$, pero para cada $n$ el valor $g(f(x))$ tiende a infinito cuando $x\uparrow1$, mientras $g(f_n(x))$ permanece acotado por $n$. La diferencia no está uniformemente acotada.

Interpretar la bola como cerrada tampoco salva el enunciado en dimensión infinita. En $V=\ell^2$, sea $(e_k)$ la sucesión de vectores canónicos. En la bola cerrada unidad definimos
\[
 g(u)=\sum_{k=1}^\infty k\max\{0,1-4\norm{u-e_k}_2\}.
\]
Esta función es continua. Para justificarlo, observamos que $\norm{e_j-e_k}_2=\sqrt2$ si $j\ne k$. Cada punto tiene una vecindad que intersecta a lo sumo una de las bolas $B(e_k,1/4)$: por ejemplo, si dos de esas bolas intersectaran $B(u,1/10)$, la distancia entre sus centros sería menor que $2(1/4+1/10)<\sqrt2$. Así, cerca de cada punto la suma es cero o coincide con una sola función continua. Definamos en $S=\N$
\[
 f(k)=e_k,\qquad f_n(k)=(1-1/n)e_k.
\]
Las imágenes están en la bola cerrada unidad y $\norm{f_n-f}_\infty=1/n$. Para $n>4$, $g(f(k))=k$ y $g(f_n(k))=k(1-4/n)$. La diferencia $4k/n$ no es acotada en $k$. Esto demuestra por qué continuidad sobre una bola acotada no sustituye a continuidad uniforme.
\end{solucion}

\begin{ejercicio}[Lectura 26, ejercicio 4]\label{ej:lectura-26-4}
Sean $V$ Banach y $f_n:S\to V$. Suponga que $\norm{f_{n+1}(x)-f_n(x)}\le M_n$ para todo $x$, donde $\sum_nM_n<\infty$. Demuestre que $(f_n)$ converge uniformemente.
\end{ejercicio}
\begin{solucion}
Para $q>p$, la identidad telescópica da
\[
 f_q(x)-f_p(x)=\sum_{n=p}^{q-1}(f_{n+1}(x)-f_n(x)),
 \qquad \norm{f_q(x)-f_p(x)}\le\sum_{n=p}^{q-1}M_n.
\]
El último miembro es pequeño para todo $q>p$ cuando $p$ es grande, con independencia de $x$. El criterio de Cauchy del ejercicio~\ref{ej:lectura-26-1} concluye la prueba. Además, al pasar al límite $q\to\infty$ obtenemos la estimación útil $\norm{f-f_p}_\infty\le\sum_{n=p}^\infty M_n$.
\end{solucion}

\begin{ejercicio}[Lectura 26, ejercicio 5]\label{ej:lectura-26-5}
Estudie la relación entre convergencia absoluta y uniforme:
\begin{enumerate}[label=(\roman*)]
\item En el dominio $[0,1]$, para $x>0$ sea $u_n(x)=\sin(\pi/x)\mathbf1_{[1/(n+1),1/n]}(x)$ y defina $u_n(0)=0$. Analice la convergencia absoluta puntual de $\sum_nu_n$ y la convergencia uniforme.
\item Demuestre que $\sum_{n=1}^\infty(-1)^n(x^2+n)/n^2$ converge uniformemente en $[0,1]$, pero no absolutamente en ningún punto.
\end{enumerate}
\end{ejercicio}
\begin{solucion}
\begin{enumerate}[label=(\roman*)]
\item En cada punto $x>0$ hay a lo sumo dos intervalos $[1/(n+1),1/n]$ que contienen a $x$. Si hay dos, $x=1/m$ y $\sin(\pi/x)=\sin(m\pi)=0$. Por tanto, en cada punto hay a lo sumo un término no nulo, y $\sum_n|u_n(x)|<\infty$. En cero todos los términos son cero. La suma puntual es $s(x)=\sin(\pi/x)$ para $x>0$ y $s(0)=0$.

Para cada $n$, el punto $x_n=1/(n+1/2)$ está en el interior del intervalo que soporta $u_n$ y satisface $|u_n(x_n)|=1$. Así $\norm{u_n}_\infty=1$. Una serie uniformemente convergente debe tener términos que tiendan uniformemente a cero, pues $u_n=s_n-s_{n-1}$; esta condición falla. Por tanto, la serie no converge uniformemente. La propia sucesión $(u_n)$ converge puntualmente a cero y tampoco converge uniformemente. Esta distinción aclara el uso de «convergencia absoluta de la sucesión» en la lectura: la convergencia absoluta corresponde propiamente a una \emph{serie}.
\item Descomponemos la serie como
\[
 x^2\sum_{n=1}^\infty\frac{(-1)^n}{n^2}
 +\sum_{n=1}^\infty\frac{(-1)^n}{n}.
\]
La primera converge uniformemente por el criterio M, pues $|x^2/n^2|\le1/n^2$. La segunda es una serie numérica convergente por el criterio de Leibniz y sus sumas parciales, vistas como funciones constantes de $x$, convergen uniformemente. Para verlo directamente, el error de la segunda serie es a lo sumo $1/(N+1)$, independientemente de $x$. La suma de dos errores que tienden uniformemente a cero también tiende uniformemente a cero. Sin embargo,
\[
 \sum_{n=1}^\infty\left|(-1)^n\frac{x^2+n}{n^2}\right|
 =\sum_{n=1}^\infty\left(\frac{x^2}{n^2}+\frac1n\right)=\infty
\]
para todo $x\in[0,1]$.
\end{enumerate}
\end{solucion}

\begin{ejercicio}[Lectura 26, ejercicio 6]\label{ej:lectura-26-6}
Si $f_n:X\to V$ son continuas y convergen uniformemente a $f$, demuestre que $f_n(x_n)\to f(x)$ siempre que $x_n\to x$. ¿Vale el recíproco?
\end{ejercicio}
\begin{solucion}
El límite uniforme $f$ es continuo por el teorema~\ref{ch12:limite-continuo}, cuyo argumento de tres términos también vale sin acotación global. Entonces
\[
 \norm{f_n(x_n)-f(x)}
 \le\norm{f_n(x_n)-f(x_n)}+\norm{f(x_n)-f(x)}\longrightarrow0.
\]
El primer término está controlado por la convergencia uniforme y el segundo por continuidad de $f$.

El recíproco falla en un espacio no compacto. En $X=\R$ tomemos
\[
 f_n(t)=\max\{0,1-|t-n|\},\qquad f(t)=0.
\]
Si $x_n\to x$, la sucesión $(x_n)$ es acotada. Para $n$ suficientemente grande, $|x_n-n|>1$, y por ello $f_n(x_n)=0$. Se cumple la propiedad pedida, pero $\norm{f_n}_\infty=1$ para todo $n$.

En un espacio métrico compacto sí vale el recíproco. Si no hubiera convergencia uniforme, existirían $\eps>0$, índices estrictamente crecientes $n_j$ y puntos $y_j$ con $\norm{f_{n_j}(y_j)-f(y_j)}\ge\eps$. La propiedad para las sucesiones constantes implica primero $f_n(y)\to f(y)$ en cada punto. Además, implica continuidad de $f$: si $z_j\to z$ y $f(z_j)\not\to f(z)$, elegimos índices crecientes $m_j$ tales que $\norm{f_{m_j}(z_j)-f(z_j)}<1/j$, y definimos $x_{m_j}=z_j$, $x_n=z$ en los demás índices. Entonces $x_n\to z$ y la propiedad obliga a $f_{m_j}(z_j)\to f(z)$, contradicción. Por compacidad podemos extraer $y_{j_k}\to y$. Definimos $x_{n_{j_k}}=y_{j_k}$ y $x_n=y$ en los demás índices. La propiedad da $f_{n_{j_k}}(y_{j_k})\to f(y)$; la continuidad recién probada da $f(y_{j_k})\to f(y)$. Contradice la cota inferior $\eps$.
\end{solucion}

\begin{ejercicio}[Lectura 26, ejercicio 7: Dirichlet uniforme]\label{ej:lectura-26-7}
Sean $V$ Banach, $u_n:S\to V$ y $v_n:S\to\R$. Suponga que, para algún $M<\infty$,
\[
 \sup_{x\in S}\left\|\sum_{k=1}^nu_k(x)\right\|\le M\quad(n\ge1),
\]
que $v_n\rightrightarrows0$ y que $v_{n+1}(x)\le v_n(x)$ para todo $x$. Demuestre que $\sum_nv_nu_n$ converge uniformemente.
\end{ejercicio}
\begin{solucion}
La monotonía y el límite cero implican $v_n(x)\ge0$. Es importante que las sumas parciales de $u_n$ estén acotadas por un \emph{mismo} $M$, independiente tanto de $n$ como de $x$; la acotación puntual no basta para esta prueba.

Fijemos $p\le q$ y pongamos $A_k(x)=\sum_{j=p}^ku_j(x)$, con $A_{p-1}=0$. Como $A_k$ es diferencia de dos sumas parciales, $\norm{A_k(x)}\le2M$. La identidad de sumación por partes es
\[
 \sum_{k=p}^qv_k u_k
 =v_qA_q+\sum_{k=p}^{q-1}(v_k-v_{k+1})A_k.
\]
Para verificarla, sustituimos $u_k=A_k-A_{k-1}$ y agrupamos el coeficiente de cada $A_k$: para $p\le k<q$ es $v_k-v_{k+1}$ y para $k=q$ es $v_q$. Como todos esos coeficientes son no negativos,
\[
 \left\|\sum_{k=p}^qv_k(x)u_k(x)\right\|
 \le2M\left(v_q(x)+\sum_{k=p}^{q-1}(v_k(x)-v_{k+1}(x))\right)
 =2Mv_p(x).
\]
La convergencia uniforme $v_p\to0$ hace pequeña esta cota para todo $q\ge p$ y todo $x$. El criterio de Cauchy uniforme concluye la prueba. Cuando el espacio de escalares de $V$ es $\C$, los multiplicadores $v_n$ siguen siendo reales: el orden $v_{n+1}\le v_n$ no es un orden de números complejos arbitrarios.
\end{solucion}

\begin{ejercicio}[Lectura 26, ejercicio 8: condensación uniforme de Cauchy]\label{ej:lectura-26-8}
Sea $(u_n)$ una sucesión de funciones reales sobre $S$, decreciente respecto de $n$ en cada punto. Demuestre que $\sum_{n=1}^\infty u_n$ y $\sum_{k=1}^\infty 2^ku_{2^k}$ convergen uniformemente o dejan de hacerlo simultáneamente.
\end{ejercicio}
\begin{solucion}
Primero atendemos al signo, que el enunciado no especifica. Si para algún $x$ y $n_0$ se tiene $u_{n_0}(x)<0$, la monotonía implica $u_n(x)\le u_{n_0}(x)<0$ para todo $n\ge n_0$. La primera serie diverge en ese punto y los términos $2^ku_{2^k}(x)$ de la segunda ni siquiera tienden a cero. Ninguna puede converger uniformemente. Por tanto, solo queda considerar el caso $u_n(x)\ge0$ para todos $n,x$.

Para cada $k\ge0$, la monotonía produce las desigualdades de bloques
\[
 2^ku_{2^{k+1}}(x)
 \le\sum_{n=2^k}^{2^{k+1}-1}u_n(x)
 \le2^ku_{2^k}(x).
\]
Si la serie condensada converge uniformemente, dado $\eps>0$ elegimos $K$ tal que toda suma de sus términos con índices al menos $K$ sea menor que $\eps$, uniformemente en $x$. Para $q\ge p\ge2^K$, sea $L$ tal que $q<2^{L+1}$. Por no negatividad,
\[
 0\le\sum_{n=p}^qu_n(x)
 \le\sum_{k=K}^L\sum_{n=2^k}^{2^{k+1}-1}u_n(x)
 \le\sum_{k=K}^L2^ku_{2^k}(x)<\eps.
\]
Así la primera serie satisface el criterio de Cauchy uniforme.

Recíprocamente, para $k\ge1$ cada término del bloque $2^{k-1}\le n<2^k$ es al menos $u_{2^k}$. Por ello
\[
 2^ku_{2^k}(x)\le2\sum_{n=2^{k-1}}^{2^k-1}u_n(x),
\]
y, sumando desde $k=K$ hasta $L$,
\[
 0\le\sum_{k=K}^L2^ku_{2^k}(x)
 \le2\sum_{n=2^{K-1}}^{2^L-1}u_n(x).
\]
Si la primera serie converge uniformemente, la última cota puede hacerse menor que $\eps$ para todo $L\ge K$ y todo $x$. La serie condensada es entonces uniformemente de Cauchy. Omitir o incluir su término $k=0$ no altera la convergencia.
\end{solucion}

\begin{ejercicio}[Ampliación: iteraciones que revelan la solución]\label{ej:ampl-picard-1}
Resuelva $u'=u+t$, $u(0)=1$, en $[0,2]$ mediante Picard. Partiendo de $u_0=1$, encuentre una fórmula para $u_n$ y una cota uniforme explícita del error. Explique por qué la estimación directa del operador no basta para aplicar Banach en la norma uniforme.
\end{ejercicio}
\begin{solucion}
La función $f(t,x)=x+t$ es continua y Lipschitz en $x$ con $L=1$. Con $T=2$ la estimación directa da $LT=2$. Más aún, para dos funciones cuya diferencia es la constante $1$, la diferencia de sus imágenes es $t$, de norma $2$; el operador no es contracción en esa norma.

La recurrencia es $u_{n+1}(t)=1+\int_0^t(u_n(s)+s)\,ds$. Afirmamos que, para $n\ge0$,
\[
 u_n(t)=2\sum_{k=0}^n\frac{t^k}{k!}-1-t+\frac{t^{n+1}}{(n+1)!}.
\]
Para $n=0$ el lado derecho es $1$. Si la fórmula vale para $n$, sumar $s$ cancela el término $-s$, e integrar da
\[
 u_{n+1}(t)=1+2\sum_{k=0}^n\frac{t^{k+1}}{(k+1)!}-t
 +\frac{t^{n+2}}{(n+2)!},
\]
que es la misma fórmula con $n+1$. La serie exponencial converge uniformemente en $[0,2]$ por el criterio M, y $t^{n+1}/(n+1)!\to0$ uniformemente. Por tanto
\[
 u(t)=2e^t-1-t.
\]
La derivación directa verifica $u'=2e^t-1=u+t$ y $u(0)=1$. Finalmente,
\[
 |u(t)-u_n(t)|
 \le2\sum_{k=n+1}^{\infty}\frac{2^k}{k!}
 \le2e^2\frac{2^{n+1}}{(n+1)!}.
\]
En la última desigualdad se usa $(n+1+j)!\ge(n+1)!j!$. Es una cota válida para todo $t$; no presupone que ya conozcamos el punto donde ocurre el mayor error.
\end{solucion}

\begin{ejercicio}[Ampliación: una ecuación no lineal en un intervalo largo]\label{ej:ampl-picard-2}
Para $u'=\sin u+t$, $u(0)=0$, en $[0,2]$, verifique las hipótesis de Picard, calcule $u_1,u_2$ desde $u_0=0$ y obtenga cotas del error mediante la norma $\norm v_1=\sup_{0\le t\le2}e^{-t}|v(t)|$ y mediante el factorial.
\end{ejercicio}
\begin{solucion}
El valor medio da $|\sin x-\sin y|\le|x-y|$, de modo que $L=1$. La función es continua en todo el cilindro y el teorema asegura una solución única. Las dos primeras aproximaciones son
\[
 u_1(t)=t^2/2,\qquad
 u_2(t)=t^2/2+\int_0^t\sin(s^2/2)\,ds.
\]
No necesitamos una primitiva elemental de la última integral: define una función continua y puede aproximarse mediante integración.

Para $\beta=1$, la constante de contracción es $q=1-e^{-2}$. Como la derivada de $e^{-t}t^2/2$ es $e^{-t}t(2-t)/2\ge0$ en $[0,2]$,
$\norm{u_1-u_0}_1=2e^{-2}$. La cota de Banach y la equivalencia de normas dan
\[
 \norm{u-u_n}_\infty
 \le e^2\frac{q^n}{1-q}\norm{u_1-u_0}_1
 =2e^2(1-e^{-2})^n.
\]
Por otra parte, $M=\max_{[0,2]}|f(t,0)|=2$, y la proposición~\ref{ampl:picard-error} da
\[
 \norm{u-u_n}_\infty\le4e^2\frac{2^n}{(n+1)!}.
\]
Ambas cotas garantizan convergencia. La segunda utiliza la estructura integral repetida y no sólo una constante de contracción; su factorial ofrece un control más fino para índices grandes.
\end{solucion}

\begin{ejercicio}[Ampliación: estabilidad frente a datos perturbados]\label{ej:ampl-picard-3}
Sean $u,v$ las soluciones en $I$ de $u'=f(t,u)$, $v'=g(t,v)$ con $u(t_0)=\xi$, $v(t_0)=\eta$. Suponga las hipótesis de Picard para ambas ecuaciones, que $f$ tiene constante Lipschitz $L$, y que $\norm{f(t,x)-g(t,x)}_\infty\le\delta$ para todo $t,x$. Demuestre, con $r=|t-t_0|$, que para $L>0$,
\[
 \norm{u(t)-v(t)}_\infty
 \le e^{Lr}\norm{\xi-\eta}_\infty
       +\delta\frac{e^{Lr}-1}{L}.
\]
Obtenga la versión para $L=0$ y deduzca dependencia continua respecto del dato inicial.
\end{ejercicio}
\begin{solucion}
Trabajamos primero a la derecha de $t_0$. Ponga $h(r)=\norm{u(t_0+r)-v(t_0+r)}_\infty$ y $A=\norm{\xi-\eta}_\infty$. Restar las ecuaciones integrales y añadir y restar $f(t,v(t))$ produce
\[
 h(r)\le A+\delta r+L\int_0^r h(s)\,ds=:H(r).
\]
La función $h$ es continua, luego $H$ es derivable y
$H'(r)=\delta+Lh(r)\le\delta+LH(r)$. Multiplicar por $e^{-Lr}$ da
\[
 (e^{-Lr}H(r))'\le\delta e^{-Lr}.
\]
Integrando entre $0$ y $r$, con $H(0)=A$, obtenemos
$H(r)\le Ae^{Lr}+\delta(e^{Lr}-1)/L$ si $L>0$. Como $h\le H$, esta es la cota buscada. A la izquierda de $t_0$, usamos $h(r)=\norm{u(t_0-r)-v(t_0-r)}_\infty$ y la orientación de la integral desaparece al tomar normas; se obtiene la misma desigualdad.

Para $L=0$, la desigualdad integral inicial ya dice $h(r)\le A+\delta r$. Si $f=g$, entonces $\delta=0$, y en todo el intervalo
\[
 \norm{u-v}_\infty\le e^{LT}\norm{\xi-\eta}_\infty.
\]
Por ello, datos iniciales próximos producen soluciones uniformemente próximas. La estimación también cuantifica el efecto de modificar la ecuación, siempre que la perturbación esté uniformemente controlada.
\end{solucion}

\begin{ejercicio}[Ampliación: versión local y explosión de una solución]\label{ej:ampl-picard-4}
Sea $f$ continua en $[t_0-a,t_0+a]\times\{x:\norm{x-\xi}_\infty\le b\}$, con $a,b>0$, acotada allí por $M$ y Lipschitz en $x$ con constante $L$. Si $0<h\le a$, $Mh\le b$ y $Lh<1$, pruebe existencia y unicidad en $[t_0-h,t_0+h]$ entre las soluciones que permanecen en esa bola. Aplique el resultado a $u'=u^2$, $u(0)=1$, y explique por qué no puede concluirse existencia en todo $[0,2]$.
\end{ejercicio}
\begin{solucion}
Sea $J=[t_0-h,t_0+h]$ y
\[
 E=\{u\in C(J,\R^d):\norm{u-\xi}_\infty\le b\}.
\]
Es no vacío porque contiene la función constante $\xi$, y es cerrado en el Banach $C(J,\R^d)$, por continuidad de la norma. Por tanto es completo. El operador integral $S$ está definido en $E$, pues todas las parejas $(t,u(t))$ permanecen en el cilindro. Además
\[
 \norm{Su-\xi}_\infty\le Mh\le b,\qquad
 \norm{Su-Sv}_\infty\le Lh\norm{u-v}_\infty.
\]
Así $S:E\to E$ es contracción. Su único punto fijo proporciona la solución y cualquier otra solución que permanezca en la bola sería otro punto fijo de $S$.

Para $f(t,x)=x^2$, tome $a=b=1$, $\xi=1$. En $0\le x\le2$ tenemos $|f|\le4$ y $|x^2-y^2|\le4|x-y|$, de modo que $M=L=4$. Por ejemplo, $h=1/8$ cumple todas las condiciones. La función $u(t)=1/(1-t)$ resuelve la ecuación cerca de cero. De hecho, cualquier solución en un intervalo a la derecha de cero es positiva: su derivada es $u^2\ge0$ y $u(0)=1$. Entonces $(1/u)'=-1$, y por integración $1/u(t)=1-t$. Una solución en $[0,2]$ daría, en $t=1$, $1/u(1)=0$, imposible para un valor real finito positivo. La condición Lipschitz global del teorema~\ref{ampl:picard-teorema} falla: la derivada $2x$ no está acotada en $\R$.
\end{solucion}

\begin{ejercicio}[Ampliación: continuidad sin unicidad]\label{ej:ampl-picard-5}
Para el problema $u'=2\sqrt{|u|}$, $u(0)=0$, en $[0,1]$, pruebe que para cada $c\in[0,1]$,
\[
 u_c(t)=\begin{cases}0,&0\le t\le c,\\(t-c)^2,&c\le t\le1\end{cases}
\]
es una solución. Identifique exactamente qué hipótesis de unicidad falta, incluso si se restringe el estado a un intervalo compacto alrededor de cero.
\end{ejercicio}
\begin{solucion}
Las dos fórmulas coinciden en $t=c$. Sus derivadas son $0$ a la izquierda y $2(t-c)$ a la derecha, que también coinciden en $c$ y son continuas. Así $u_c$ es $C^1$, satisface el dato inicial y cumple $u_c'=2\sqrt{|u_c|}$ en ambos tramos y en su unión. Distintos tiempos $c<d$ producen funciones distintas: para $c<t<d$ se tiene $u_c(t)>0=u_d(t)$. Incluyendo $c=1$ obtenemos también la solución idénticamente cero.

La función $x\mapsto2\sqrt{|x|}$ es continua, pero no es Lipschitz en ninguna vecindad de cero, pues para $x>0$,
\[
 \frac{|2\sqrt{x}-0|}{|x-0|}=\frac2{\sqrt{x}}\longrightarrow\infty.
\]
Acotar el estado no repara esa falta local. El ejemplo muestra que continuidad del campo y existencia de soluciones no garantizan unicidad; en Picard el control uniforme de las diferencias respecto del estado tiene una función independiente de la continuidad.
\end{solucion}

\chapter{Equicontinuidad y teorema de Arzelà--Ascoli}

\begin{ejercicio}[Lectura 27, ejercicio 1]\label{ej:lectura-27-1}
Demuestre que $\mathcal A\subset C(X,\R)$ es compacta si y solo si es cerrada, acotada en norma y equicontinua.
\end{ejercicio}
\begin{solucion}
Si $\mathcal A$ es compacta, es cerrada y acotada por las propiedades generales de espacios métricos; además es equicontinua por Arzelà--Ascoli. Recíprocamente, la acotación en norma implica acotación puntual porque $|f(x)|\le\norm f_\infty$. La equicontinuidad y la acotación puntual hacen compacta su clausura. Como $\mathcal A$ es cerrada, coincide con esa clausura y es compacta. La condición de cerrado no puede omitirse si se desea compacidad y no solo compacidad relativa.
\end{solucion}

\begin{ejercicio}[Lectura 27, ejercicio 2]\label{ej:lectura-27-2}
Sea $v:C(X,\R)\times X\to\R$ la evaluación $v(f,x)=f(x)$.
\begin{enumerate}[label=(\roman*)]
\item Demuestre que $v$ es continua.
\item Para $\mathcal A\subset C(X,\R)$ y $\eps>0$, demuestre que
\[
 U_\eps=\{((f,x),(g,y)):|f(x)-g(y)|<\eps\}
\]
es abierto en $(\mathcal A\times X)^2$ y contiene la diagonal. Si $\mathcal A$ es compacta, deduzca la necesidad de Arzelà--Ascoli.
\end{enumerate}
\end{ejercicio}
\begin{solucion}
\begin{enumerate}[label=(\roman*)]
\item Fijemos $(f,x)$ y $\eps>0$. Por continuidad de $f$, existe $\delta>0$ con $|f(y)-f(x)|<\eps/2$ si $d(x,y)<\delta$. Si además $\norm{g-f}_\infty<\eps/2$, entonces
\[
 |v(g,y)-v(f,x)|\le|g(y)-f(y)|+|f(y)-f(x)|<\eps.
\]
Esto demuestra continuidad para la topología producto. No necesitamos una estimación de continuidad común para todas las funciones del espacio: trabajamos cerca de una función fija $f$.
\item La aplicación $((f,x),(g,y))\mapsto|v(f,x)-v(g,y)|$ es continua por el inciso anterior y por continuidad de la resta y del valor absoluto. Por tanto, $U_\eps$ es la preimagen del abierto $(-\infty,\eps)$ y es abierto. Si $(f,x)=(g,y)$, la diferencia vale cero; así la diagonal está contenida en $U_\eps$.

Supongamos ahora que $\mathcal A$ es compacta y no vacía. El producto $K=\mathcal A\times X$, con métrica
$\rho((f,x),(g,y))=\max\{\norm{f-g}_\infty,d(x,y)\}$, es compacto. Justificamos la propiedad de vecindad de la diagonal usada en el ejercicio: si no existiera $\delta>0$ tal que $\rho(p,q)<\delta$ implicara $(p,q)\in U_\eps$, podríamos elegir $p_n,q_n\in K$ con $\rho(p_n,q_n)<1/n$ y $(p_n,q_n)\notin U_\eps$. Una subsucesión de $p_n$ converge a algún $p\in K$, y la desigualdad triangular obliga a que la correspondiente subsucesión de $q_n$ converja al mismo $p$. El complemento de $U_\eps$ es cerrado, por lo que contendría a $(p,p)$, contradiciendo que $U_\eps$ incluye la diagonal.

Para $p=(f,x)$ y $q=(f,y)$ tenemos $\rho(p,q)=d(x,y)$. Así, el $\delta$ anterior asegura $|f(x)-f(y)|<\eps$ simultáneamente para toda $f\in\mathcal A$, lo que prueba equicontinuidad uniforme. Para cada $x$, la imagen de $\mathcal A$ por la evaluación es compacta y, por tanto, acotada. Si la familia inicial es solo relativamente compacta, aplicamos este argumento a su clausura y restringimos después a la familia.
\end{enumerate}
\end{solucion}

\begin{ejercicio}[Lectura 27, ejercicio 3]\label{ej:lectura-27-3}
En $C^1([a,b],\R)$ consideramos $\norm f_{C^1}=\norm f_\infty+\norm{f'}_\infty$. Demuestre que la inclusión $J:C^1([a,b])\to C([a,b])$, $Jf=f$, es lineal y compacta.
\end{ejercicio}
\begin{solucion}
La linealidad se verifica punto por punto y $\norm{Jf}_\infty\le\norm f_{C^1}$ demuestra continuidad. Sea $\mathcal E$ un conjunto acotado para la norma $C^1$; existe $M$ con $\norm f_\infty+\norm{f'}_\infty\le M$ para toda $f\in\mathcal E$. En particular, $|f(x)|\le M$. El teorema del valor medio da
$|f(x)-f(y)|\le M|x-y|$ para todos $x,y\in[a,b]$. Si $M>0$, elegir $\delta=\eps/M$ demuestra equicontinuidad; si $M=0$, todas las funciones son cero. Por Arzelà--Ascoli, $J(\mathcal E)$ tiene clausura compacta. Esto es la definición de operador compacto.

Como complemento, la norma usada convierte $C^1([a,b])$ en Banach. Si $(f_n)$ es de Cauchy para ella, existen límites uniformes $f_n\to f$ y $f_n'\to g$ en $C([a,b])$. El teorema~\ref{ch26:derivacion}, junto con continuidad en los extremos, implica $f'=g$ en el interior; las derivadas laterales en los extremos se obtienen pasando al límite en los cocientes o integrando $g$. Concretamente, $f_n(x)-f_n(a)=\int_a^x f_n'$ y el límite uniforme da $f(x)-f(a)=\int_a^xg$, de donde $f\in C^1([a,b])$. Las dos convergencias uniformes equivalen a $\norm{f_n-f}_{C^1}\to0$.
\end{solucion}

\begin{ejercicio}[Lectura 27, ejercicio 4]\label{ej:lectura-27-4}
Sea $\mathcal A\subset C([a,b],\R)$ una familia cuyas funciones son derivables en $(a,b)$. Si $|f'(x)|\le M$ para todas las funciones y todos los puntos interiores, demuestre que $\mathcal A$ es equicontinua.
\end{ejercicio}
\begin{solucion}
Para $x<y$ en $[a,b]$, la función $f$ es continua en $[x,y]$ y derivable en $(x,y)$, incluso si $x=a$ o $y=b$. El teorema del valor medio proporciona $\xi\in(x,y)$ tal que $f(y)-f(x)=f'(\xi)(y-x)$. Así $|f(y)-f(x)|\le M|y-x|$. Si $M>0$, el radio $\delta=\eps/(2M)$ funciona para toda la familia y todos los puntos. Si $M=0$, la misma estimación implica que cada función es constante, y cualquier $\delta>0$ funciona. Las constantes pueden variar sin estar acotadas entre funciones: equicontinuidad por sí sola no implica acotación puntual.
\end{solucion}

\begin{ejercicio}[Lectura 27, ejercicio 5: versión compleja]\label{ej:lectura-27-5}
Sea $\mathcal A\subset C(X,\C)$ y ponga $\mathcal A_1=\{\operatorname{Re}f:f\in\mathcal A\}$, $\mathcal A_2=\{\operatorname{Im}f:f\in\mathcal A\}$.
\begin{enumerate}[label=(\roman*)]
\item Pruebe que $\mathcal A$ es equicontinua si y solo si lo son $\mathcal A_1$ y $\mathcal A_2$.
\item Demuestre que $\Phi:C(X,\R)\times C(X,\R)\to C(X,\C)$, $\Phi(u,v)=u+iv$, es continua.
\item Deduzca Arzelà--Ascoli complejo: $\mathcal A$ es relativamente compacta si y solo si es equicontinua y cada $\mathcal A(x)$ es acotado.
\end{enumerate}
\end{ejercicio}
\begin{solucion}
\begin{enumerate}[label=(\roman*)]
\item Las desigualdades $|\operatorname{Re}z|,|\operatorname{Im}z|\le|z|$ demuestran que un radio válido para la familia compleja funciona para sus dos partes. Recíprocamente, dados $x$ y $\eps$, elegimos radios que controlen las diferencias de las partes real e imaginaria por $\eps/2$ y tomamos el mínimo. La desigualdad $|z|\le|\operatorname{Re}z|+|\operatorname{Im}z|$ produce el control $\eps$ para toda la familia compleja.
\item Para $u,v,u',v'$ reales continuas,
\[
 \norm{\Phi(u,v)-\Phi(u',v')}_\infty
 \le\norm{u-u'}_\infty+\norm{v-v'}_\infty.
\]
Así $\Phi$ es Lipschitz para la norma producto suma y, en particular, continua.
\item Supongamos equicontinuidad y acotación puntual de $\mathcal A$. Sus partes real e imaginaria tienen esas mismas propiedades, así que $K_j=\cl{\mathcal A_j}$ es compacto por el teorema real. El producto $K_1\times K_2$ es compacto, y su imagen por $\Phi$ es un compacto de $C(X,\C)$ que contiene a $\mathcal A$. Esa imagen es cerrada, por ser compacta en un espacio métrico. Por ello contiene también a $\cl{\mathcal A}$, que es un subconjunto cerrado de un compacto y resulta compacto.

Recíprocamente, las aplicaciones $f\mapsto\operatorname{Re}f$ y $f\mapsto\operatorname{Im}f$ son continuas y llevan la clausura compacta de $\mathcal A$ a conjuntos compactos reales. Las clausuras de $\mathcal A_1,\mathcal A_2$ son entonces compactas. El teorema real y el inciso (i) proporcionan equicontinuidad de $\mathcal A$. Para cada $x$, la evaluación compleja es continua y su imagen sobre la clausura es compacta en $\C$, luego acotada.
\end{enumerate}
\end{solucion}

\chapter{Extensión y aproximación de funciones continuas}

\begin{ejercicio}[Lectura 12, ejercicio 1]\label{ej:lectura-12-1}
Sin usar el lema de Urysohn, separe dos cerrados disjuntos $A,B\subset X$ por abiertos disjuntos que los contengan.
\end{ejercicio}
\begin{solucion}
Si alguno es vacío, elegimos para él el abierto vacío y para el otro el abierto $X$. Si ambos son no vacíos, tomamos
\[
 U=\{x:\dist(x,A)<\tfrac12\dist(x,B)\},\qquad
 W=\{x:\dist(x,B)<\tfrac12\dist(x,A)\}.
\]
Son abiertos porque son conjuntos de positividad o negatividad estricta de diferencias de funciones continuas. Si $a\in A$, entonces $\dist(a,A)=0$ y $\dist(a,B)>0$, por ser $B$ cerrado y disjunto de $A$; de ahí $a\in U$. El mismo argumento da $B\subset W$. Un punto en ambos abiertos satisfaría $\dist(x,A)<\frac14\dist(x,A)$, imposible para un número no negativo. Así $U\cap W=\varnothing$.
\end{solucion}

\begin{ejercicio}[Lectura 12, ejercicio 2]\label{ej:lectura-12-2}\label{ch12:soporte}
Defina el soporte de una función por
\[
 \supp(f)=\cl{\{x\in X:f(x)\ne0\}}.
\]
Sean $V$ cerrado y $U$ abierto con $V\subset U$.
\begin{enumerate}[label=(\roman*)]
\item Construya $\varphi:X\to[0,1]$ continua, igual a $1$ en $V$, y con $\supp(\varphi)\subset U$.
\item Examine la afirmación de que toda $g:V\to[a,b]$ continua tiene una extensión continua $h:X\to[a,b]$ con soporte en $U$, y resuelva la versión válida.
\end{enumerate}
\end{ejercicio}
\begin{solucion}
\textbf{(i)} Si $V=\varnothing$, tomamos $\varphi=0$; si $U=X$, tomamos $\varphi=1$. En otro caso aplicamos Urysohn a los cerrados disjuntos $V$ y $X\setminus U$ y obtenemos $u:X\to[0,1]$ con $u=0$ en $V$ y $u=1$ en $X\setminus U$. Definimos
\[
 \varphi(x)=\max\{0,1-3u(x)\}.
\]
Es continua, vale uno en $V$, y sus puntos no nulos están en $\{u<1/3\}$. Al tomar clausuras,
\[
 \supp(\varphi)\subset\{u\leq1/3\}\subset U,
\]
pues $u=1$ fuera de $U$. El margen entre $1/3$ y $1$ es lo que garantiza la inclusión para la \emph{clausura} del conjunto de no anulación, no sólo para ese conjunto.

\textbf{(ii)} El enunciado requiere una corrección cuando $0\notin[a,b]$ y $U\ne X$. Por ejemplo, si $X=\R$, $V=\{0\}$, $U=(-1,1)$ y $g(0)=1\in[1,2]$, toda función $h:\R\to[1,2]$ es no nula en todo $\R$ y su soporte es $\R$, imposible de contener en $U$. El texto original tampoco hace explícita la igualdad $h|_V=g$, que aquí incluimos para expresar el problema de extensión.

La versión válida conserva todas las propiedades si $0\in[a,b]$. Por Tietze, extendemos $g$ a $\widetilde g:X\to[a,b]$. Tomamos la función $\varphi$ del inciso (i) y ponemos $h=\varphi\widetilde g$. Es continua y coincide con $g$ en $V$. Como $0\leq\varphi\leq1$ y el intervalo contiene a cero, $h(x)$ es una combinación convexa de $\widetilde g(x)$ y cero y permanece en $[a,b]$. Además, $\{h\ne0\}\subset\{\varphi\ne0\}$ y, por clausura, $\supp(h)\subset\supp(\varphi)\subset U$.

Sin exigir $0\in[a,b]$, el mismo procedimiento da siempre una extensión con soporte en $U$ y valores en $[\min\{a,0\},\max\{b,0\}]$. Si se desea conservar exactamente el intervalo original, otra formulación válida fija un $c\in[a,b]$ y exige soporte de $h-c$ en $U$: la función $h=c+\varphi(\widetilde g-c)$ permanece en $[a,b]$, coincide con $g$ en $V$ y vale $c$ fuera del soporte de $\varphi$.
\end{solucion}

\begin{ejercicio}[Lectura 12, ejercicio 3]\label{ej:lectura-12-3}
Sean $U_1,\ldots,U_n$ abiertos que cubren $X$. Construya funciones continuas $\varphi_i:X\to[0,1]$ con $\supp(\varphi_i)\subset U_i$ y $\sum_{i=1}^n\varphi_i=1$.
\end{ejercicio}
\begin{solucion}
Si $X=\varnothing$, las únicas funciones con ese dominio satisfacen todas las condiciones. Supongamos $X\ne\varnothing$ y $n\geq1$. Para cada $i$ definimos
\[
 r_i(x)=\begin{cases}
 \min\{1,\dist(x,X\setminus U_i)\},&U_i\ne X,\\
 1,&U_i=X.
 \end{cases}
\]
Son continuas y no negativas. Si $U_i\ne X$, el complemento es cerrado no vacío y $r_i(x)>0$ exactamente cuando $x\in U_i$; esta afirmación también vale para $U_i=X$. Por la propiedad de cubierta, $S(x)=\sum_i r_i(x)>0$ en todo punto. Pongamos $w_i=r_i/S$. Son continuas, no negativas y suman uno.

No basta tomar directamente $w_i$ como partición subordinada en el sentido del soporte: su conjunto no nulo está en $U_i$, pero su clausura podría alcanzar la frontera. Introducimos un margen:
\[
 V_i=\{x:w_i(x)>1/(2n)\},\qquad
 g_i(x)=\max\{w_i(x)-1/(2n),0\}.
\]
Para cada $x$, algún $w_i(x)$ es al menos $1/n$, ya que la suma es uno. Por ello los $V_i$ cubren $X$ y $G=\sum_i g_i$ es estrictamente positiva. Además,
\[
 \cl{V_i}\subset\{w_i\geq1/(2n)\}\subset\{w_i>0\}=U_i.
\]
Ésta realiza explícitamente la reducción de la cubierta sugerida en la lectura original. Finalmente, definimos $\varphi_i=g_i/G$. Son continuas, están entre cero y uno y suman uno. Su conjunto de no anulación es $V_i$, luego $\supp(\varphi_i)=\cl{V_i}\subset U_i$. Hemos construido una partición finita de la unidad con todas las inclusiones de soporte justificadas.
\end{solucion}

\begin{ejercicio}[Lectura 28, ejercicio 1: aproximación de Weierstrass]\label{ej:lectura-28-1}
Sea $f:[a,b]\to\F$ continua. Demuestre que existe una sucesión de polinomios con coeficientes en $\F$ que converge uniformemente a $f$.
\end{ejercicio}
\begin{solucion}
Si $a=b$, basta el polinomio constante $f(a)$. Supongamos $a<b$ y primero $\F=\R$. Las restricciones de polinomios reales a $[a,b]$ forman una subálgebra de $C([a,b],\R)$. La función constante $1$ garantiza que no se anula en ningún punto, y el polinomio $p(x)=x$ separa puntos. Por el teorema~\ref{ch28:stone}, para cada $n$ existe un polinomio $p_n$ con $\norm{f-p_n}_\infty<1/n$. Esa elección da la sucesión requerida.

Si $f$ es compleja, escribimos $f=u+iv$ con $u,v$ reales continuas. Elegimos polinomios reales $p_n,q_n$ con $\norm{u-p_n}_\infty<1/(2n)$ y $\norm{v-q_n}_\infty<1/(2n)$. El polinomio complejo $r_n=p_n+iq_n$ satisface
$\norm{f-r_n}_\infty\le\norm{u-p_n}_\infty+\norm{v-q_n}_\infty<1/n$.
\end{solucion}

\begin{ejercicio}[Lectura 28, ejercicio 2: determinación por momentos]\label{ej:lectura-28-2}
Sea $f:[0,1]\to\R$ continua y suponga que $\int_0^1 f(x)x^n\,dx=0$ para todo entero $n\ge0$. Demuestre que $f=0$.
\end{ejercicio}
\begin{solucion}
Por linealidad de la integral, $\int_0^1fp=0$ para cualquier polinomio real $p$, porque $p$ es una combinación finita de monomios. El ejercicio~\ref{ej:lectura-28-1} proporciona polinomios $p_n$ con $\norm{p_n-f}_\infty\to0$. La función $f$ es acotada por continuidad en un compacto, y
\[
 \left|\int_0^1f(x)p_n(x)\,dx-\int_0^1f(x)^2\,dx\right|
 \le\int_0^1|f(x)|\,|p_n(x)-f(x)|\,dx
 \le\norm f_\infty\norm{p_n-f}_\infty\longrightarrow0.
\]
La primera integral es cero para todo $n$, de modo que $\int_0^1f^2=0$.

Si existiera $x_0$ con $f(x_0)\ne0$, por continuidad habría una vecindad relativa de $x_0$ en $[0,1]$ donde $|f(x)|>|f(x_0)|/2$. Esa vecindad contiene un intervalo de longitud positiva, incluso cuando $x_0$ es un extremo. La integral de $f^2$ sobre ese intervalo sería al menos su longitud multiplicada por $|f(x_0)|^2/4>0$, contradiciendo $\int_0^1f^2=0$. Por tanto, $f$ se anula en cada punto.
\end{solucion}

\begin{ejercicio}[Lectura 28, ejercicio 3: Stone--Weierstrass complejo]\label{ej:lectura-28-3}
Sea $X$ un espacio métrico compacto. Una subálgebra compleja $A\subset C(X,\C)$ es \emph{autoadjunta} si $\overline f\in A$ para toda $f\in A$. Demuestre que, si $A$ es autoadjunta, separa puntos y no se anula en ningún punto, entonces es densa en $C(X,\C)$.
\end{ejercicio}
\begin{solucion}
\index{subalgebra@subálgebra!autoadjunta}
Consideremos $A_\R=\{f\in A:f(X)\subset\R\}$. Es una subálgebra \emph{real} de $C(X,\R)$: las combinaciones lineales reales y los productos conservan tanto la pertenencia a $A$ como los valores reales. Para $f\in A$, la autoadjunción y la estructura de espacio vectorial complejo dan
\[
 \operatorname{Re}f=\frac{f+\overline f}{2}\in A_\R,
 \qquad\operatorname{Im}f=\frac{f-\overline f}{2i}\in A_\R.
\]
Si $x\ne y$, elegimos $f\in A$ con $f(x)\ne f(y)$. Los números complejos no pueden tener simultáneamente iguales partes real e imaginaria si son distintos; por tanto, $\operatorname{Re}f$ o $\operatorname{Im}f$ separa esos puntos. Si para un $x$ elegimos $f\in A$ con $f(x)\ne0$, alguna de sus dos partes es no nula en $x$. Así, $A_\R$ satisface las hipótesis del teorema real y es densa en $C(X,\R)$.

Para $h=u+iv\in C(X,\C)$ y $\eps>0$, elegimos $a,b\in A_\R$ con $\norm{u-a}_\infty<\eps/2$ y $\norm{v-b}_\infty<\eps/2$. Entonces $a+ib\in A$ y
\[
 \norm{h-(a+ib)}_\infty\le\norm{u-a}_\infty+\norm{v-b}_\infty<\eps.
\]
Esto prueba la densidad compleja.
\end{solucion}

\begin{ejercicio}[Lectura 28, ejercicio 4: por qué hace falta conjugación]\label{ej:lectura-28-4}
Sea $\mathbb T=\{z\in\C:|z|=1\}$ y sea $A$ el álgebra de restricciones a $\mathbb T$ de los polinomios complejos en $z$:
\[
 p(e^{i\theta})=\sum_{k=0}^N a_ke^{ik\theta}.
\]
Demuestre que $A$ separa puntos y no se anula en ningún punto, pero no es densa en $C(\mathbb T,\C)$.
\end{ejercicio}
\begin{solucion}
La función constante $1$ pertenece a $A$, así que no hay ceros comunes. La función $p(z)=z$ pertenece a $A$ y toma valores distintos en puntos distintos de $\mathbb T$; por tanto, separa puntos. La suma y el producto de dos polinomios en potencias no negativas de $z$ siguen siendo de ese tipo, de modo que $A$ es una subálgebra.

Definamos el funcional
\[
 L(h)=\frac1{2\pi}\int_0^{2\pi}h(e^{i\theta})e^{i\theta}\,d\theta
 \qquad(h\in C(\mathbb T,\C)).
\]
Es lineal y satisface $|L(h)|\le\norm h_\infty$, por la desigualdad para integrales del capítulo~\ref{chap:26}; por ello es continuo. Para un monomio $h(z)=z^k$, con $k\ge0$, calculamos
\[
 L(z^k)=\frac1{2\pi}\int_0^{2\pi}e^{i(k+1)\theta}\,d\theta
 =\frac{e^{2\pi i(k+1)}-1}{2\pi i(k+1)}=0.
\]
Por linealidad, $L(p)=0$ para todo $p\in A$. La continuidad implica que también se anula en $\cl A$: si $p_n\to h$ uniformemente, $L(h)=\lim_nL(p_n)=0$.

Sin embargo, la función continua $h(z)=\overline z$ satisface
\[
 L(h)=\frac1{2\pi}\int_0^{2\pi}e^{-i\theta}e^{i\theta}\,d\theta=1.
\]
Así $h\notin\cl A$. Incluso obtenemos una cota cuantitativa: para todo $p\in A$,
$1=|L(h-p)|\le\norm{h-p}_\infty$. Como el polinomio cero da error uno, la distancia de $h$ a $A$ es exactamente uno. Esta álgebra no es autoadjunta, pues contiene $z$ y no contiene $\overline z$.
\end{solucion}

\chapter{Producto interno y proyección ortogonal}

\begin{ejercicio}[Lectura 29, ejercicio 1]\label{ej:lectura-29-1}
Caracterizar la igualdad en Cauchy--Schwarz.
\end{ejercicio}
\begin{solucion}
Si uno de los vectores es cero, hay igualdad y el conjunto es linealmente dependiente. Si $y\ne0$, la prueba del teorema \ref{ch29:cs} muestra exactamente
\[
 \Big\|x-\frac{\ip{x}{y}}{\norm{y}^{2}}y\Big\|^{2}
 =\norm{x}^{2}-\frac{|\ip{x}{y}|^{2}}{\norm{y}^{2}}.
\]
La igualdad equivale a que el miembro izquierdo sea cero, es decir, a que $x$ sea múltiplo de $y$. Recíprocamente, si $x=\alpha y$, ambos miembros de Cauchy--Schwarz son $|\alpha|\norm{y}^{2}$.
\end{solucion}

\begin{ejercicio}[Lectura 29, ejercicio 2: corrección del recíproco]\label{ej:lectura-29-2}
Analizar si $\norm{v_1+\cdots+v_n}^{2}=\sum_j\norm{v_j}^{2}$ implica ortogonalidad por pares en un espacio real.
\end{ejercicio}
\begin{solucion}
Para dos vectores sí: \eqref{ch29:expansion} implica $2\ip{v_1}{v_2}=0$. Para tres o más es falso. En $\R$, con producto usual, tómense $v_1=1$, $v_2=1$, $v_3=-1/2$. La norma al cuadrado de su suma es $9/4$, igual a $1+1+1/4$, pero ninguno de los productos por pares es cero. En general la igualdad solamente afirma $\sum_{i<j}\ip{v_i}{v_j}=0$; distintos términos pueden cancelarse. Una versión correcta exige la igualdad de Pitágoras para \emph{cada par}. Otra exige $\norm{\sum_j t_jv_j}^{2}=\sum_j t_j^{2}\norm{v_j}^{2}$ para todos los escalares reales $t_j$: basta elegir dos coeficientes iguales a uno y los demás cero. Sobre $\C$, incluso la igualdad para un solo par solo fuerza $\operatorname{Re}\ip{x}{y}=0$; se necesita además aplicarla al par $x,iy$ para anular la parte imaginaria.
\end{solucion}

\begin{ejercicio}[Lectura 29, ejercicio 3: identidad del paralelogramo]\label{ej:lectura-29-3}
Demostrar que una norma procede de un producto interno si y solo si
\[
 \norm{x+y}^{2}+\norm{x-y}^{2}=2\norm{x}^{2}+2\norm{y}^{2}.
\]
Aplicarlo a (i) $\ell^{1}$ y (ii) la norma uniforme en $B(S,W)$.
\end{ejercicio}
\begin{solucion}
Los factores $2$ son necesarios: con $y=0$ se detecta inmediatamente su omisión en el enunciado original. La necesidad ya se probó. Para la suficiencia escribamos $q(x)=\norm{x}^{2}$ y, considerando primero el espacio vectorial real subyacente,
\[
 B(x,y)=\frac{q(x+y)-q(x-y)}4.
\]
Es simétrica porque $q(-z)=q(z)$, continua y satisface $B(x,x)=q(x)$. El paralelogramo, aplicado a $(x+y,z)$ y $(x-y,z)$ y restando, da
\[
 B(x+z,y)+B(x-z,y)=2B(x,y).
\]
Con $z=x$ obtenemos $B(2x,y)=2B(x,y)$. Sustituyendo después $x=(u+v)/2$, $z=(u-v)/2$, se deduce $B(u+v,y)=B(u,y)+B(v,y)$. La aditividad implica homogeneidad para enteros y racionales; por continuidad y aproximación racional, también para todos los reales. La simetría da bilinealidad. Así $B$ es el producto interno buscado sobre $\R$.

Sobre $\C$, $q(ix)=q(x)$ implica $B(ix,iy)=B(x,y)$ y, sustituyendo $y$ por $-iy$, $B(ix,y)=-B(x,iy)$. Definamos
\[
 \ip{x}{y}=B(x,y)+iB(x,iy).
\]
La bilinealidad real y las dos identidades anteriores dan $\ip{ix}{y}=i\ip{x}{y}$ y, por ello, linealidad compleja en la primera entrada. Además, $B(y,ix)=B(ix,y)=-B(x,iy)$, lo que prueba simetría conjugada. Finalmente $B(x,ix)=-B(ix,x)$ y la simetría obligan a $B(x,ix)=0$; así $\ip{x}{x}=q(x)$.

\emph{(i)} En $\ell^{1}$, para $x=e_1$ y $y=e_2$, el miembro izquierdo del paralelogramo vale $2^2+2^2=8$, mientras el derecho vale $4$.

\emph{(ii)} Se necesita una salvedad: si $S$ contiene al menos dos puntos y $W\ne\{0\}$, elegimos $w\in W$ de norma uno y funciones $f,g$ que valen $w$ en dos puntos distintos, respectivamente, y cero en los restantes. Entonces $\norm{f\pm g}_\infty=1$ y $\norm{f}_\infty=\norm{g}_\infty=1$; los dos miembros valen $2$ y $4$. Si $S$ es un solo punto, $B(S,W)$ es isométrico a $W$, y la afirmación general sería falsa cuando la norma de $W$ sí procede de un producto interno. Si $W=\{0\}$, o si $S$ es vacío y se identifica el espacio de funciones con el espacio cero, la norma también es hilbertiana.
\end{solucion}

\begin{ejercicio}[Lectura 29, ejercicio 4]\label{ej:lectura-29-4}
Demostrar la continuidad conjunta del producto interno.
\end{ejercicio}
\begin{solucion}
Si $x_n\to x$ y $y_n\to y$, la sucesión $(x_n)$ está acotada, digamos por $C$. Entonces
\[
 |\ip{x_n}{y_n}-\ip{x}{y}|
 \leq |\ip{x_n}{y_n-y}|+|\ip{x_n-x}{y}|
 \leq C\norm{y_n-y}+\norm{x_n-x}\norm{y}\longrightarrow0.
\]
Como el dominio es métrico, este criterio secuencial demuestra continuidad. La estimación también proporciona directamente la prueba con $\eps$ y $\delta$ al restringir $\norm{x'-x}<1$, pues entonces $\norm{x'}\leq\norm{x}+1$.
\end{solucion}

\begin{ejercicio}[Lectura 29, ejercicio 5: identidad polar]\label{ej:lectura-29-5}
\index{polarizacion@polarización, identidad de}
Si $\zeta$ es una raíz primitiva de la unidad de orden $n\geq3$, probar
\[
 \ip{x}{y}=\frac1n\sum_{k=0}^{n-1}\zeta^k\norm{x+\zeta^ky}^{2}.
\]
\end{ejercicio}
\begin{solucion}
La expansión del sumando es
\[
 \zeta^k\norm{x+\zeta^ky}^{2}
 =\zeta^k(\norm{x}^{2}+\norm{y}^{2})+\ip{x}{y}+\zeta^{2k}\ip{y}{x}.
\]
Las sumas de $\zeta^k$ y de $\zeta^{2k}$ son cero: ambas son sumas geométricas, sus razones son distintas de uno y la potencia $n$-ésima de cada razón es uno. Para la segunda razón se utiliza $n\geq3$. Solo permanecen $n$ copias de $\ip{x}{y}$.
\end{solucion}

\begin{ejercicio}[Lectura 29, ejercicio 6]\label{ej:lectura-29-6}
Sea $h:[a,b]\to\F$ continua y sin ceros, con $a<b$. Determinar cuándo $\ip{f}{g}_h=\int_a^b f\bar g h$ es producto interno, y discutir el caso continuo a trozos.
\end{ejercicio}
\begin{solucion}
Si $h>0$, la sesquilinealidad y la simetría conjugada se verifican integrando; la positividad definida se obtiene en un intervalo donde $|f|$ y $h$ tienen cotas inferiores positivas.

Recíprocamente, la realidad de $\ip{f}{f}_h$ implica $\int|f|^{2}\operatorname{Im}h=0$ para toda $f$ continua. Si $\operatorname{Im}h$ fuera positiva en un punto, lo sería con una cota positiva en un intervalo relativo; una función triangular no nula soportada allí daría integral positiva, contradicción. El mismo argumento aplicado a $-\operatorname{Im}h$ descarta valores negativos. Así $h$ es real. Si $h$ fuera negativa en algún punto, la misma función triangular daría $\ip{f}{f}_h<0$. Por tanto $h\geq0$; al no tener ceros, $h>0$.

Para una función acotada continua a trozos, sus valores en los puntos de discontinuidad no afectan a la integral. Por ejemplo, $h=1$ salvo $h(c)=-1$ en un punto interior produce el producto interno usual aunque no sea positiva en todo punto. La caracterización general es: $h$ debe ser real y no negativa casi en todas partes, y $\int_J h>0$ para todo intervalo relativo $J$ de longitud positiva. La necesidad de la primera condición se prueba con funciones triangulares alrededor de puntos de continuidad; los posibles puntos excepcionales son finitos. Si $\int_Jh=0$, una función continua no nula soportada en el interior de $J$ tiene norma cero. Para la suficiencia, si $f\ne0$, existe un intervalo $J$ y $c>0$ con $|f|\geq c$ en $J$, luego $\int|f|^{2}h\geq c^{2}\int_Jh>0$. Aquí «casi en todas partes» puede leerse como «excepto en un conjunto de longitud cero»; en los ejemplos basta un conjunto finito.
\end{solucion}

\begin{ejercicio}[Lectura 29, ejercicio 7]\label{ej:lectura-29-7}
\index{ell2@\ensuremath{\ell^2}!ponderado}
Para $w_n>0$, sea $\ell^{2}_w=\{x:\sum_n w_n|x_n|^{2}<\infty\}$. Probar (i) que es espacio vectorial, (ii) que $\ip{x}{y}_w=\sum_nw_nx_n\bar y_n$ es producto interno y (iii) que es completo.
\end{ejercicio}
\begin{solucion}
\emph{(i)} La desigualdad $|\alpha x_n+\beta y_n|^{2}\leq2|\alpha|^{2}|x_n|^{2}+2|\beta|^{2}|y_n|^{2}$ prueba la estabilidad al sumar y multiplicar por escalares, después de multiplicar por $w_n$ y sumar.
\emph{(ii)} Definamos $U(x)_n=\sqrt{w_n}x_n$. Es una biyección lineal de $\ell^{2}_w$ a $\ell^{2}$, con inversa $U^{-1}(z)_n=z_n/\sqrt{w_n}$. La serie del producto converge absolutamente por Cauchy--Schwarz en $\ell^{2}$ y es precisamente $\ip{Ux}{Uy}_{\ell^{2}}$, por lo que hereda todos los axiomas.
\emph{(iii)} Si $(x^{(j)})$ es de Cauchy en la norma ponderada, $(Ux^{(j)})$ es de Cauchy en $\ell^{2}$, converge a algún $z$, y $x^{(j)}\to U^{-1}z$. No se requiere que los pesos estén acotados ni separados de cero: la transformación $U$ incorpora exactamente esos pesos.
\end{solucion}

\begin{ejercicio}[Lectura 29, ejercicio 8]\label{ej:lectura-29-8}
\index{ell2@\ensuremath{\ell^2}!indices arbitrarios@índices arbitrarios}
Para cualquier conjunto $S$, definir $\ell^{2}(S)$ por $\sum_{s\in S}|f(s)|^{2}<\infty$, con la suma de la definición \ref{ch29:suma}. Probar (i) la estructura vectorial, (ii) la existencia del producto $\ip{f}{g}=\sum_{s\in S}f(s)\overline{g(s)}$ y (iii) la completitud.
\end{ejercicio}
\begin{solucion}
Los conjuntos finitos en la definición deben tener índices distintos: permitir repeticiones haría que cualquier valor no nulo pudiera sumarse arbitrariamente muchas veces.

\emph{(i)} Para cada $F\subset S$ finito,
\[
 \sum_F|\alpha f+\beta g|^{2}
 \leq2|\alpha|^{2}\norm{f}_2^{2}+2|\beta|^{2}\norm{g}_2^{2}.
\]
Tomar el supremo prueba que $\alpha f+\beta g\in\ell^{2}(S)$.

\emph{(ii)} El lema \ref{ch29:soporte} muestra que los soportes de $f$ y $g$ son numerables. Sobre su unión, Cauchy--Schwarz finita da $\sum_F|f(s)g(s)|\leq\norm{f}_2\norm{g}_2$. La serie es absolutamente convergente, independiente de la enumeración. Su linealidad y simetría conjugada se deducen de esas propiedades para series absolutamente convergentes; su valor en $(f,f)$ es la suma de los cuadrados. Si vale cero, cada $|f(s)|^{2}$ es cero.

\emph{(iii)} Sea $(f_j)$ de Cauchy. Para cada $s$, $|f_j(s)-f_k(s)|\leq\norm{f_j-f_k}_2$, así que existe $f(s)=\lim_j f_j(s)$. La sucesión está acotada por algún $C$. Para todo $F$ finito, al pasar al límite en una suma finita,
\[
 \sum_{s\in F}|f(s)|^{2}=\lim_j\sum_{s\in F}|f_j(s)|^{2}\leq C^{2}.
\]
Por tanto $f\in\ell^{2}(S)$. Dado $\eps>0$, elija $N$ tal que $\norm{f_j-f_k}_2<\eps$ para $j,k\geq N$. Fijando $j\geq N$ y un $F$ finito, y haciendo $k\to\infty$, resulta
\[
 \sum_{s\in F}|f_j(s)-f(s)|^{2}\leq\eps^{2}.
\]
Tomar el supremo sobre $F$ demuestra $\norm{f_j-f}_2\leq\eps$. Esto prueba completitud sin suponer que $S$ sea numerable.
\end{solucion}

\begin{ejercicio}[Ampliación: aproximar una parábola por una recta]\label{ej:ampl-gram-1}
En $C([0,1],\R)$ con $\ip fg=\int_0^1f(t)g(t)\,dt$, encuentre el polinomio $a+bt$ que minimiza $\int_0^1(t^2-a-bt)^2\,dt$. Calcule el error mínimo y verifique directamente la ortogonalidad del residuo.
\end{ejercicio}
\begin{solucion}
Para $a_1=1$, $a_2=t$ y $u=t^2$, los momentos $\int_0^1t^k\,dt=1/(k+1)$ dan
\[
 G=\begin{pmatrix}1&1/2\\1/2&1/3\end{pmatrix},\qquad
 \alpha=\begin{pmatrix}1/3\\1/4\end{pmatrix}.
\]
El sistema es $a+b/2=1/3$, $a/2+b/3=1/4$. Restar la mitad de la primera ecuación a la segunda da $b/12=1/12$, luego $b=1$ y $a=-1/6$. Así $p(t)=t-1/6$. El residuo $r=t^2-t+1/6$ satisface
\[
 \int_0^1r=\frac13-\frac12+\frac16=0,\qquad
 \int_0^1tr=\frac14-\frac13+\frac1{12}=0.
\]
Es perpendicular a $1$ y a $t$, y por linealidad a cualquier recta. El error cuadrático mínimo es
\[
 \norm{u-p}^2=\frac15-\alpha^*\begin{pmatrix}-1/6\\1\end{pmatrix}
 =\frac15-\left(-\frac1{18}+\frac14\right)=\frac1{180}.
\]
La distancia, sin elevar al cuadrado, es $1/\sqrt{180}$. Aunque $C([0,1])$ no es completo para esta norma integral, la proyección sobre un subespacio de dimensión finita existe.
\end{solucion}

\begin{ejercicio}[Ampliación: dónde va el conjugado]\label{ej:ampl-gram-2}
En $\C^2$, con $\ip xy=x_1\overline{y_1}+x_2\overline{y_2}$, proyecte $u=(1,1)$ sobre la recta compleja generada por $a=(1,i)$. Compruebe la norma del error y explique por qué usar $\ip au/\norm a^2$ como coeficiente daría una respuesta incorrecta.
\end{ejercicio}
\begin{solucion}
La matriz de Gram tiene una sola entrada $G=\norm a^2=2$. El dato es $\alpha=\ip ua=1-i$, luego
\[
 c=\frac{1-i}{2},\qquad p=ca=\left(\frac{1-i}{2},\frac{1+i}{2}\right).
\]
El residuo es $r=((1+i)/2,(1-i)/2)$ y
\[
 \ip ra=\frac{1+i}{2}-i\frac{1-i}{2}=0,\qquad
 \norm r^2=\frac{|1+i|^2+|1-i|^2}{4}=1.
\]
Si se usa $c'=(1+i)/2$, entonces
$\ip{u-c'a}{a}=(1-i)-2c'=-2i\ne0$. El problema no es una elección estética de notación: con producto interno lineal en la primera entrada, el coeficiente que se extrae de $\ip{ca}{a}$ es $c$, y por eso la ecuación correcta es $c\norm a^2=\ip ua$.
\end{solucion}

\begin{ejercicio}[Ampliación: tres datos y dos parámetros]\label{ej:ampl-gram-3}
Ajuste una recta $y=a+bt$ a los datos $(0,1)$, $(1,2)$ y $(2,2)$ minimizando la suma de los cuadrados de los tres errores. Calcule los parámetros, los valores ajustados y el error mínimo. Explique por qué el problema tiene una solución única aunque ninguna recta pase por los tres datos.
\end{ejercicio}
\begin{solucion}
Las columnas $a_1=(1,1,1)$ y $a_2=(0,1,2)$ son independientes. Con $u=(1,2,2)$,
\[
 G=\begin{pmatrix}3&3\\3&5\end{pmatrix},\qquad
 \alpha=\begin{pmatrix}5\\6\end{pmatrix}.
\]
Restar las ecuaciones $3a+3b=5$ y $3a+5b=6$ da $b=1/2$, y después $a=7/6$. Los valores ajustados son
\[
 p=\left(\frac76,\frac53,\frac{13}6\right),\qquad
 r=u-p=\left(-\frac16,\frac13,-\frac16\right).
\]
El residuo tiene suma cero y producto con $(0,1,2)$ igual a $1/3-2/6=0$. Su norma cuadrada es $1/36+1/9+1/36=1/6$. Para cualquier otro ajuste $v$ en el mismo plano, Pitágoras da $\norm{u-v}^2=1/6+\norm{p-v}^2$, estrictamente mayor salvo si $v=p$. La independencia de las columnas asegura que $p$ corresponde a un único par $(a,b)$. Los datos no son colineales porque las pendientes entre puntos consecutivos son $1$ y $0$; la proyección responde al problema de aproximación aun cuando el de interpolación sea imposible.
\end{solucion}

\begin{ejercicio}[Ampliación: cambiar de base sin cambiar la proyección]\label{ej:ampl-gram-4}
Sea $a_1,\ldots,a_m$ una lista independiente y $C=(C_{ij})$ una matriz invertible. Defina $b_j=\sum_i a_iC_{ij}$. Pruebe que las matrices de Gram y los datos de proyección satisfacen $G_b=C^*G_aC$ y $\alpha_b=C^*\alpha_a$, y relacione las coordenadas del mismo vector proyectado. Ilustre el resultado con $a_1=(1,0)$, $a_2=(1,\eps)$, $\eps>0$, y $u=(0,1)$ en $\R^2$.
\end{ejercicio}
\begin{solucion}
Por linealidad en la primera entrada y conjugación en la segunda,
\[
 \ip{b_j}{b_i}=\sum_{k,l}C_{kj}\overline{C_{li}}\ip{a_k}{a_l}
 =(C^*G_aC)_{ij}.
\]
Asimismo $\ip u{b_i}=\sum_l\overline{C_{li}}\ip u{a_l}=(C^*\alpha_a)_i$. Si $c$ resuelve $G_ac=\alpha_a$ y $d=C^{-1}c$, entonces
\[
 G_bd=C^*G_aCC^{-1}c=C^*\alpha_a=\alpha_b.
\]
Además $\sum_jd_jb_j=\sum_i(Cd)_ia_i=\sum_ic_ia_i$. Las coordenadas cambian, pero el vector proyectado es el mismo.

En el ejemplo,
\[
 G=\begin{pmatrix}1&1\\1&1+\eps^2\end{pmatrix},\qquad
 \alpha=\begin{pmatrix}0\\\eps\end{pmatrix},\qquad
 c=\begin{pmatrix}-1/\eps\\1/\eps\end{pmatrix}.
\]
La proyección es todo $u$, pues las dos columnas generan $\R^2$. Sin embargo, sus coordenadas crecen cuando $\eps\downarrow0$ y se cancelan en la primera componente. Tomando la base $b_1=a_1$, $b_2=(a_2-a_1)/\eps=(0,1)$, obtenemos $G_b=I$ y coordenadas $(0,1)$. El fenómeno distingue la geometría del subespacio de la calidad de las coordenadas elegidas para describirlo.
\end{solucion}

\begin{ejercicio}[Ampliación: listas dependientes y coeficientes no únicos]\label{ej:ampl-gram-5}
Permita ahora que $a_1,\ldots,a_m$ sean dependientes. Defina $A:\F^m\to V$ por $Ac=\sum_jc_ja_j$. Pruebe que $\ker G=\ker A$, que $\operatorname{rango}G=\dim\spanop\{a_j\}$ y que $Gc=\alpha$ sigue teniendo soluciones para $\alpha_i=\ip u{a_i}$. Muestre que todas producen el mismo mejor aproximante. Calcule el sistema para $a_1=1$, $a_2=t$, $a_3=1+t$ y $u=t^2$, con producto integral en $[0,1]$.
\end{ejercicio}
\begin{solucion}
Si $Ac=0$, entonces $(Gc)_i=\ip{Ac}{a_i}=0$ para cada $i$. Recíprocamente, si $Gc=0$, la identidad $\norm{Ac}^2=c^*Gc$ implica $Ac=0$. Por el teorema de rango y nulidad en el dominio finito $\F^m$,
\[
 \operatorname{rango}G=m-\dim\ker G
 =m-\dim\ker A=\dim\operatorname{im}A.
\]
Para una matriz hermitiana, $\operatorname{im}G=(\ker G)^\perp$: la inclusión se obtiene de $z^*Gv=(Gz)^*v=0$ si $Gz=0$, y ambos subespacios tienen dimensión $m-\dim\ker G$. Si $z\in\ker G=\ker A$, entonces
\[
 z^*\alpha=\sum_i\overline{z_i}\ip u{a_i}=\ip u{Az}=0.
\]
Por tanto $\alpha\in(\ker G)^\perp=\operatorname{im}G$, y el sistema tiene solución. Dos soluciones difieren en $\ker G=\ker A$, así que dan el mismo vector $Ac$. El cálculo del teorema~\ref{ampl:gram-proyeccion} prueba que $u-Ac$ es perpendicular a cada $a_i$ y a su espacio generado; por ello ese vector es el mejor aproximante único.

En el ejemplo,
\[
 G=\begin{pmatrix}
 1&1/2&3/2\\1/2&1/3&5/6\\3/2&5/6&7/3
 \end{pmatrix},\qquad
 \alpha=\begin{pmatrix}1/3\\1/4\\7/12\end{pmatrix}.
\]
El núcleo de $A$ está generado por $(-1,-1,1)$, pues $c_1+c_2t+c_3(1+t)=0$ equivale a $c_1+c_3=c_2+c_3=0$. Una solución es $(-1/6,1,0)$ por el ejercicio~\ref{ej:ampl-gram-1}. Todas son
\[
 c=(-1/6-s,\,1-s,\,s),\qquad s\in\R,
\]
y todas representan $t-1/6$. La dependencia hace singular a Gram y elimina la unicidad de las coordenadas, pero no la unicidad del vector proyectado.
\end{solucion}

\begin{ejercicio}[Lectura 30, ejercicio 1]\label{ej:lectura-30-1}
Sea $C$ un subconjunto convexo, completo y no vacío de un espacio con producto interno. Probar que para cada $x$ existe un único $c\in C$ tal que $\norm{x-c}=\dist(x,C)$.
\end{ejercicio}
\begin{solucion}
Escribamos $d=\dist(x,C)$ y elijamos $c_n\in C$ con $\norm{x-c_n}\to d$. La convexidad asegura $(c_n+c_k)/2\in C$. La misma identidad del paralelogramo de la prueba del teorema \ref{ch30:proyeccion} da
\[
 \norm{c_n-c_k}^{2}\leq2\norm{x-c_n}^{2}+2\norm{x-c_k}^{2}-4d^{2}\to0.
\]
Como $C$ es completo, $c_n\to c\in C$; por continuidad, $\norm{x-c}=d$. Si $c'$ fuera otro minimizador, el paralelogramo daría
\[
 \norm{c-c'}^{2}=4d^{2}-4\left\|x-\frac{c+c'}2\right\|^{2}\leq0.
\]
Así $c=c'$. Esta prueba utiliza solamente que los puntos medios permanecen en $C$; la completitud garantiza que el límite también permanece.
\end{solucion}

\begin{ejercicio}[Lectura 30, ejercicio 2: Parseval sobre un subespacio]\label{ej:lectura-30-2}
Sea $M$ ortogonalmente complementado y $E\subset M$ ortonormal con $\spanop E$ denso en $M$. Para $x\in V$ y $m=P_Mx$, probar
\[
 \norm{m}^{2}=\sum_{e\in E}|\ip{x}{e}|^{2}.
\]
\end{ejercicio}
\begin{solucion}
Como $x-m\perp M$, $\ip{x}{e}=\ip{m}{e}$. Para $F\subset E$ finito, pongamos $p_F=\sum_{e\in F}\ip{m}{e}e$. Entonces $m-p_F\perp\spanop F$ y
\[
 \norm{m-p_F}^{2}=\norm{m}^{2}-\sum_{e\in F}|\ip{m}{e}|^{2}.
\]
Dado $\eps>0$, la densidad proporciona $q\in\spanop E$ con $\norm{m-q}<\eps$. El vector $q$ es una combinación de un conjunto finito $F_0$. Para todo $F\supset F_0$, $p_F$ es el mejor aproximante en $\spanop F$, por lo que $\norm{m-p_F}\leq\norm{m-q}<\eps$. Así las sumas finitas se acercan a $\norm{m}^{2}$ desde abajo. Su supremo es exactamente ese valor. La prueba no exige que $V$ o $M$ sean completos.
\end{solucion}

\begin{ejercicio}[Lectura 30, ejercicio 3: tres perpendiculares]\label{ej:lectura-30-3}
Sean $M\subset N$ subespacios ortogonalmente complementados. Probar $P_MP_Nx=P_Mx$.
\end{ejercicio}
\begin{solucion}
Pongamos $m=P_Mx$ y $n=P_Nx$. Como $x-n\in N^{\perp}\subset M^{\perp}$ y $x-m\in M^{\perp}$, su diferencia $n-m=(x-m)-(x-n)$ pertenece a $M^{\perp}$. Puesto que $m\in M$, la caracterización del mejor aproximante dice que $m=P_Mn$. También $P_NP_Mx=P_Mx$, porque $P_Mx\in M\subset N$. Se obtiene así $P_MP_N=P_NP_M=P_M$.
\end{solucion}

\begin{ejercicio}[Lectura 30, ejercicio 4: complementos ortogonales]\label{ej:lectura-30-4}
Para subespacios $M,N$, estudiar
\begin{enumerate}[label=(\roman*)]
 \item $(M+N)^{\perp}=M^{\perp}\cap N^{\perp}$;
 \item $(M\cap N)^{\perp}=\cl{M^{\perp}+N^{\perp}}$;
 \item $M\subset(M^{\perp})^{\perp}$.
\end{enumerate}
En (ii) suponemos que $M,N$ son subespacios cerrados de un espacio de Hilbert. Explicaremos por qué no basta, en general, suponer que $M,N$ sean distinguidos en un espacio incompleto. Los incisos (i) y (iii) no requieren cerradura ni completitud.
\end{ejercicio}
\begin{solucion}
\emph{(i)} Si $x\perp M+N$, tomando sucesivamente $n=0$ y $m=0$ obtenemos $x\perp M$ y $x\perp N$. Si estas dos condiciones se cumplen, $\ip{x}{m+n}=\ip{x}{m}+\ip{x}{n}=0$, y $x\perp M+N$.

\emph{(iii)} Si $m\in M$ y $y\in M^{\perp}$, $\ip{m}{y}=\overline{\ip{y}{m}}=0$, por lo que $m\in(M^{\perp})^{\perp}$.

\emph{(ii), caso Hilbert.} Si $M,N$ son distinguidos, son cerrados y satisfacen $M^{\perp\perp}=M$, $N^{\perp\perp}=N$. Por (i),
\[
 (M^{\perp}+N^{\perp})^{\perp}=M\cap N.
\]
Tomamos complementos otra vez y aplicamos el corolario \ref{ch30:doble}. Esto produce exactamente el cierre que aparece en el enunciado. No se puede suprimir: las sumas de subespacios cerrados no son necesariamente cerradas.

\emph{Por qué hace falta cuidar el espacio ambiente.} Damos un contraejemplo completo a la generalidad prehilbertiana. Sea $A$ el álgebra real de funciones sobre $[0,1]$ generada por $t$ y $(1+t^{2})^{-1}$; sus elementos son funciones racionales continuas. En $V=A\times A$ definamos
\[
 \ip{(f,g)}{(u,v)}=f(0)u(0)+g(0)v(0)
 +\int_0^1(fu+gv)\,dt.
\]
Es un producto interno: la anulación de la integral de $f^{2}+g^{2}$, por continuidad, obliga a $f=g=0$. Los operadores
\[
 P(f,g)=(f,0),\qquad
 Q(f,g)=\frac{f+tg}{1+t^{2}}(1,t)
\]
preservan $V$, son idempotentes y sus matrices en cada $t$ son simétricas. Por integración, $\ip{Px}{y}=\ip{x}{Py}$ y $\ip{Qx}{y}=\ip{x}{Qy}$. Sus imágenes $M=P(V)$ y $N=Q(V)$ son distinguidas: $x=Px+(x-Px)$ y el segundo sumando es perpendicular a la imagen, y lo mismo para $Q$.

Ahora $M=\{(f,0):f\in A\}$ y $N=\{(b,tb):b\in A\}$. La intersección es cero, porque $tb(t)=0$ para $t>0$ implica $b=0$ también en $0$ por continuidad. Además,
\[
 M^{\perp}=\{(0,g):g\in A\},\qquad
 N^{\perp}=\{(-tb,b):b\in A\},
\]
pues son los núcleos de $P$ y $Q$, respectivamente. Todo vector de su suma tiene primera componente cero en $0$. Esta condición se conserva al tomar límites en $V$, ya que $|f(0)|\leq\norm{(f,g)}$. Por tanto $(1,0)\notin\cl{M^{\perp}+N^{\perp}}$, mientras $(M\cap N)^{\perp}=V$. La igualdad original requiere, pues, una hipótesis adicional; el contexto Hilbert proporciona la versión utilizada en los capítulos~\ref{cap:nuevo15} y~\ref{cap:nuevo16}.
\end{solucion}

\begin{ejercicio}[Lectura 30, ejercicio 5]\label{ej:lectura-30-5}
Probar que $P_M$ es lineal, continuo, idempotente y simétrico respecto al producto interno. Determinar su norma.
\end{ejercicio}
\begin{solucion}
Si $x=m_x+n_x$ y $y=m_y+n_y$ son las descomposiciones en $M\oplus M^{\perp}$, entonces
\[
 \alpha x+\beta y=(\alpha m_x+\beta m_y)+(\alpha n_x+\beta n_y).
\]
La unicidad de la descomposición da $P_M(\alpha x+\beta y)=\alpha P_Mx+\beta P_My$. Por Pitágoras, $\norm{P_Mx}\leq\norm{x}$, por lo que es continuo y $\norm{P_M}\leq1$. Si $M\ne\{0\}$, tomamos $m\in M$ de norma uno: $P_Mm=m$, y $\norm{P_M}=1$. Si $M=\{0\}$, la norma es cero, excepción que debe añadirse al enunciado original.

Como $P_Mx\in M$ y la proyección de un vector de $M$ es él mismo, $P_M^{2}=P_M$. Finalmente,
\[
 \ip{P_Mx}{y}=\ip{m_x}{m_y+n_y}=\ip{m_x}{m_y}
 =\ip{m_x+n_x}{m_y}=\ip{x}{P_My}.
\]
En particular, la proyección perpendicular es una contracción: no amplifica distancias, pues $\norm{P_Mx-P_My}\leq\norm{x-y}$.
\end{solucion}

\begin{ejercicio}[Lectura 30, ejercicio 6]\label{ej:lectura-30-6}
Sea $P:V\to V$ lineal, $P^{2}=P$, y $\ip{Px}{y}=\ip{x}{Py}$ para todos $x,y$. Probar que $P$ es una proyección ortogonal.
\end{ejercicio}
\begin{solucion}
Sea $M=P(V)$. Para $m=Py\in M$,
\[
 \ip{x-Px}{m}=\ip{x-Px}{Py}=\ip{P(x-Px)}{y}=0.
\]
Por tanto $Px\in M$ y $x-Px\in M^{\perp}$: la caracterización del teorema \ref{ch30:caracterizacion} da $Px=P_Mx$. El ejercicio~\ref{ej:lectura-30-5} demuestra automáticamente su continuidad; no era necesario suponerla. Además $\ker P=M^{\perp}$: si $Px=0$, la descomposición dice $x\in M^{\perp}$; si $x\in M^{\perp}$, la unicidad de la descomposición obliga a $Px=0$.
\end{solucion}

\chapter{Bases ortonormales y desarrollos de Fourier}

\begin{ejercicio}[Ampliación: una base para las sucesiones convergentes]\label{ej:ampl-schauder-1}
En el espacio $c$ con la norma uniforme pruebe que $(\mathbf1,e_1,e_2,\ldots)$ es una base de Schauder. Determine los coeficientes, el error de las sumas parciales y cotas para los funcionales coordenados. Este ejercicio adapta la propuesta de encontrar una base de $c$ en \cite{max-schauder}; la construcción y la solución se desarrollan aquí.
\end{ejercicio}
\begin{solucion}
Para $x\in c$, escriba $L=\lim_nx_n$. Definimos
\[
 P_Nx=L\mathbf1+\sum_{n=1}^N(x_n-L)e_n.
\]
En las primeras $N$ coordenadas, $P_Nx$ coincide con $x$; en las restantes vale $L$. Por ello
$\norm{x-P_Nx}_\infty=\sup_{n>N}|x_n-L|\to0$.

Si una expansión con coeficiente inicial $a_0$ y coeficientes posteriores $a_n$ converge a $x$, la función límite de sucesiones $\ell:c\to\F$, $\ell(y)=\lim_ky_k$, es continua porque $|\ell(y)|\le\norm y_\infty$. Aplicada a cada suma parcial da $a_0$, luego $a_0=L$. Después, evaluar en la coordenada $n$ da $x_n=a_0+a_n$, así que $a_n=x_n-L$. Hemos probado unicidad, no sólo encontrado una representación.

Los funcionales satisfacen $|a_0|\le\norm x_\infty$ y $|a_n|\le2\norm x_\infty$. Además $\norm{P_Nx}_\infty\le\norm x_\infty$, pues $|L|\le\norm x_\infty$ y todas las demás coordenadas de $P_Nx$ son coordenadas de $x$. La componente constante recoge precisamente la información que desaparece al pasar al cociente $c/c_0$ del ejercicio~\ref{ej:ampl-cocientes-2}.
\end{solucion}

\begin{ejercicio}[Ampliación: independencia y totalidad no bastan]\label{ej:ampl-schauder-2}
En $\ell^2$ defina $b_1=e_1$ y $b_n=e_1+n^{-1}e_n$ para $n\ge2$. Pruebe que la lista es linealmente independiente y total, pero no es una base de Schauder. Examine el vector $x=(0,1/2,1/3,\ldots)$. Ejercicio propio de esta edición.
\end{ejercicio}
\begin{solucion}
Si una combinación finita $\sum_{n=1}^Na_nb_n$ es cero, su coordenada $k\ge2$ es $a_k/k$, luego $a_k=0$ para $2\le k\le N$. Su primera coordenada deja entonces $a_1=0$. La lista es independiente. Además $e_1=b_1$ y $e_n=n(b_n-b_1)$ para $n\ge2$. Su espacio generado contiene todos los vectores coordenados y, por tanto, es denso en $\ell^2$.

El vector indicado pertenece a $\ell^2$. Si tuviera una expansión convergente $x=\sum_na_nb_n$, evaluar en la coordenada $k\ge2$ obligaría a $a_k/k=1/k$, es decir, $a_k=1$. La primera coordenada de la suma parcial número $N$ sería $a_1+N-1$, que no converge al valor $0$ de $x$. La convergencia en $\ell^2$ implica convergencia de cada coordenada, así que hemos obtenido una contradicción. Aunque pueden elegirse combinaciones finitas que aproximen $x$, sus coeficientes no se ensamblan en una expansión ordenada fija.
\end{solucion}

\begin{ejercicio}[Ampliación: separabilidad y una obstrucción]\label{ej:ampl-schauder-3}
Pruebe que todo espacio con base de Schauder es separable. Demuestre después directamente que $\ell^\infty$ no es separable y concluya que no admite ninguna base de Schauder numerable, no sólo que su lista coordenada falla. La primera implicación y esta aplicación desarrollan la discusión de \cite{max-schauder}.
\end{ejercicio}
\begin{solucion}
Sea $(b_n)$ una base. Considere las combinaciones finitas de los $b_n$ con coeficientes racionales, o racionales complejos en el caso complejo. Forman un conjunto numerable $D$: para cada longitud hay un producto finito de conjuntos numerables, y después se toma una unión numerable. Dado $x$ y $\eps>0$, elija una suma parcial $s_N$ con $\norm{x-s_N}<\eps/2$. Aproximamos sus coeficientes $a_j$ por racionales $q_j$ de modo que
\[
 \sum_{j=1}^N|a_j-q_j|\norm{b_j}<\eps/2.
\]
La desigualdad triangular da un elemento de $D$ a distancia menor que $\eps$ de $x$. Así $D$ es denso.

En $\ell^\infty$, las sucesiones con valores en $\{0,1\}$ forman un conjunto no numerable: una supuesta lista se contradice cambiando el dígito diagonal. Dos sucesiones distintas de ese conjunto tienen distancia uniforme $1$, por lo que las bolas abiertas de radio $1/3$ centradas en ellas son disjuntas. Si existiera un conjunto denso numerable, cada una de esas bolas contendría uno de sus elementos; elegir el primero de una enumeración produciría una inyección del conjunto no numerable de centros en $\N$, contradicción. Así $\ell^\infty$ no es separable y la primera parte descarta cualquier base de Schauder secuencial.
\end{solucion}

\begin{ejercicio}[Ampliación: transportar una base a otra]\label{ej:ampl-schauder-4}
Sea $T:V\to W$ una biyección lineal con $\norm{Tx}\le C\norm x$ y $\norm{T^{-1}y}\le D\norm y$. Demuestre que una base de Schauder $(b_n)$ de $V$ se transforma en la base $(Tb_n)$ de $W$. Aplique el resultado en $\ell^2$ a $T=I+\lambda R$, donde $R(x_1,x_2,\ldots)=(0,x_1,x_2,\ldots)$ y $|\lambda|<1$. Construya la inversa y calcule los coeficientes de la base $e_n+\lambda e_{n+1}$. Ejercicio propio de esta edición.
\end{ejercicio}
\begin{solucion}
Para $y\in W$, expanda $x=T^{-1}y=\sum_na_nb_n$. La cota de $T$ permite pasar al límite: $y=\sum_na_nTb_n$. Si otra lista de coeficientes produjera $y$, aplicar $T^{-1}$, que también es continua por su cota, daría otra expansión de $x$, contradiciendo unicidad.

En $\ell^2$, $R$ es una isometría, pues insertar un cero no cambia la suma de cuadrados. Para cada $x$ la serie
\[
 Ux=\sum_{k=0}^{\infty}(-\lambda)^kR^kx
\]
converge absolutamente en norma, ya que la suma de normas es a lo sumo $\norm x_2/(1-|\lambda|)$. La completitud de $\ell^2$ garantiza su existencia. Sumar series y multiplicar por escalares muestra linealidad, y $\norm{Ux}_2\le\norm x_2/(1-|\lambda|)$ muestra continuidad. En las sumas finitas se verifica
\[
 (I+\lambda R)\sum_{k=0}^N(-\lambda)^kR^k
 =I-(-\lambda R)^{N+1}.
\]
El resto aplicado a $x$ tiene norma $|\lambda|^{N+1}\norm x_2\to0$. Como los factores conmutan, la misma identidad vale en el orden inverso. Pasando al límite obtenemos $TU=UT=I$, de modo que $U=T^{-1}$.

Las coordenadas $a_n$ de $x$ respecto de $Te_n=e_n+\lambda e_{n+1}$ son las coordenadas de $Ux$. En la coordenada $n$ sólo contribuyen $k=0,\ldots,n-1$, por lo que
\[
 a_n=\sum_{k=0}^{n-1}(-\lambda)^k x_{n-k}.
\]
Hemos construido una base de Schauder sin ortogonalidad: si $\lambda\ne0$, el producto de dos vectores consecutivos de la lista es $\lambda$, que no se anula. La convergencia y unicidad provienen aquí de un cambio de coordenadas continuo con inversa continua.
\end{solucion}

\begin{ejercicio}[Ampliación: una base cuyo orden importa]\label{ej:ampl-schauder-5}
En $c_0$ con norma uniforme, sea $s_n=e_1+\cdots+e_n$. Pruebe que $(s_n)$ es una base de Schauder, calcule sus coeficientes y la norma del operador de suma parcial. Usando $x_n=(-1)^{n-1}/n$, construya una permutación de esta base que deje de ser base de Schauder. Ejercicio propio de esta edición.
\end{ejercicio}
\begin{solucion}
Ponga $a_n=x_n-x_{n+1}$. La suma parcial $P_Nx=\sum_{n=1}^Na_ns_n$ tiene, para $k\le N$, coordenada
\[
 \sum_{n=k}^N(x_n-x_{n+1})=x_k-x_{N+1},
\]
y para $k>N$ tiene coordenada cero. Por tanto
\[
 \norm{x-P_Nx}_\infty
 =\max\{|x_{N+1}|,\sup_{k>N}|x_k|\}
 =\sup_{k>N}|x_k|\longrightarrow0.
\]
La diferencia entre las coordenadas $k$ y $k+1$ de una combinación de los $s_n$ es exactamente el coeficiente de $s_k$, si ese vector aparece, y cero en caso contrario. Como ambas evaluaciones son continuas, cualquier expansión convergente de $x$ obliga a ese coeficiente a ser $x_k-x_{k+1}$. Esto prueba unicidad.

La fórmula de $P_Nx$ da $\norm{P_Nx}_\infty\le2\norm x_\infty$. Para obtener igualdad, tome un vector de soporte finito con $x_1=1$, $x_{N+1}=-1$ y las otras coordenadas cero. Tiene norma $1$ y la primera coordenada de $P_Nx$ es $2$. Por tanto $\norm{P_N}=2$ para todo $N\ge1$.

Para el vector especificado,
\[
 a_n=(-1)^{n-1}\left(\frac1n+\frac1{n+1}\right).
\]
La suma de los términos positivos diverge a $+\infty$ y la de los valores absolutos de los negativos también diverge: cada una domina una subserie armónica sobre índices pares o impares. Construimos una permutación así: en la etapa $k$, añadimos al menos el siguiente término positivo aún no utilizado, y cuantos positivos adicionales sean necesarios para que la suma escalar supere $k$; después añadimos el siguiente término negativo aún no utilizado. La divergencia positiva permite completar cada etapa. Se usa todo término positivo, porque se añade al menos uno por etapa en su orden, y todo negativo, porque se añade exactamente uno por etapa en su orden. Es, pues, una permutación de todos los índices.

Justo antes de cada término negativo hay una suma parcial escalar mayor que $k$. En la correspondiente serie de vectores $\sum a_{\pi(n)}s_{\pi(n)}$, la primera coordenada de toda suma parcial es precisamente esa suma escalar, puesto que cada $s_j$ tiene primera coordenada $1$. Las sumas parciales no pueden converger en norma.

Falta descartar una expansión con otros coeficientes después de permutar la lista. El argumento de la diferencia entre las coordenadas $k$ y $k+1$ sigue funcionando: una vez que aparece $s_k$, esa diferencia de las sumas parciales es su coeficiente y ningún otro vector la modifica. Cualquier expansión convergente de $x$ debe asignar a $s_k$ el mismo $a_k$ calculado antes. Por tanto, para esta permutación, $x$ no tiene expansión convergente y la lista permutada no es base de Schauder. La ortogonalidad del caso anterior evitaba exactamente este problema de orden.
\end{solucion}

\begin{ejercicio}[Lectura 31, ejercicio 1: Gram--Schmidt]\label{ej:lectura-31-1}
\index{Gram--Schmidt, procedimiento de}
Para una familia finita o numerable linealmente independiente $v_1,v_2,\dots$ en un espacio con producto interno, definir
\[
 w_1=v_1,\qquad
 w_{n+1}=v_{n+1}-\sum_{k=1}^{n}\frac{\ip{v_{n+1}}{w_k}}{\norm{w_k}^{2}}w_k.
\]
Probar la ortogonalidad, la no nulidad, la igualdad de los espacios generados y la existencia de una base ortonormal en todo espacio prehilbertiano separable.
\end{ejercicio}
\begin{solucion}
Procedemos por inducción. $w_1\ne0$ por independencia. Supongamos que $w_1,\dots,w_n$ son no nulos, ortogonales y que generan lo mismo que $v_1,\dots,v_n$. Para $j\leq n$,
\[
 \ip{w_{n+1}}{w_j}
 =\ip{v_{n+1}}{w_j}
 -\sum_{k=1}^{n}\frac{\ip{v_{n+1}}{w_k}}{\norm{w_k}^{2}}\ip{w_k}{w_j}=0.
\]
Si $w_{n+1}=0$, la fórmula expresaría $v_{n+1}$ en el espacio generado por $v_1,\dots,v_n$, contradicción. La misma fórmula expresa $w_{n+1}$ mediante $v_1,\dots,v_{n+1}$ y, despejándola, expresa $v_{n+1}$ mediante $w_1,\dots,w_{n+1}$. Esto prueba la igualdad de los espacios generados en cada etapa y, tomando uniones, para las familias enteras. La familia $(w_n)$ es independiente porque es ortogonal y no tiene vectores cero. Finalmente, $e_n=w_n/\norm{w_n}$ es ortonormal y conserva el espacio generado.

Si $V$ es separable y no nulo, elijamos una sucesión densa $(d_n)$. Recorremos esa lista y conservamos $d_n$ únicamente cuando no pertenece al espacio generado por los elementos ya conservados. Cada elemento descartado pertenece al espacio generado por los conservados anteriores. La familia resultante es independiente, finita o numerable, y su espacio generado contiene toda la lista densa. Gram--Schmidt produce una familia ortonormal total. Si $V=\{0\}$, se utiliza la familia vacía. No se necesita completitud porque cada etapa utiliza solamente sumas finitas.
\end{solucion}

\begin{ejercicio}[Lectura 31, ejercicio 2: primeros polinomios de Legendre]\label{ej:lectura-31-2}
\index{Legendre, polinomios de}
Aplicar Gram--Schmidt a $1,x,x^{2},x^{3},x^{4}$ con $\ip{f}{g}=\int_{-1}^{1}f(x)\overline{g(x)}\,dx$.
\end{ejercicio}
\begin{solucion}
Usaremos $\int_{-1}^{1}x^{2j+1}\,dx=0$ y $\int_{-1}^{1}x^{2j}\,dx=2/(2j+1)$. La paridad evita productos entre un polinomio par y uno impar. Los residuos sin normalizar son
\begin{align*}
 w_0&=1, &\norm{w_0}^{2}&=2,\\
 w_1&=x, &\norm{w_1}^{2}&=2/3,\\
 w_2&=x^{2}-\tfrac13, &\norm{w_2}^{2}&=8/45,\\
 w_3&=x^{3}-\tfrac35x, &\norm{w_3}^{2}&=8/175,\\
 w_4&=x^{4}-\tfrac67x^{2}+\tfrac3{35}, &\norm{w_4}^{2}&=128/11025.
\end{align*}
En efecto, la proyección de $x^{2}$ sobre $1$ es $(2/3)/2=1/3$. La proyección de $x^{3}$ sobre $x$ tiene coeficiente $(2/5)/(2/3)=3/5$. Para $x^{4}$, la proyección sobre $1$ es $1/5$ y
\[
 \frac{\ip{x^{4}}{w_2}}{\norm{w_2}^{2}}
 =\frac{2/7-2/15}{8/45}=\frac67.
\]
Así $w_4=x^{4}-1/5-(6/7)(x^{2}-1/3)$, que es la expresión escrita. Las normas se obtienen expandiendo los cuadrados; por ejemplo,
\[
 \norm{w_2}^{2}=\frac25-\frac49+\frac29=\frac8{45},\qquad
 \norm{w_3}^{2}=\frac27-\frac{12}{25}+\frac6{25}=\frac8{175}.
\]
Para $w_4$, usar la ortogonalidad de sus proyecciones reduce el cálculo a
\[
 \norm{w_4}^{2}=\frac29-\left(\frac15\right)^2 2
 -\left(\frac67\right)^2\frac8{45}=\frac{128}{11025}.
\]
La familia ortonormal resultante es
\begin{align*}
 e_0&=\frac1{\sqrt2}, &
 e_1&=\sqrt{\frac32}\,x,\\
 e_2&=\sqrt{\frac52}\,\frac{3x^{2}-1}{2}, &
 e_3&=\sqrt{\frac72}\,\frac{5x^{3}-3x}{2},\\
 e_4&=\frac3{\sqrt2}\,\frac{35x^{4}-30x^{2}+3}{8}.
\end{align*}
Los factores polinómicos con valor uno en $x=1$ son $P_0,\dots,P_4$, los primeros polinomios de Legendre. Si continuamos Gram--Schmidt sobre todos los monomios, obtenemos una familia ortonormal total en $C([-1,1])$ para esta norma: los polinomios son uniformemente densos por el teorema de aproximación y $\norm{f}_2\leq\sqrt2\norm{f}_\infty$. Las cinco funciones escritas son solo el comienzo de esa familia.
\end{solucion}

\begin{ejercicio}[Lectura 31, ejercicio 3]\label{ej:lectura-31-3}
Probar que $(e^{inx})_{n\in\Z}$ es ortonormal en $C([-\pi,\pi])$ con el producto interno normalizado.
\end{ejercicio}
\begin{solucion}
Con $\ip{f}{g}=(2\pi)^{-1}\int_{-\pi}^{\pi}f\bar g$, se tiene
\[
 \ip{e^{inx}}{e^{imx}}=\frac1{2\pi}\int_{-\pi}^{\pi}e^{i(n-m)x}\,dx.
\]
Si $n=m$, la integral es $2\pi$ y el resultado es uno. Si $k=n-m\ne0$, una primitiva es $e^{ikx}/(ik)$; sus valores en los extremos coinciden porque $e^{ik\pi}=e^{-ik\pi}=(-1)^k$. La integral es cero. Sin el factor $1/(2\pi)$, las funciones seguirían siendo ortogonales, pero tendrían norma $\sqrt{2\pi}$ y habría que dividirlas por ese número para normalizarlas.
\end{solucion}

\begin{ejercicio}[Lectura 31, ejercicio 4: Müntz en norma integral]\label{ej:lectura-31-4}
\index{Muntz@Müntz, teorema de}\index{Gram, determinante de}
Sean $1\leq n_1<n_2<\cdots$ enteros y $M_k=\spanop\{x^{n_1},\dots,x^{n_k}\}\subset C([0,1])$, con $\norm{f}_2^{2}=\int_0^1|f|^{2}$. Probar que $\bigcup_kM_k$ es denso para esta norma si y solo si $\sum_k1/n_k=\infty$. Desarrollar (i) el cociente de determinantes de Gram y la fórmula de distancia; (ii) el criterio de las series logarítmicas con sus hipótesis; (iii) el análisis del producto resultante.
\end{ejercicio}
\begin{solucion}
La norma es esencial: se trata de densidad en norma integral, no uniforme. Todos los monomios utilizados se anulan en cero, por lo que no podrían aproximar uniformemente la constante uno. En norma integral, un valor en un solo punto no impide aproximar.

\emph{(i), reducción a monomios.} La densidad implica $\dist(x^m,M_k)\to0$ para cada entero $m\geq0$. Recíprocamente, si todos los monomios admiten tales aproximaciones, también todo polinomio: para $p=\sum_{j=0}^{r}a_jx^j$, se eligen aproximantes $q_j$ con errores suficientemente pequeños y se toma un mismo $M_k$ que contenga todos ellos; entonces $\norm{p-\sum a_jq_j}_2\leq\sum|a_j|\norm{x^j-q_j}_2$. La aproximación uniforme por polinomios implica aproximación en norma integral, y demuestra densidad en $C([0,1])$.

Si en la reducción se consideran únicamente $m\geq1$, todavía es suficiente: los polinomios $1-(1-x)^N$, sin término constante, aproximan a uno en norma integral porque
\[
 \int_0^1|(1-x)^N|^{2}\,dx=\frac1{2N+1}\longrightarrow0.
\]
Así la convención para $\N$ no altera esta reducción, pero la suma $\sum1/n_k$ exige exponentes positivos.

\emph{La fórmula de Gram.} Para vectores independientes $u_1,\dots,u_k$, sea $G=(\ip{u_j}{u_i})$, con la convención de la sección~\ref{ampl:gram}, y sea $u$ otro vector. Restar del último vector su proyección $p$ sobre $\spanop(u_1,\dots,u_k)$ equivale a operaciones de filas y columnas que no cambian el determinante de Gram: se resta a la última columna la combinación correspondiente y a la última fila su conjugada. El último vector queda $r=u-p$, perpendicular a todos los anteriores. La nueva matriz es diagonal por bloques, con bloques $G$ y $\norm r^{2}$. Por ello
\[
 \dist(u,\spanop(u_1,\dots,u_k))^{2}
 =\frac{\det\operatorname{Gram}(u_1,\dots,u_k,u)}{\det G}.
\]
El denominador es positivo: si $c\ne0$, $\sum c_iu_i\ne0$, y la forma cuadrática de Gram es su norma al cuadrado. En particular, $G$ es invertible.

Para monomios de exponentes reales $a_i>-1/2$, las entradas son $1/(a_i+a_j+1)$. Usaremos la identidad de Cauchy
\begin{equation}\label{ch31:cauchydet}
 \det\left(\frac1{x_i+y_j}\right)_{i,j=1}^{r}
 =\frac{\prod_{i<j}(x_j-x_i)(y_j-y_i)}{\prod_{i,j}(x_i+y_j)}.
\end{equation}
He aquí una justificación algebraica. Para $r=1$ es inmediata. Consideremos ambos miembros como funciones racionales de $x_r$, manteniendo las demás variables fijas y distintas, sin denominadores nulos. Sus únicas singularidades posibles son simples, en $x_r=-y_j$. El coeficiente del término $1/(x_r+y_j)$, llamado aquí \emph{residuo}, se calcula mediante el límite $\lim_{x_r\to-y_j}(x_r+y_j)R(x_r)$; no se necesita teoría de integración compleja. Expandiendo el determinante por su última fila, ese coeficiente es $(-1)^{r+j}$ veces el determinante de orden $r-1$ obtenido suprimiendo la fila $r$ y la columna $j$.

En el cociente de productos, cancelamos el factor $x_r+y_j$ y sustituimos $x_r=-y_j$. Los factores $\prod_{i<r}(x_r-x_i)$ se convierten en $(-1)^{r-1}\prod_{i<r}(x_i+y_j)$ y cancelan los factores restantes de la columna $j$ del denominador. Al dividir los factores de la segunda Vandermonde que contienen $y_j$ por $\prod_{l\ne j}(y_l-y_j)$, queda el signo $(-1)^{j-1}$. El signo total es $(-1)^{r-1}(-1)^{j-1}=(-1)^{r+j}$; el resto es el cociente de orden $r-1$, igual al determinante correspondiente por inducción. Así los coeficientes de todas las singularidades coinciden. La diferencia tiene denominador factorizado en los factores $x_r+y_j$ y ninguno subsiste después de cancelar, por lo que es un polinomio. Ambos miembros tienden a cero cuando $x_r\to\infty$, de modo que ese polinomio es cero. Esto prueba la identidad por inducción; los casos restantes permitidos se obtienen por continuidad.

Tomando $x_i=a_i$, $y_j=a_j+1$, obtenemos
\[
 \det\operatorname{Gram}(x^{a_1},\dots,x^{a_r})
 =\frac{\prod_{i<j}(a_j-a_i)^2}
 {\prod_i(2a_i+1)\prod_{i<j}(a_i+a_j+1)^2}.
\]
Al añadir el exponente $m$ y dividir los dos determinantes, todos los factores antiguos se cancelan y queda
\begin{equation}\label{ch31:muntzdist}
 \dist(x^m,M_k)^2
 =\frac1{2m+1}\prod_{i=1}^{k}\frac{(n_i-m)^2}{(n_i+m+1)^2}.
\end{equation}
Si $m=n_i$ para algún índice, ambos lados son cero una vez que ese monomio pertenece a $M_k$; no es necesario aplicar la fórmula de determinante a una familia independiente ampliada en ese caso.

\emph{(ii), el criterio logarítmico.} Para $a_n\geq0$, las series $\sum a_n$ y $\sum\log(1+a_n)$ convergen simultáneamente. En efecto, $\log(1+t)=\int_0^t(1+s)^{-1}\,ds$ da, para $0\leq t\leq1$,
\[
 \frac t2\leq\log(1+t)\leq t.
\]
Si $a_n\to0$, estas estimaciones comparan las colas. Si $a_n$ no tiende a cero, existe $\delta>0$ superado por infinitos términos y ambas series de términos no negativos divergen. La hipótesis de signo no puede omitirse. Por ejemplo, con
\[
 a_{2n-1}=\frac1{\sqrt{n+1}},\qquad a_{2n}=-\frac1{\sqrt{n+1}},
\]
la serie de $a_n$ converge a cero, pero la suma de cada par de logaritmos es $\log(1-1/(n+1))=\log(n/(n+1))$; las sumas pares de logaritmos son $-\log(N+1)$ y divergen a $-\infty$.

\emph{(iii), conclusión.} Fijemos $m\geq0$ que no figure entre los $n_i$. Como $n_i\to\infty$, salvo un número finito de factores tenemos $n_i>m$ y
\[
 \frac{n_i-m}{n_i+m+1}=1-b_i,
 \qquad b_i=\frac{2m+1}{n_i+m+1}\in(0,1).
\]
Cuando $b_i\leq1/2$, la fórmula integral para el logaritmo da
\[
 b_i\leq-\log(1-b_i)=\int_0^{b_i}\frac{dt}{1-t}\leq2b_i.
\]
Además $b_i/(1/n_i)\to2m+1>0$. Tomando logaritmos en \eqref{ch31:muntzdist},
\[
 \log\dist(x^m,M_k)^2=C_m+2\sum_{i=i_0}^{k}\log(1-b_i),
\]
con $C_m$ finito. Por las comparaciones anteriores, esta expresión tiende a $-\infty$ si y solo si $\sum_i1/n_i$ diverge. En tal caso la distancia tiende a cero para todo $m$ ausente; para los exponentes presentes la distancia ya es cero desde alguna etapa. La reducción de (i) prueba densidad.

Si $\sum_i1/n_i<\infty$, tomemos $m=0$, que no aparece porque todos los $n_i$ son positivos. Los logaritmos convergen a un número finito, así que \eqref{ch31:muntzdist} tiene límite estrictamente positivo: la constante uno no puede aproximarse. Esto prueba la necesidad. También aclara una precisión del inciso original: para $m$ que ya aparece en la lista, la distancia es cero independientemente de la convergencia de la serie; la equivalencia logarítmica se aplica a exponentes ausentes.
\end{solucion}

\begin{ejercicio}[Lectura 31, ejercicio 5]\label{ej:lectura-31-5}
Probar que un espacio de Hilbert es separable si y solo si tiene una base ortonormal a lo sumo numerable.
\end{ejercicio}
\begin{solucion}
Si es separable, el procedimiento con una lista densa del ejercicio~\ref{ej:lectura-31-1} produce la base requerida. Si tiene una base $E$ finita o numerable, consideremos todas sus combinaciones finitas con coeficientes racionales en el caso real, o de la forma $a+ib$ con $a,b\in\Q$ en el caso complejo. Ese conjunto es numerable: hay una cantidad numerable de conjuntos finitos de índices y de listas finitas de coeficientes de ese tipo. Es denso en $\spanop E$, porque en una combinación finita pueden aproximarse sus coeficientes y Pitágoras controla el error por la suma de sus cuadrados. Como $\spanop E$ es denso en $H$, también lo es ese conjunto numerable. El espacio cero tiene el conjunto denso $\{0\}$.

Además, cualquier conjunto ortonormal de un espacio separable es numerable: las bolas de radio $1/3$ centradas en sus elementos son disjuntas, ya que la distancia entre dos elementos distintos es $\sqrt2$. Cada bola contiene un elemento de una lista densa; asignando el primero de la lista que pertenece a ella se obtiene una inyección en $\N$.
\end{solucion}

\begin{ejercicio}[Lectura 31, ejercicio 6]\label{ej:lectura-31-6}
Calcular los cinco primeros coeficientes respecto de $e_0,\dots,e_4$ del ejercicio~\ref{ej:lectura-31-2} para (i) $x^2+x+1$, (ii) $\sin(\pi x)$, (iii) $\cos(\pi x)$ y (iv) $|x|$.
\end{ejercicio}
\begin{solucion}
Los resultados, en el orden $c_j=\ip{f}{e_j}$, son
\[
\begin{array}{c|ccccc}
 f&c_0&c_1&c_2&c_3&c_4\\ \hline
 x^2+x+1&4\sqrt2/3&\sqrt{2/3}&\frac23\sqrt{2/5}&0&0\\[2pt]
 \sin(\pi x)&0&\sqrt6/\pi&0&\sqrt{14}(\pi^{-1}-15\pi^{-3})&0\\[2pt]
 \cos(\pi x)&0&0&-3\sqrt{10}/\pi^2&0&\frac3{\sqrt2}(-20\pi^{-2}+210\pi^{-4})\\[2pt]
 |x|&1/\sqrt2&0&\frac14\sqrt{5/2}&0&-1/(8\sqrt2)
\end{array}
\]
Justificamos los cálculos para que la tabla no oculte el procedimiento.

\emph{(i)} Integrando los monomios,
\[
 c_0=\frac1{\sqrt2}\left(\frac23+2\right)=\frac{4\sqrt2}3,
 \qquad c_1=\sqrt{\frac32}\frac23=\sqrt{\frac23}.
\]
Como $\int P_2=0$ y $xP_2$ es impar,
$c_2=\sqrt{5/2}\int x^2P_2=\sqrt{5/2}(4/15)=\frac23\sqrt{2/5}$.
Los polinomios $e_3,e_4$ son ortogonales a todos los de grado menor, por construcción, y los dos coeficientes restantes son cero.

\emph{(ii)} La paridad anula $c_0,c_2,c_4$. Integración por partes da
\[
 \int_{-1}^{1}x\sin(\pi x)\,dx=\frac2\pi,
 \qquad
 \int_{-1}^{1}x^3\sin(\pi x)\,dx=\frac2\pi-\frac{12}{\pi^3}.
\]
Para la segunda integral, la primera integración produce $2/\pi+(3/\pi)\int x^2\cos(\pi x)$ y otra integración da $\int x^2\cos(\pi x)=-4/\pi^2$. Sustituir en $P_3=(5x^3-3x)/2$ da $\int P_3\sin(\pi x)=2/\pi-30/\pi^3$, de donde salen los dos coeficientes no nulos.

\emph{(iii)} La paridad anula $c_1,c_3$, y $\int_{-1}^{1}\cos(\pi x)=0$ anula $c_0$. Las integraciones por partes anteriores y una adicional dan
\[
 \int x^2\cos(\pi x)=-\frac4{\pi^2},\qquad
 \int x^4\cos(\pi x)=-\frac8{\pi^2}+\frac{48}{\pi^4},
\]
donde ambas integrales se toman de $-1$ a $1$. Para la última, integrar una vez produce $-(4/\pi)\int x^3\sin(\pi x)$, que utiliza el cálculo previo. Insertar estos momentos en $P_2$ y $P_4$ produce $-6/\pi^2$ y $-20/\pi^2+210/\pi^4$, respectivamente; luego se multiplican por los factores de normalización.

\emph{(iv)} Los coeficientes impares se anulan. Para los pares,
\[
 \int_{-1}^{1}|x|\,dx=1,\qquad
 \int_{-1}^{1}|x|P_2(x)\,dx=\int_0^1x(3x^2-1)\,dx=\frac14,
\]
\[
 \int_{-1}^{1}|x|P_4(x)\,dx
 =\frac14\left(\frac{35}{6}-\frac{30}{4}+\frac32\right)=-\frac1{24}.
\]
Los factores $1/\sqrt2$, $\sqrt{5/2}$ y $3/\sqrt2$ dan la última fila.
\end{solucion}

\begin{ejercicio}[Lectura 31, ejercicio 7]\label{ej:lectura-31-7}
Si $E$ es ortonormal y $x\in H$, estimar el número $n_k$ de vectores $e\in E$ con $|\ip{x}{e}|>1/k$.
\end{ejercicio}
\begin{solucion}
Por Bessel, el conjunto es finito. Si contiene $n_k>0$ elementos, sumando las desigualdades estrictas de sus cuadrados obtenemos
\[
 \frac{n_k}{k^2}<\sum_{|\ip{x}{e}|>1/k}|\ip{x}{e}|^{2}\leq\norm{x}^{2},
\]
y, por tanto, $n_k<k^2\norm{x}^{2}$. Si $n_k=0$ y $x\ne0$, la misma conclusión estricta vale porque el lado derecho es positivo. Si $x=0$, entonces $n_k=0$ y la desigualdad estricta original se convierte en $0<0$, que es falsa. La formulación válida para todo $x$ es $n_k\leq k^2\norm{x}^{2}$, con desigualdad estricta cuando $x\ne0$.
\end{solucion}

\chapter{Representación de Riesz y operadores adjuntos}

\begin{ejercicio}[Lectura 32, ejercicio 1]\label{ej:lectura-32-1}
Probar las reglas (i) $(\alpha S+\beta T)^*=\bar\alpha S^*+\bar\beta T^*$; (ii) $(ST)^*=T^*S^*$; (iii) $T^{**}=T$; (iv) $I^*=I$; (v) $(T^{-1})^*=(T^*)^{-1}$ para $T$ invertible.
\end{ejercicio}
\begin{solucion}
En todas las verificaciones utilizamos la unicidad del adjunto, una vez identificada la transformación que satisface su relación definitoria.

\emph{(i)} Para $x\in H_1,y\in H_2$,
\[
 \ip{(\alpha S+\beta T)x}{y}
 =\alpha\ip{x}{S^*y}+\beta\ip{x}{T^*y}
 =\ip{x}{(\bar\alpha S^*+\bar\beta T^*)y}.
\]
\emph{(ii)} Para $T:H_1\to H_2$, $S:H_2\to H_3$,
$\ip{STx}{z}=\ip{Tx}{S^*z}=\ip{x}{T^*S^*z}$.
\emph{(iii)} La identidad $\ip{Tx}{y}=\ip{x}{T^*y}$, al conjugar, se convierte en $\ip{T^*y}{x}=\ip{y}{Tx}$, que caracteriza al adjunto de $T^*$.
\emph{(iv)} $\ip{Ix}{y}=\ip{x}{Iy}$.
\emph{(v)} De $T^{-1}T=TT^{-1}=I$, las reglas (ii) y (iv) dan
$T^*(T^{-1})^*=I$ y $(T^{-1})^*T^*=I$. Así $(T^{-1})^*$ es la inversa bilateral de $T^*$. Aquí invertible en $\Bcal(H)$ significa que la inversa es acotada; para una biyección acotada entre espacios de Banach esa propiedad se deduce del teorema de la aplicación abierta.
\end{solucion}

\begin{ejercicio}[Lectura 32, ejercicio 2]\label{ej:lectura-32-2}
\index{operador!normal}
Si $T$ es normal, es decir, $T^*T=TT^*$, probar $\norm{Tx}=\norm{T^*x}$ y que $T-\alpha I$ es normal para todo $\alpha\in\C$.
\end{ejercicio}
\begin{solucion}
La identidad de adjunción da
\[
 \norm{Tx}^{2}=\ip{x}{T^*Tx}=\ip{x}{TT^*x}=\norm{T^*x}^{2}.
\]
Ambas normas son no negativas, por lo que sus raíces coinciden. Además,
\begin{align*}
 (T-\alpha I)^*(T-\alpha I)
 &=T^*T-\alpha T^*-\bar\alpha T+|\alpha|^2I,\\
 (T-\alpha I)(T-\alpha I)^*
 &=TT^*-\bar\alpha T-\alpha T^*+|\alpha|^2I.
\end{align*}
La normalidad de $T$ hace iguales estas expresiones.
\end{solucion}

\begin{ejercicio}[Lectura 32, ejercicio 3]\label{ej:lectura-32-3}
\index{operador!autoadjunto}\index{operador!positivo}
Probar que $T+T^*$ y $TT^*$ son autoadjuntos, esto es, coinciden con sus adjuntos.
\end{ejercicio}
\begin{solucion}
Las reglas del ejercicio~\ref{ej:lectura-32-1} dan $(T+T^*)^*=T^*+T^{**}=T^*+T$ y $(TT^*)^*=(T^*)^*T^*=TT^*$. De manera adicional, $TT^*$ es positivo en el sentido de que $\ip{TT^*x}{x}=\norm{T^*x}^{2}\geq0$; ser autoadjunto, por sí solo, no exige esta desigualdad.
\end{solucion}

\begin{ejercicio}[Lectura 32, ejercicio 4: isometrías y unitarios]\label{ej:lectura-32-4}
\index{operador!unitario}\index{isometria@isometría!lineal}
Analizar la afirmación: $\ip{Tx}{Ty}=\ip{x}{y}$ para todo $x,y$ si y solo si $T$ es invertible y $T^{-1}=T^*$.
\end{ejercicio}
\begin{solucion}
Tal como está escrita en dimensión infinita, la afirmación es falsa. El corrimiento $R$ preserva productos internos, pero ningún vector tiene imagen $e_1$. La condición escrita define una \emph{isometría lineal}; un operador \emph{unitario} es una isometría lineal \emph{suprayectiva}.

Si $T$ preserva productos, entonces $\ip{x}{T^*Ty}=\ip{Tx}{Ty}=\ip{x}{y}$ para todo $x$, y $T^*Ty=y$. Por tanto $T^*T=I$. Esto da una inversa por la izquierda, sin asegurar una inversa por la derecha. Si además $T$ es suprayectivo, es biyectivo, pues una isometría tiene núcleo cero. Su inversa preserva normas, así que es acotada. Multiplicando $T^*T=I$ por $T^{-1}$ se obtiene $T^*=T^{-1}$.

Recíprocamente, si $T$ es invertible y $T^*=T^{-1}$, entonces
$\ip{Tx}{Ty}=\ip{x}{T^*Ty}=\ip{x}{y}$; es una isometría suprayectiva. Una formulación equivalente es $T^*T=TT^*=I$: la primera igualdad da inyectividad y preservación del producto, y la segunda expresa cada $y$ como $T(T^*y)$, dando suprayectividad.
\end{solucion}

\begin{ejercicio}[Lectura 32, ejercicio 5]\label{ej:lectura-32-5}
Probar que la proyección ortogonal $P_W$ sobre un subespacio cerrado es autoadjunta.
\end{ejercicio}
\begin{solucion}
Descompongamos $x=w_x+n_x$, $y=w_y+n_y$ en $W\oplus W^{\perp}$. Entonces
\[
 \ip{P_Wx}{y}=\ip{w_x}{w_y+n_y}=\ip{w_x}{w_y}
 =\ip{w_x+n_x}{w_y}=\ip{x}{P_Wy}.
\]
La unicidad del adjunto da $P_W^*=P_W$.
\end{solucion}

\begin{ejercicio}[Lectura 32, ejercicio 6]\label{ej:lectura-32-6}
Si $T^*T=0$, probar $T=0$. Decidir si la conclusión vale cuando $TT^*=0$.
\end{ejercicio}
\begin{solucion}
Si $T^*T=0$, para cada $x$ se tiene $\norm{Tx}^{2}=\ip{x}{T^*Tx}=0$, de modo que $Tx=0$. Si $TT^*=0$, aplicamos lo anterior al operador $T^*$: $\norm{T^*x}^{2}=\ip{x}{TT^*x}=0$, así que $T^*=0$ y $T=T^{**}=0$.
\end{solucion}

\begin{ejercicio}[Lectura 32, ejercicio 7]\label{ej:lectura-32-7}
Demostrar que $\norm{Tv}=\norm T\norm v$ si y solo si $T^*Tv=\norm T^{2}v$.
\end{ejercicio}
\begin{solucion}
Sea $A=\norm T^{2}I-T^*T$. Es autoadjunto y para todo $x$,
\[
 q(x):=\ip{Ax}{x}=\norm T^{2}\norm x^{2}-\norm{Tx}^{2}\geq0.
\]
La igualdad de normas propuesta equivale a $q(v)=0$. Demostraremos que esta igualdad implica $Av=0$, un paso que necesita justificación y no se sigue solo de que $A$ sea autoadjunto.

Fijemos $z$ y escribamos $a=\ip{Av}{z}$. Para cualquier escalar $t$, la autoadjunción da
\[
0\leq q(v+tz)=2\operatorname{Re}(\bar t a)+|t|^{2}q(z).
\]
Si $a\ne0$, elegimos $t=-sa$ con $s>0$. El miembro derecho es $|a|^{2}(-2s+s^{2}q(z))$, negativo para $s$ suficientemente pequeño; contradicción. Por tanto $\ip{Av}{z}=0$ para todo $z$ y, con $z=Av$, se obtiene $Av=0$.

Recíprocamente, si $T^*Tv=\norm T^{2}v$, entonces $\norm{Tv}^{2}=\ip{v}{T^*Tv}=\norm T^{2}\norm v^{2}$; tomando raíces se obtiene la igualdad. Para $v=0$ ambas condiciones se cumplen trivialmente; para hablar de que el operador «alcanza su norma» suele exigirse un vector unitario.
\end{solucion}

\begin{ejercicio}[Lectura 32, ejercicio 8]\label{ej:lectura-32-8}
Probar $\ker T=(T^*(H))^{\perp}$.
\end{ejercicio}
\begin{solucion}
Para $x\in H$,
\[
 x\in(T^*(H))^{\perp}
 \Longleftrightarrow \ip{x}{T^*y}=0\text{ para todo }y
 \Longleftrightarrow \ip{Tx}{y}=0\text{ para todo }y
 \Longleftrightarrow Tx=0.
\]
La última equivalencia se prueba tomando $y=Tx$ en una dirección; en la otra, todos los productos son cero. También resulta $\cl{T^*(H)}=(\ker T)^{\perp}$ al tomar complementos y utilizar el doble complemento en Hilbert. El cierre no debe omitirse si no se sabe que el rango es cerrado.
\end{solucion}

\begin{ejercicio}[Lectura 32, ejercicio 9: polarización de una forma]\label{ej:lectura-32-9}
Sea $\xi$ sesquilineal y $\widehat\xi(v)=\xi(v,v)$. Probar (i)
\[
 \xi(u,v)=\sum_{k=1}^{4}i^k\widehat\xi\left(\frac{u+i^kv}{2}\right),
\]
y (ii) que la forma cuadrática determina la forma sesquilineal.
\end{ejercicio}
\begin{solucion}
\emph{(i)} Por sesquilinealidad,
\[
 \widehat\xi\left(\frac{u+i^kv}{2}\right)
 =\frac14\bigl(\xi(u,u)+\xi(v,v)+i^{-k}\xi(u,v)+i^k\xi(v,u)\bigr).
\]
Multiplicando por $i^k$ y sumando, las sumas $\sum_{k=1}^{4}i^k$ y $\sum_{k=1}^{4}i^{2k}$ son cero, mientras que el coeficiente de $\xi(u,v)$ es $4/4=1$. No se utilizó simetría conjugada de $\xi$: la identidad vale para toda forma sesquilineal.

\emph{(ii)} Si dos formas tienen la misma función diagonal, los cuatro términos del lado derecho coinciden para cada par $(u,v)$; la identidad prueba que las formas son iguales. La hipótesis compleja es relevante: para formas bilineales reales generales, la parte antisimétrica se anula en la diagonal y no puede recuperarse de ella.
\end{solucion}

\begin{ejercicio}[Lectura 32, ejercicio 10: multiplicación]\label{ej:lectura-32-10}
Sea $h\in C([a,b],\C)$, $a<b$, y $T_hf=hf$ en $C([a,b])$ con producto integral. Hallar (i) el adjunto, (ii) cuándo es normal, (iii) cuándo es autoadjunto y (iv) cuándo es unitario.
\end{ejercicio}
\begin{solucion}
Este espacio no es completo para la norma integral. La existencia general de adjuntos de Hilbert no se aplica automáticamente, pero en este caso podemos construir el adjunto explícitamente dentro del espacio. La acotación se obtiene de $\norm{hf}_2\leq\norm h_\infty\norm f_2$.

\emph{(i)} Para $f,g$ continuas,
\[
 \ip{T_hf}{g}=\int_a^b hf\bar g
 =\int_a^b f\,\overline{\bar h g}=\ip{f}{T_{\bar h}g}.
\]
El mismo cálculo funciona si el producto está dividido por $b-a$. Por tanto $T_h^*=T_{\bar h}$; la unicidad se prueba usando la positividad del producto interno como antes.

\emph{(ii)} Para toda $h$, $T_h^*T_h=T_{|h|^{2}}=T_hT_h^*$, así que siempre es normal.

\emph{(iii)} Es autoadjunto si y solo si $h=\bar h$, es decir, si $h$ toma valores reales. La necesidad se ve aplicando $T_h=T_{\bar h}$ a la función constante uno; la suficiencia se lee directamente del adjunto calculado.

\emph{(iv)} Es unitario si y solo si $|h(x)|=1$ para todo $x$. Si esto sucede, preserva productos y tiene inversa $T_{\bar h}$. Recíprocamente, la preservación de la norma implica
\[
 \int_a^b (|h|^{2}-1)|f|^{2}=0\quad\text{para toda }f\in C([a,b]).
\]
Si la función real continua $|h|^{2}-1$ no fuera cero, tendría signo constante y módulo separado de cero en un intervalo. Una función triangular no nula soportada allí daría una integral no nula, contradicción. En este ejemplo, a diferencia del corrimiento, la preservación de normas ya obliga a una inversa explícita dentro de $C([a,b])$.
\end{solucion}
